% !TeX program = xelatex
% =====================================================================
%  main.tex —— 文档根文件
%
%  本文件不写正文，只做三件事：装宏包、按序 \input、印全局部件。
%  "一个插件一个模块章"的模块化就落在 modules/ 下的 \input 清单上：
%  增删一个模块 = 增删一行。
% =====================================================================

\documentclass[
  UTF8,
  12pt,
  oneside,
  openany,
  fontset=windows
]{ctexbook}

% =====================================================================
%  packages.tex —— 宏包加载与引擎设定
%  引擎：XeLaTeX（中文经 ctex，字体走 Windows 字体集）
%  加载顺序有讲究：biblatex -> hyperref -> cleveref -> imakeidx
%  请勿在此文件写正文，也不要写版式参数（那些属于 style.tex）
% =====================================================================

% ---- 页面几何 --------------------------------------------------------
\usepackage[
  top=25mm, bottom=25mm,
  left=28mm, right=25mm,
  headsep=6mm, footskip=10mm
]{geometry}

% ---- 数学 ------------------------------------------------------------
\usepackage{amsmath}
\usepackage{amssymb}
\usepackage{mathtools}
\usepackage{bm}

% ---- 表格与浮动体 ----------------------------------------------------
\usepackage{booktabs}
\usepackage{array}
\usepackage{tabularx}
\usepackage{longtable}
\usepackage{multirow}
\usepackage{float}
\usepackage{caption}

% ---- 列表 ------------------------------------------------------------
\usepackage{enumitem}

% ---- 可断行的等宽标识符 ----------------------------------------------
%  xurl 让 \path 可以在任意标点处断行。本文档的 \pid（模块标识，
%  形如 physics.quantum.analytic_tw）用它来避免窄表格列里的溢出盒。
%  ⚠ 只给 \pid 用：\path 是 verbatim 扫描，会\emph{吞掉空格}，
%  因此含空格的 \akey/\fpath 继续用 \detokenize（见 style.tex）。
\usepackage{xurl}

% ---- 颜色与图形 ------------------------------------------------------
\usepackage{xcolor}
\usepackage{graphicx}

%  调色板必须在 hyperref 之前定义：hyperref 的 colorlinks 会直接引用这些名字。
\definecolor{NemoLink}{HTML}{1F4E79}   % 深蓝：内部链接
\definecolor{NemoCite}{HTML}{7A3E00}   % 赭石：文献引用
\definecolor{NemoUrl}{HTML}{0B6E4F}    % 深绿：外部链接
\definecolor{NemoRule}{HTML}{4A4A4A}   % 灰：表格线
\definecolor{NemoCardBg}{HTML}{F2F5F8} % 浅蓝灰：模块卡片底纹
\definecolor{NemoKeyBg}{HTML}{EDEDED}  % 浅灰：代码/键名底纹
\definecolor{NemoWarn}{HTML}{9C2A00}   % 暗红：边界与风险提示

% ---- 单位 ------------------------------------------------------------
\usepackage{siunitx}
\sisetup{
  detect-all,
  per-mode = symbol,
  range-phrase = {\text{--}},
  range-units = single,
  separate-uncertainty = true
}

% ---- 参考文献（全局唯一一份，样式为国标 GB/T 7714-2015）--------------
\usepackage[backend=biber,style=gb7714-2015,sorting=gb7714-2015]{biblatex}
\addbibresource{refs.bib}

% ---- 超链接 ----------------------------------------------------------
\usepackage[
  unicode=true,
  bookmarksnumbered=true,
  bookmarksopen=true,
  colorlinks=true,
  linkcolor=NemoLink,
  citecolor=NemoCite,
  urlcolor=NemoUrl,
  filecolor=NemoUrl,
  pdftitle={Huntian: GaN 射频 HEMT 轻量化 TCAD 内核},
  pdfsubject={模块化技术文档},
  pdfkeywords={GaN, HEMT, TCAD, Schroedinger-Poisson, drift-diffusion}
]{hyperref}

% ---- 智能交叉引用 ----------------------------------------------------
%  注意：这里\emph{不}加载 cleveref —— 它与 ctex 不兼容，第二遍编译会报
%  "Extra \endcsname"（最小复现已确认，且与 hyperref 无关）。
%  本文档改用自己的引用宏，见 preamble/refs.tex。

% ---- 索引 ------------------------------------------------------------
%  P0 用 texindy（UTF-8 + 英文排序规则）；中文条目一律带 ASCII 排序键
%  （写成 \index{3method@方法!...}），因此排序完全由文档控制，
%  与 makeindex / xindy 的选择无关。详见 §"索引与术语约定"。
\usepackage{imakeidx}
\makeindex[
  title={索引},
  intoc=true,
  program=texindy,
  options={-C utf8 -L english},
  columns=2,
  columnsep=1.2em
]

% =====================================================================
%  style.tex —— 版式、字体、页眉页脚、题注、交叉引用名称
%  本文件只放"长什么样"，不放"是什么"。
% =====================================================================

% ---- 字体 ------------------------------------------------------------
%  ctexbook 已按 fontset=windows 配好中文（宋体正文 / 黑体标题）。
%  西文与数学保持 Computer Modern 家族，与中文宋体搭配稳定，
%  且避免 unicode-math 带来的额外字体依赖。
\setmainfont{Latin Modern Roman}
\setsansfont{Latin Modern Sans}
\setmonofont{Latin Modern Mono}[Scale=0.92]

% ---- 行距与段落 ------------------------------------------------------
\linespread{1.35}
\setlength{\parskip}{0.35em}

% ---- 断行宽容度 ------------------------------------------------------
%  中英混排 + 等宽长标识符（physics.quantum.analytic_tw 之类）无法连字符
%  断开，容易产生 overfull box。提高容忍度并给足紧急伸展量，
%  让 TeX 用字间距吸收，而不是把字排出边界。
\tolerance=3000
\emergencystretch=5em
\hbadness=10000   % 欠满盒多为中英混排的正常现象，不再刷屏
%  溢出的报告阈值：2pt（≈0.7mm）以下不报。这是一个\emph{声明的容差}，
%  不是掩盖缺陷 —— 在窄表格列里塞等宽长标识符时，亚毫米级溢出
%  在视觉上不可辨，而它每次都会淹没日志里真正需要看的行。
%  当前文档在此阈值下为 0 个溢出盒。
\hfuzz=2pt

% ---- 页眉页脚 --------------------------------------------------------
\usepackage{fancyhdr}
\pagestyle{fancy}
\fancyhf{}
\fancyhead[L]{\small\itshape\nouppercase{\leftmark}}
\fancyhead[R]{\small\thepage}
\renewcommand{\headrulewidth}{0.4pt}
\renewcommand{\footrulewidth}{0pt}
\fancypagestyle{plain}{%
  \fancyhf{}%
  \fancyfoot[C]{\small\thepage}%
  \renewcommand{\headrulewidth}{0pt}%
}

% ---- 章标题 ----------------------------------------------------------
%  不引入 titlesec，也不碰 tocloft：ctexbook 的章样式与中文排版已协调，
%  改它们容易踩坑，而本文档的结构一致性靠 preamble/modulecard 保证。

% ---- 题注 ------------------------------------------------------------
\captionsetup{
  font=small,
  labelfont=bf,
  labelsep=quad,
  skip=6pt,
  hypcap=false        % 本文档用非浮动的 \captionof，hypcap 无意义
}
\captionsetup[table]{position=top}

% ---- 页眉高度 --------------------------------------------------------
%  fancyhdr 的 \headheight 警告：给足高度，避免它每页刷屏。
\setlength{\headheight}{15pt}

% ---- 表格线宽 --------------------------------------------------------
\setlength{\aboverulesep}{0.35ex}
\setlength{\belowrulesep}{0.35ex}
\heavyrulewidth=1.0pt
\lightrulewidth=0.5pt

% ---- 代码/键名 -------------------------------------------------------
%  \pid{} 用于模块 id、\fpath{} 用于文件路径：两者只含点、下划线、斜杠，
%  \emph{不含空格与花括号}，因此走 \path 使其可在标点处断行 ——
%  模块 id 与路径都很长，在窄表格列里必然溢出。
%  \akey{} 用 \detokenize：它的内容可能含空格与花括号
%  （如 layers[].{role, material}），而 \path 会吞掉空格，故不适用。
%  两者都要求内容写"裸"字符（裸下划线、裸花括号、裸百分号），
%  不要写 \_ 、\{ —— \detokenize 会把它们连反斜杠一起印出来。
%  ⚠ 用 \DeclareRobustCommand 而非 \newcommand：\path 是脆弱的（fragile），
%  直接放进 \caption 这类\emph{移动参数}会报
%  "\url used in a moving argument"。声明为 robust 之后，
%  写进 .lot/.toc 的是 \fpath{...} 本身，展开推迟到排版时，问题消失。
\DeclareRobustCommand{\pid}[1]{\texttt{\path{#1}}}
\DeclareRobustCommand{\fpath}[1]{\texttt{\path{#1}}}
\newcommand{\akey}[1]{\texttt{\detokenize{#1}}}

% ---- 术语首现 --------------------------------------------------------
%  \term{中文}{English}：首次出现时给出英文对照，之后只用中文。
\newcommand{\term}[2]{#1（#2）}

% ---- 边界与风险 ------------------------------------------------------
\newcommand{\boundary}[1]{%
  \par\smallskip\noindent
  \textcolor{NemoWarn}{\rule[-0.2ex]{2pt}{1.1em}}\hspace{0.6em}%
  \begin{minipage}[t]{\dimexpr\linewidth-1.2em\relax}%
  \small #1%
  \end{minipage}\par\smallskip
}

% ---- 交叉引用名称 ----------------------------------------------------
%  本文档不用 cleveref（与 ctex 不兼容，见 preamble/refs.tex 文件头）。
%  引用一律通过 preamble/refs.tex 中的类型化宏书写。

% ---- 目录深度 --------------------------------------------------------
\setcounter{tocdepth}{2}      % 收到 subsection
\setcounter{secnumdepth}{3}   % 编号到 subsubsection

% =====================================================================
%  math.tex —— 全局数学宏
%  规则：正文中的物理量一律走本文件的宏，禁止裸写下标；
%        符号表（frontmatter/notation.tex）与本文件必须一一对应。
% =====================================================================

% ---- 正体算子 --------------------------------------------------------
\DeclareMathOperator{\diag}{diag}
\DeclareMathOperator{\tridiag}{tridiag}
\DeclareMathOperator{\rank}{rank}
\DeclareMathOperator{\spann}{span}
\DeclareMathOperator*{\argmin}{arg\,min}
\DeclareMathOperator*{\argmax}{arg\,max}

% ---- 微分 ------------------------------------------------------------
\newcommand{\dd}{\mathrm{d}}                       % 直立的 d
\newcommand{\der}[2]{\frac{\dd #1}{\dd #2}}        % 全导数
\newcommand{\pder}[2]{\frac{\partial #1}{\partial #2}}

% ---- 常数与数 --------------------------------------------------------
\newcommand{\ii}{\mathrm{i}}                       % 虚数单位
\newcommand{\ee}{\mathrm{e}}                       % 自然对数底
\newcommand{\kb}{k_{\mathrm{B}}}                   % 玻尔兹曼常数
\newcommand{\kT}{k_{\mathrm{B}}T}                  % 热能
\newcommand{\epsO}{\varepsilon_{0}}                % 真空介电常数
\newcommand{\epsr}{\varepsilon_{\mathrm{r}}}       % 相对介电常数
%  面的等效相对介电常数。单独定义一个宏，因为它带上标：写成
%  \epsr^{\mathrm{face}}_i 会因 \epsr 自带下标而触发 "Double subscript"。
\newcommand{\epsface}{\varepsilon^{\mathrm{face}}}

%  ⚠ 网格节点索引必须走下面这三个宏（逗号下标形式）。
%  \Ec 与 \epsr 本身已经带下标，再写 \Ec_i 就是"同一个基上的两个下标"，
%  TeX 会报 Double subscript。这是本文件里最容易重犯的一类错误，
%  新增带角标的物理量宏时请一并在此登记。
\newcommand{\Eci}[1]{E_{\mathrm{c},#1}}                    % E_{c,i}
\newcommand{\Ecifl}[1]{E_{\mathrm{c},#1}^{\mathrm{full}}}  % E_{c,i}^{full}
\newcommand{\epsri}[1]{\varepsilon_{\mathrm{r},#1}}        % eps_{r,i}
\newcommand{\hb}{\hbar}                            % 约化普朗克常数

% ---- 物理量 ----------------------------------------------------------
\newcommand{\Ec}{E_{\mathrm{c}}}                   % 导带底
\newcommand{\Ev}{E_{\mathrm{v}}}                   % 价带顶
\newcommand{\Ef}{E_{\mathrm{F}}}                   % 费米能级
\newcommand{\Ei}{E_{i}}                            % 第 i 个亚带能级
\newcommand{\Ej}{E_{j}}                            % 第 j 个亚带能级（微扰求和用）
\newcommand{\ns}{n_{\mathrm{s}}}                   % 二维面密度
\newcommand{\nsu}{n_{\mathrm{s}}^{\max}}           % 面密度上界 sigma_eff/e
\newcommand{\meff}{m^{*}}                          % 有效质量
\newcommand{\mO}{m_{0}}                            % 自由电子质量
\newcommand{\phib}{\phi_{\mathrm{b}}}              % 肖特基势垒高度
\newcommand{\dEc}{\Delta E_{\mathrm{c}}}           % 导带带阶
\newcommand{\sigpol}{\sigma_{\mathrm{pol}}}        % 极化面电荷
\newcommand{\sigeff}{\sigma_{\mathrm{eff}}}        % 屏蔽后有效面电荷
\newcommand{\Jnn}{J_{n}}                           % 电子电流密度
\newcommand{\mumu}{\mu}                            % 迁移率
\newcommand{\NC}{N_{\mathrm{C}}}                   % 导带有效态密度
%  带有效质量修正与节点索引的导带态密度。同样必须整体定义：
%  写成 \NC\,_{\text{eff},i} 会因 \NC 自带下标而触发 Double subscript。
\newcommand{\NCeff}[1]{N_{\mathrm{C,eff},#1}}      % N_{C,eff,i}

% ---- 记号算子 --------------------------------------------------------
\newcommand{\avg}[1]{\left\langle #1 \right\rangle}   % 期望值
\newcommand{\abs}[1]{\left\lvert #1 \right\rvert}     % 绝对值
\newcommand{\norm}[1]{\left\lVert #1 \right\rVert}    % 范数
\newcommand{\T}{^{\mathsf{T}}}                        % 转置
\newcommand{\dg}{^{\dagger}}                          % Moore-Penrose 伪逆
\newcommand{\mat}[1]{\bm{#1}}                         % 矩阵

% ---- 单位 ------------------------------------------------------------
%  用法：\U{cm^{-2}}  ->  直立罗马体的单位
%  用 \ensuremath 包裹，使其在文本模式与数学模式中都能用
%  （表格与符号表里大量按文本模式书写，必须两种模式通吃）。
\newcommand{\U}[1]{\ensuremath{\,\mathrm{#1}}}

% =====================================================================
%  refs.tex —— 交叉引用宏
%
%  为什么不用 cleveref：
%    cleveref 与 ctex 不兼容。最小复现是
%        \documentclass{ctexbook}\usepackage{cleveref}
%        \begin{document}\cref{s1}\par\section{X}\label{s1}\end{document}
%    第二遍编译必报 "Extra \endcsname"。已确认与 hyperref 无关
%    （ctexbook+cleveref 已复现；book+hyperref+cleveref 正常）。
%    zref-clever 与 ctex 兼容，但不提供中文引用名，而本文档要求
%    "第 10 章""式 (3)""表 2"这样的中文形态。
%
%  因此本文档自带一组\emph{按类型区分}的引用宏，建立在 \hyperref + \ref*
%  之上：零宏包风险，输出形态完全可控，且链接照常可用。
%
%  纪律：正文中禁止写死"见 §4.2"这类字面编号，一律使用下列宏。
%        新增引用类型时在此处扩展，不要在正文里就地拼字符串。
% =====================================================================

%  \nref{label}{text} —— 带超链接的文本（\ref* 不产生嵌套链接）
\newcommand{\nref}[2]{\hyperref[#1]{#2}}

%  章：第 10 章
\newcommand{\crefchap}[1]{\nref{#1}{第\,\ref*{#1}\,章}}
%  节：第 10.2 节
\newcommand{\crefsec}[1]{\nref{#1}{第\,\ref*{#1}\,节}}
%  式：式 (3)
\newcommand{\crefeq}[1]{式~\eqref{#1}}
%  表：表 2
\newcommand{\creftab}[1]{\nref{#1}{表~\ref*{#1}}}
%  图：图 1
\newcommand{\creffig}[1]{\nref{#1}{图~\ref*{#1}}}
%  附录：附录 A
\newcommand{\crefapp}[1]{\nref{#1}{附录~\ref*{#1}}}
%  算法/代码块
\newcommand{\creflst}[1]{\nref{#1}{代码~\ref*{#1}}}

% =====================================================================
%  modulecard.tex —— 模块卡片与可标定参数表
%
%  这是"一个插件 = 一个模块章"在排版上的机械载体。
%  每个模块章的 §x.1 必须以 modulecard 开头，字段顺序固定：
%      类别 / 能力 / 依赖 artifact / 产出 artifact / 验证状态 / 源码
%  任何一章缺字段，或字段顺序不同，都视为模板违规（见 P6 校验脚本计划）。
%
%  ⚠ 实现约束：与 preamble/verify.tex 文件头所列完全相同。
%  核心一条：不用 \begin{tabularx}/\begin{table}/\multicolumn，
%  一律用不带环境记账的 \tabular ... \endtabular 原语 + \captionof。
%  真因（tabularx.sty:53 读 \@currenvir）见该文件头，勿回退。
% =====================================================================

% ---- 模块卡片（非浮动，就地排版）-------------------------------------
%  用法：\begin{modulecard}{module_id}{module_version} ... \end{modulecard}
\newenvironment{modulecard}[2]{%
  \par\medskip\noindent
  \textbf{模块卡片}\quad\pid{#1}\quad{\small(#2)}%
  \par\nobreak\smallskip\noindent
  \begingroup
  \setlength{\tabcolsep}{6pt}%
  \renewcommand{\arraystretch}{1.22}%
  \tabular{@{}
    >{\bfseries\raggedright\arraybackslash}p{30mm}
    p{\dimexpr\textwidth-30mm-2\tabcolsep\relax}@{}}
  \toprule
}{%
  \bottomrule
  \endtabular
  \endgroup
  \par\medskip
}

%  卡片行
\newcommand{\cardrow}[2]{#1 & #2 \\}
%  卡片内分隔线
\newcommand{\cardsep}{\midrule}

% ---- 可标定参数表（非浮动）-------------------------------------------
%  用法：\begin{paramtable}{题注}{label} ... \end{paramtable}
%  列：参数名 | 默认值 | 单位 | 含义与取值依据
\newenvironment{paramtable}[2]{%
  \par\medskip\noindent
  \begingroup
  \small
  \captionof{table}{#1}\label{#2}%
  \par\nobreak\smallskip\noindent
  \setlength{\tabcolsep}{5pt}%
  \renewcommand{\arraystretch}{1.18}%
  \tabular{@{}p{38mm} p{20mm} p{16mm}
    p{\dimexpr\textwidth-74mm-6\tabcolsep\relax}@{}}
  \toprule
  \textbf{参数} & \textbf{默认值} & \textbf{单位} & \textbf{含义与取值依据}\\
  \midrule
}{%
  \bottomrule
  \endtabular
  \par\endgroup
  \medskip
}

%  参数行
\newcommand{\prow}[4]{#1 & #2 & #3 & #4 \\}
\newcommand{\calib}{\textsuperscript{\textcolor{NemoCite}{$\ast$}}}
\newcommand{\calibnote}{%
  \par\smallskip\noindent\small
  \textsuperscript{\textcolor{NemoCite}{$\ast$}}~标记者为\emph{可标定量}：
  其默认值未经 L2 基准（Hall / XRD / C--V 实测）验证，使用前必须自行标定。%
}

% =====================================================================
%  verify.tex —— 验证基准表与边界清单
%
%  每个模块章的 §x.5 用 veriftable，§x.6 用 bounds。
%  三列固定为：基准 | 判据 | 实测。
%  "判据"列写测试断言的门槛，"实测"列写文档记录的实际观测值 —— 两者
%  必须同时给出，否则读者无法判断余量，也无法在版本演化后复查。
%
%  ⚠ 实现约束（真因已查明，勿回退）：
%  本环境在 begin 代码里开表、在 end 代码里闭表，配对跨越了
%  \newenvironment 的边界。凡是"配对敏感"的构造一律不能用：
%
%    * \begin{tabularx}...\end{tabularx} —— 不可用。
%      tabularx.sty:53 写的是
%          \def\tabularx#1{\edef\TX@{\@currenvir}...\TX@get@body}
%      它把"当前环境名"存进 \TX@，然后一路扫描直到 \end{<那个名字>}。
%      在自定义环境里 \@currenvir 是 veriftable 而不是 tabularx，
%      于是它去找 \end{veriftable}，把整章内容吞进表格。
%      即便改用 \tabularx 原语也一样，因为它读的是 \@currenvir。
%
%    * \begin{table} / \@float —— 同样有环境边界记账问题。
%
%  结论：一律使用**不带环境记账的底层原语** \tabular ... \endtabular，
%  配合非浮动的 \captionof（caption 宏包）。列宽用 \dimexpr 精确算到
%  \textwidth，因此不需要 tabularx 的自适应 X 列。
%  附带好处：表就地排版，永远跟着它所属的段落走，不会漂到别页。
% =====================================================================

% ---- 表注 ------------------------------------------------------------
%  ⚠ 表注必须是表格里的一\emph{行}，不能是"表格末尾的一条空指令"。
%  写成 \def 存文本、表后再印（旧做法）会踩到 TeX 的一个死角：
%  在 \tabular 里，\\ 之后紧跟的下一个记号若不是 \noalign/\omit，
%  LaTeX 就认为新行已经开始，于是随后的 \bottomrule 报
%  "Misplaced \noalign"。
%  因此这里让 \vnote 自己排成一行、并自己以 \\ 收尾：
%  它既是最后一行，也仍然是"完整的一行"，\bottomrule 便落在合法的
%  \cr 之后。普通 tabular 里 \multicolumn 是安全的 —— 之前出事的是
%  tabularx 的 X 列（见 modulecard.tex 文件头），两者不要混淆。
\newcommand{\vnote}[1]{%
  \multicolumn{3}{@{}p{\textwidth}@{}}{\footnotesize #1}\\}

% ---- 验证基准表（非浮动）---------------------------------------------
%  用法：\begin{veriftable}{题注}{label} ... \end{veriftable}
\newenvironment{veriftable}[2]{%
  \par\medskip\noindent
  \begingroup
  \small
  \captionof{table}{#1}\label{#2}%
  \par\nobreak\smallskip\noindent
  \setlength{\tabcolsep}{5pt}%
  \renewcommand{\arraystretch}{1.18}%
  \tabular{@{}p{\dimexpr\textwidth-72mm-4\tabcolsep\relax} p{34mm} p{38mm}@{}}
  \toprule
  \textbf{基准} & \textbf{判据} & \textbf{实测}\\
  \midrule
}{%
  \bottomrule
  \endtabular
  \par\endgroup
  \medskip
}

%  基准行：基准名 | 判据 | 实测
\newcommand{\vrow}[3]{#1 & #2 & #3 \\}

%  分组标题行：不用 \multicolumn（tabular 里虽可用，但保持与列宽一致）
\newcommand{\vgroup}[1]{\textit{#1} & & \\}

% ---- 边界清单 --------------------------------------------------------
%  用法：\begin{bounds} \item ... \end{bounds}
%  注意：不要把它命名为 limits —— \limits 是 TeX 原语，会撞名。
\newenvironment{bounds}{%
  \begin{itemize}[leftmargin=1.7em,itemsep=0.28em,topsep=0.35em,parsep=0pt]
}{%
  \end{itemize}
}

%  "已知边界"小节标题的统一样式
\newcommand{\limitsheading}{\medskip\noindent\textbf{适用边界与已知局限}\par\nobreak\smallskip}


% ---- 全局元数据 ------------------------------------------------------
\title{\bfseries Huntian\\[0.3em]
       \large GaN 射频 HEMT 轻量化 TCAD 内核\\[0.15em]
       \large 模块化技术文档}
\author{Huntian 项目组}
\date{文档版本 0.1（P0 骨架）\\ 2026 年 10 月}

\begin{document}

% =====================================================================
\frontmatter
% =====================================================================
\maketitle
\tableofcontents

% =====================================================================
%  abstract.tex —— 摘要
% =====================================================================

\chapter*{摘要}
\addcontentsline{toc}{chapter}{摘要}
\markboth{摘要}{摘要}

本文档系统记述 \textbf{Huntian} —— 一套面向 GaN 射频高电子迁移率晶体管
（high electron mobility transistor, HEMT）的轻量化、模块化工艺计算机辅助设计
（technology computer-aided design, TCAD）内核 —— 的计算原理、离散格式与验证证据。

Huntian 的架构主张是\emph{模块即插件}：每一个物理模型、每一个数值服务都以独立
插件的形式挂在注册中心上，彼此之间不通过代码调用而只通过命名数据通道
（artifact）交换数据。因此，本文档采用与之同构的组织方式：\textbf{一个插件对应一个
模块章}，各章自洽地给出该模块的控制方程、离散化、算法、参数依据与验证基准，
而架构层与验证方法学集中在前置的总论部分。

每章遵循固定的六段式骨架：模块契约、控制方程与物理假设、离散化与数值算法、
可标定参数、验证与基准、误差来源与适用边界。这种一致性并非形式主义 ——
它使读者可以在任意两章之间横向比较\emph{同一类证据}，也使文档本身可以被机械化
校验（章首的模块标识与代码中的 \texttt{MANIFEST} 逐字段比对）。

本文档对验证证据采取严格的二分表述：凡进入"实测"列的数字，均可在仓库的测试
套件或基准记录中回查；凡未经验证者，一律在"适用边界"一节明确列出，不作粉饰。
截至本文档当前版本，仓库测试套件共 \textbf{426} 个用例，全部通过。

\vspace{1em}
\noindent\textbf{关键词：} GaN HEMT；TCAD；自洽 Schrödinger--Poisson 方程；
二维电子气；漂移扩散；Scharfetter--Gummel 离散；Newton 法；精确 Jacobian；
插件架构；可复现性

\cleardoublepage

% =====================================================================
%  notation.tex —— 全局符号表（唯一真源）
%
%  纪律：正文中出现的每一个符号必须在此登记；此表出现的每一个符号
%        必须至少在一章中被使用。preamble/math.tex 的宏与本表一一对应。
% =====================================================================

\chapter*{符号表}
\addcontentsline{toc}{chapter}{符号表}
\markboth{符号表}{符号表}

本表为全文符号的唯一真源。各模块章一律沿用本表定义，不得就地重新定义符号；
若某模块确需引入新符号，须同时在本表登记。

\vspace{0.5em}
\small
\begin{longtable}{@{}p{24mm} p{72mm} p{30mm} p{22mm}@{}}
\toprule
\textbf{符号} & \textbf{含义} & \textbf{单位} & \textbf{首次出现}\\
\midrule
\endfirsthead
\toprule
\textbf{符号} & \textbf{含义} & \textbf{单位} & \textbf{首次出现}\\
\midrule
\endhead
\bottomrule
\endfoot

\multicolumn{4}{@{}l@{}}{\itshape 基本常数（CODATA 2018）}\\
$q$                     & 元电荷 & \U{C} & \S\ref{sec:sp1d-model}\\
$\kb$                   & 玻尔兹曼常数 & \U{J/K} & \S\ref{sec:sp1d-model}\\
$\hb$                   & 约化普朗克常数 & \U{J\cdot s} & \S\ref{sec:sp1d-model}\\
$\epsO$                 & 真空介电常数 & \U{F/m} & \S\ref{sec:sp1d-model}\\
$\mO$                   & 自由电子质量 & \U{kg} & \S\ref{sec:sp1d-model}\\
$\kT$                   & 热能，$T$ 为晶格温度 & \U{eV} & \S\ref{sec:sp1d-model}\\
\midrule
\multicolumn{4}{@{}l@{}}{\itshape 几何与网格}\\
$z$                     & 竖直坐标，自栅极指向体内为正 & \U{m} & \S\ref{sec:sp1d-geometry}\\
$d$                     & 势垒层厚度（金属/势垒界面到异质界面） & \U{m} & \S\ref{sec:sp1d-geometry}\\
$L$                     & 求解域总深度（异质界面到背面） & \U{m} & \S\ref{sec:sp1d-geometry}\\
$h_i$                   & 第 $i$ 个面两侧节点间距 & \U{m} & \S\ref{sec:sp1d-fv}\\
$w_i$                   & 第 $i$ 个节点的有限体积控制体宽度 & \U{m} & \S\ref{sec:sp1d-fv}\\
$c_i$                   & 面的介电系数，$c_i=\epsface_i\epsO/h_i$ & \U{F/m^2} & \S\ref{sec:sp1d-fv}\\
\midrule
\multicolumn{4}{@{}l@{}}{\itshape 静电与量子}\\
$\Ec(z)$                & 导带底能量（以 $\Ef$ 为零点） & \U{eV} & \S\ref{sec:sp1d-model}\\
$\psi_i(z)$             & 第 $i$ 个亚带的包络波函数 & \U{m^{-1/2}} & \S\ref{sec:sp1d-model}\\
$\Ei$                   & 第 $i$ 个亚带的本征能量 & \U{eV} & \S\ref{sec:sp1d-model}\\
$\Ef$                   & 费米能级（平衡态取为零点） & \U{eV} & \S\ref{sec:sp1d-model}\\
$\meff(z)$              & 有效质量（沿 $z$ 方向） & \U{kg} & \S\ref{sec:sp1d-model}\\
$m^{*}_i$               & 第 $i$ 个亚带的态密度有效质量 & \U{kg} & \S\ref{sec:sp1d-model}\\
$g_v$                   & 导带谷简并度 & --- & \S\ref{sec:sp1d-model}\\
$n_i$                   & 第 $i$ 个亚带的二维占据密度 & \U{m^{-2}} & \S\ref{sec:sp1d-occupation}\\
$n(z)$                  & 体电子密度 & \U{m^{-3}} & \S\ref{sec:sp1d-occupation}\\
$\ns$                   & 二维电子气面密度，$\ns=\sum_i n_i$ & \U{m^{-2}} & \S\ref{sec:sp1d-occupation}\\
$\nsu$                  & 面密度上界，$\nsu=\sigeff/q$ & \U{m^{-2}} & \S\ref{sec:sp1d-occupation}\\
$\sigeff$               & 屏蔽后的有效极化面电荷 & \U{C/m^2} & \S\ref{sec:sp1d-model}\\
$\phib$                 & 金属/势垒界面肖特基势垒高度 & \U{eV} & \S\ref{sec:sp1d-bc}\\
$V_g$                   & 栅极偏压 & \U{V} & \S\ref{sec:sp1d-bc}\\
$\dEc$                  & 异质界面导带带阶 & \U{eV} & \S\ref{sec:sp1d-bc}\\
$\epsr$                 & 相对介电常数 & --- & \S\ref{sec:sp1d-fv}\\
\midrule
\multicolumn{4}{@{}l@{}}{\itshape 算子与离散}\\
$\mat{L}$               & 有限体积 Poisson 算子 & --- & \S\ref{sec:sp1d-fv}\\
$\mat{A}$               & 对称化 BenDaniel--Duke 三对角 Hamilton 矩阵 & --- & \S\ref{sec:sp1d-bdd}\\
$\mat{W}$               & 控制体宽度对角矩阵 & --- & \S\ref{sec:sp1d-fv}\\
$\mat{J}$               & 耦合残差的 Jacobian & --- & \S\ref{sec:sp1d-jacobian}\\
$\mat{R}$               & 耦合残差向量 & --- & \S\ref{sec:sp1d-jacobian}\\
$\dg$                   & Moore--Penrose 伪逆 & --- & \S\ref{sec:sp1d-resolvent}\\
$N$                     & 竖直网格节点数 & --- & \S\ref{sec:sp1d-fv}\\
$\delta_{\mathrm{FD}}$  & 中心差分步长 & --- & \S\ref{sec:sp1d-verification}\\

\end{longtable}
\normalsize

\cleardoublepage

% =====================================================================
%  conventions.tex —— 全局约定
%  这些约定对全文所有模块章生效，各章不再重复声明。
% =====================================================================

\chapter*{全局约定}
\addcontentsline{toc}{chapter}{全局约定}
\markboth{全局约定}{全局约定}

\section*{单位制与坐标}

全文采用国际单位制（SI），数值结果按器件工程惯例换算并标注单位。竖直坐标 $z$
的取法在各章统一为：\emph{自金属/半导体界面指向半导体体内为正}。这一取向是
约定而非物理，但它决定了导带带阶与极化面电荷的符号；跨章引用时以此为准。

势（potential）以伏特计，能量以电子伏特计。在本文档中，一处静电势
$\varphi$ 与相应的能量 $q\varphi$ 在数值上相等（$q\varphi$ 以 eV 计），
故正文中二者常直接互用而不显式写出 $q$；凡出现此类互用时，均以单位
（\U{V} 对 \U{eV}）区分。

\section*{能量零点}

平衡态下费米能级取为能量零点，即 $\Ef=0$。薛定谔方程的本征能量 $\Ei$
与导带边 $\Ec(z)$ 均以该零点为基准，因此束缚态表现为 $\Ei<0$，
金属/势垒界面的导带边为 $\Ec(0)=q(\phib-V_g)$。

\section*{符号与命名}

\begin{itemize}[leftmargin=1.7em,itemsep=0.25em,topsep=0.3em]
  \item \textbf{物理量}：沿用前面的\emph{符号表}，全局唯一。
  \item \textbf{模块标识}：形如 \pid{physics.quantum.sp1d} 的点分小写字符串，
        与代码中 \texttt{MANIFEST.module\_id} 逐字符一致，排版为等宽体。
  \item \textbf{数据通道键}（artifact key）：形如 \akey{polarization}、
        \akey{quantum}，与代码 \texttt{ctx.put}/\texttt{ctx.require} 的实参一致。
  \item \textbf{参数名}：形如 \akey{tolerance_eV}，与插件 \texttt{ParamSpec}
        的声明一致；可标定量以 $\ast$ 标记。
  \item \textbf{文件路径}：相对仓库根目录书写，排版为等宽体。
\end{itemize}

\section*{术语}

\term{二维电子气}{two-dimensional electron gas, 2DEG}、
\term{高电子迁移率晶体管}{HEMT}、
\term{漂移扩散}{drift--diffusion, DD}、
\term{有限体积}{finite volume, FV} 等术语在首次出现时于括号内给出原文，
其后一律使用中文。物理量与算子的拉丁字母记号不译。

\section*{验证证据的表述纪律}

本文档对"验证到什么程度"作严格二分，全文一致：

\begin{itemize}[leftmargin=1.7em,itemsep=0.25em,topsep=0.3em]
  \item \textbf{判据}列记录\emph{测试断言的门槛}。它是回归门禁，不随版本漂移。
  \item \textbf{实测}列记录\emph{文档编写时实际观测到的值}。它带有效数字，
        可回查，但会随版本演化而变；引用时必须连同模块版本一并引用。
  \item 仅有一阶吻合、趋势正确或无独立基准可对照者，一律进入
        \emph{适用边界与已知局限}，不得写入"实测"列。
  \item 未标注验证状态的数值，默认视为\textbf{未验证}。
\end{itemize}

模块的验证状态取三个值之一，与代码 \texttt{ModuleManifest.verification} 一致：
\texttt{partially-verified}（存在闭式解或收敛阶证据，但绝对数值未标定）、
\texttt{verified}（在独立基准上标定通过）、\texttt{unverified}（占位实现）。

\cleardoublepage

% =====================================================================
%  outline.tex —— 文档结构总览
%  这一章是本项目"模块化"主张在文档侧的自述：9 个插件 -> 9 章，一一对应。
%  同时诚实标注当前成文进度。
% =====================================================================

\chapter*{文档结构总览}
\addcontentsline{toc}{chapter}{文档结构总览}
\markboth{文档结构总览}{文档结构总览}

本文档与代码仓库的 \fpath{huntian/plugins/} 目录一一对应：\textbf{一个插件一个模块章}。
模块章的编号、文件名与模块标识的对应关系是稳定的契约，后续版本只增不改。

\vspace{0.5em}
\small
\begin{longtable}{@{}c l l l c@{}}
\toprule
\textbf{章} & \textbf{模块标识} & \textbf{版本} & \textbf{类别} & \textbf{状态}\\
\midrule
\endfirsthead
\toprule
\textbf{章} & \textbf{模块标识} & \textbf{版本} & \textbf{类别} & \textbf{状态}\\
\midrule
\endhead
\bottomrule
\endfoot

\multicolumn{5}{@{}l@{}}{\itshape 第一部分\ \ 总论}\\
1  & 引言与文档范围                       & ---     & ---        & 待写\\
2  & 系统架构：模块即插件、服务即插件       & ---     & ---        & 待写\\
3  & 数据流与 artifact 契约               & ---     & ---        & 待写\\
4  & 验证方法学                           & ---     & ---        & 待写\\
5  & 模块索引与依赖关系                    & ---     & ---        & 待写\\
\midrule
\multicolumn{5}{@{}l@{}}{\itshape 第二部分\ \ 网格与材料}\\
6  & \pid{mesh.structured2d}             & 0.1.0   & mesh       & \textbf{已成文}\\
7  & \pid{materials.dft}                 & 0.1.0   & physics    & \textbf{已成文}\\
\midrule
\multicolumn{5}{@{}l@{}}{\itshape 第三部分\ \ 极化与量子}\\
8  & \pid{physics.polarization.ambacher} & 0.1.0   & physics    & \textbf{已成文}\\
9  & \pid{physics.quantum.analytic_tw}  & 0.1.0   & physics    & \textbf{已成文}\\
10 & \pid{physics.quantum.sp1d}          & 1.1.1   & physics    & \textbf{已成文}\\
\midrule
\multicolumn{5}{@{}l@{}}{\itshape 第四部分\ \ 输运与提取}\\
11 & \pid{transport.dd}                  & 0.1.0   & transport  & \textbf{已成文}\\
12 & \pid{extract.fom}                   & 0.1.0   & extract    & \textbf{已成文}\\
\midrule
\multicolumn{5}{@{}l@{}}{\itshape 第五部分\ \ 数值服务模块}\\
13 & \pid{nonlinear.newton}              & 1.0.0   & solver     & \textbf{已成文}\\
14 & \pid{linear.blocktridiag}           & 0.1.0   & solver     & \textbf{已成文}\\
\midrule
\multicolumn{5}{@{}l@{}}{\itshape 附录}\\
A  & 闭式解基准汇总                       & ---     & ---        & 待写\\
B  & 可标定参数总表                       & ---     & ---        & 待写\\
C  & 版本、指纹与 artifact 契约矩阵        & ---     & ---        & 待写\\
D  & 已知边界与欠账清单                    & ---     & ---        & 待写\\
\midrule
\multicolumn{5}{@{}l@{}}{\itshape 全局部件}\\
--- & 参考文献                            & ---    & ---        & 成文\\
--- & 索引                                & ---    & ---        & 成文\\

\end{longtable}
\normalsize

\vspace{0.5em}
\noindent
\textbf{关于本阶段的编号。}第五部分之前的章节尚未成文，但\emph{章号已经锁定}：
本文档当前版本只印出第 10 章，其章号即为 10 而非 1。这样做的目的是让章号成为
稳定契约 —— 后补的第 6--9 章不会使第 10 章的交叉引用失效。待 P1--P3 铺开后，
目录自会填满，而 \crefchap{ch:sp1d} 的编号不会再变。

\cleardoublepage


% =====================================================================
\mainmatter
% =====================================================================

% ---------------------------------------------------------------------
%  第一部分  总论
%  至此没有前向标签了：19 章全部由真实章节产生编号。
%  frontmatter/future_labels.tex 已随之删除。
% ---------------------------------------------------------------------
\part{总论}
% 总论第 1 章 —— 引言与文档范围
% \futurechapter{ch:intro} 已在新章启用时从 future_labels.tex 删除

\chapter{引言与文档范围}
\label{ch:intro}

\section{项目定位}

Huntian 是一套面向 GaN 射频高电子迁移率晶体管（HEMT）的\textbf{轻量化、
模块化工艺计算机辅助设计（TCAD）内核}。它的目标不是替代商业 TCAD 的
全部功能，而是回答一个更窄、也更可验证的问题：

\begin{quote}
\emph{能不能用一套\emph{核心零第三方依赖}的代码，把 GaN HEMT 从极化电荷
到电流--电压曲线的物理链完整走通，并且让这条链上的每一环都可被独立替换、
独立验证、独立追责？}
\end{quote}

这个问题的形态决定了本文档的形态。既然每一环都可替换，那么每一环都值得
单独成章来讲清它的方程、它的离散、它的参数依据，以及\textbf{它凭什么被相信}。

\subsection{量化目标}

方案给出的量化目标见 \creftab{tab:intro-goals}。把它们列在正文开头是有
用意的：本文档全篇反复区分"格式可信"与"数值可信"，而这张表正是那条分界线
的量化形式 —— 表里的每一项都需要\emph{实测标定}才能宣称达成，
而截至本文档当前版本，\textbf{没有一项已经通过实测标定}。

\begin{veriftable}{方案制定的量化精度目标（尚未实测标定）}{tab:intro-goals}
\vrow{二维电子气面密度 $\ns$}{$\pm5\%$（对霍尔测量）}{未标定}
\vrow{阈值电压 $V_{\text{th}}$}{$\pm0.15$\U{V}}{未标定}
\vrow{饱和电流 $I_{d\max}$ / 跨导 $g_m$}{$\pm10\%$}{未标定}
\vrow{$f_T$ / $f_{\max}$}{$\pm10$--$15\%$}{未实现（无 S 参数链）}
\vrow{电流坍塌幅度}{$\pm20\%$}{未实现（无陷阱模型）}
\vrow{$P_{\text{out}}$ / PAE}{$\pm10$--$15\%$}{未实现（无大信号链）}
\vrow{单偏置点求解时间（二维，$\sim2\times10^{4}$ 节点）}{$<1$\U{s}}{见 \crefchap{ch:transport_dd}}
\vrow{模块热插拔时间}{$<1$\U{s}，无需重启进程}{通过（见 \crefchap{ch:architecture}）}
\end{veriftable}

\subsection{五条设计原则}

\begin{enumerate}[leftmargin=1.9em,itemsep=0.35em,topsep=0.4em]
  \item \textbf{专用优于通用。} 只做 GaN 射频 HEMT，只做单极 $n$ 型电子输运，
        以二维截面为主。通用性不是目标，可验证性才是。
  \item \textbf{模型即插件。} 一切物理模型、求解器、后端、适配器都走统一接口，
        可插拔、可版本化、可独立发布。
  \item \textbf{数字只能来自求解器。} 任何自动化层（含语言模型）只负责翻译、
        规划、诊断与解释，\emph{绝不直接输出物理数值}。
  \item \textbf{误差可追溯。} 每个结果携带参数快照、模块版本、网格与求解器设置，
        完全可复现。
  \item \textbf{渐进增强。} 先在 Python 中原型验证物理，再把热路径重写到
        高性能内核；性能与正确性分离推进。
\end{enumerate}

第 3 条在本文档中有具体体现：本章之后的每一章，
凡是进入"实测"列的数字，都必须能追到一次具体的运行或一条具体的断言。
这条纪律不是文风偏好，而是第 3 条原则在文档层面的落实。

\section{文档的组织方式}
\label{sec:intro-organisation}
\index{1concept@概念!one-plugin-one-chapter@一个插件一个模块章}

本文档与代码仓库的 \fpath{huntian/plugins/} 目录\textbf{一一对应}：

\begin{center}
\small
\begin{tabular}{@{}p{0.30\textwidth} p{0.60\textwidth}@{}}
\toprule
\textbf{代码侧} & \textbf{文档侧}\\
\midrule
\fpath{huntian/plugins/<name>/} 一个目录 & \fpath{modules/<name>.tex} 一章\\
\texttt{MANIFEST} 里的字段 & 章首的\emph{模块卡片}（字段顺序固定）\\
模块的残差/Jacobian 代码 & §x.2 控制方程、§x.3 离散化、§x.4 求解算法\\
\fpath{tests/test_<name>.py} & §x.5 验证与基准（判据 \& 实测两列）\\
模块说明书里的"诚实边界" & §x.6 误差来源与适用边界\\
\texttt{ctx.require}/\texttt{ctx.put} 实参 & §x.7 接口契约表\\
\bottomrule
\end{tabular}
\end{center}

\subsection{七段式骨架}

每一模块章遵循同一副骨骼：\textbf{模块契约、控制方程与物理假设、离散化与
公式化、求解算法与实现、验证与基准、误差来源与适用边界、接口契约与小结}。
其中第 3、4 两节的标题会按模块类型 specialize（数值型模块写"离散化"
与"求解算法"，解析或数据型模块写"公式化"与"合并流程"），但\textbf{两节
都必须在}。

这种一致性不是形式主义。它使读者可以在任意两章之间\textbf{横向比较同一类
证据}：第九章的"判据 $>$ 实测"与第十章的"判据 $>$ 实测"说的是同一件事，
因此可以并排读。它也使文档本身可以被机械化校验
（章首的模块标识与代码中的 \texttt{MANIFEST} 逐字段比对）——
那属于后续版本的工程化工作，见 \crefchap{ch:module-index} 的说明。

\subsection{章号是冻结契约}

本文档的章号不随成文先后变化：第 6--14 章是九个模块章，
第 1--5 章是总论，附录为 A--D。在模块章尚未写成的阶段，
章号由一份\emph{前向标签}文件声明，它只推进章计数器、不排版任何内容 ——
因此后补的章节不会使已有的交叉引用失效。

这条约定的代价写在下面，供后续维护者参考：\textbf{每写成一章，
必须同时做两件事} ——在根文件里放入真正的 \texttt{\textbackslash input}，
并从标签文件里删除对应的声明行；只做一半会导致章号重复计数。
构建日志中每一条 \texttt{[huntian-doc] forward-label} 行都对应一个
"已编号但尚未成文"的章，这是本阶段的审计凭据。

\section{读者指南}

\begin{center}
\small
\begin{tabular}{@{}p{0.26\textwidth} p{0.64\textwidth}@{}}
\toprule
\textbf{如果你想知道} & \textbf{建议路径}\\
\midrule
整体怎么搭起来的 &
  \crefchap{ch:architecture} $\to$ \crefchap{ch:dataflow} $\to$
  \crefchap{ch:module-index}\\
某个模块在算什么 &
  直接翻对应章；每章 §x.2 是方程，§x.5 是"凭什么信"\\
某段代码为什么这么写 &
  对应章的 §x.3 与 §x.4；跨模块的取舍见 \crefchap{ch:dataflow}\\
数字可信到什么程度 &
  \crefchap{ch:verification-methodology} 讲判据体系；
  每章 §x.6 列该模块自己的边界\\
还欠什么 &
  \crefapp{app:debts}（附录 D）汇总全部已知边界与欠账\\
\bottomrule
\end{tabular}
\end{center}

\section{本版状态与诚实声明}
\label{sec:intro-status}

本文档覆盖的状态可以用三句话概括。

\textbf{其一，骨架与九个模块章已成文。} 注册中心、上下文、Deck、流水线
四条地基全部落地；九个插件各自成章，合计 19 章 4 附录。

\textbf{其二，格式可信与数值可信是两件事，本文档始终分开陈述。}
已经验证到闭式解或收敛阶的，是\emph{离散格式}与\emph{导数}；
尚未验证的，是\emph{绝对数值} —— 后者受制于未经标定的材料常数与迁移率。
凡引用一个数字，本文档同时给出它的\emph{判据}（测试断言的门槛）
与\emph{实测}（编写时观测到的值），二者缺一不可。

\textbf{其三，本文档本身也接受同一条纪律。} 参考文献逐条经 DOI 核验并
记录核验日期；未能核验的一律不进入正文。核验过程中发现的、仓库内的
引用错误亦如实记录（见 \crefchap{ch:dft} 与参考文献条目 \texttt{ambacher2002}），
而不是沿袭。

\cleardoublepage
                 % 第 1 章   引言与文档范围
% 总论第 2 章 —— 系统架构

\chapter{系统架构：模块即插件、服务即插件}
\label{ch:architecture}
\index{1concept@概念!plugin architecture@插件架构}

本章记述承载全部物理模型的那套骨架。它回答的问题是：
\emph{为什么九个功能差异极大的模块（从网格生成到稠密线性代数）
能共享同一套装配、调度与校验机制}。

答案由三条地基构成，它们各自只做一件事，且都\emph{不含任何物理}：

\begin{center}
\small
\begin{tabular}{@{}p{0.20\textwidth} p{0.30\textwidth} p{0.42\textwidth}@{}}
\toprule
\textbf{地基} & \textbf{职责} & \textbf{不做什么}\\
\midrule
注册中心 & 发现、版本、依赖、参数校验、热插拔 & 不调度、不执行\\
SimContext & 命名数据通道、溯源、指纹 & 不解释数据内容\\
流水线 & 由 manifest 推导执行顺序 & 不硬编码任何模块\\
\bottomrule
\end{tabular}
\end{center}

\section{模块契约}
\label{sec:arch-contract}
\index{1concept@概念!manifest@模块清单}

一个模块就是一个实现了 \texttt{IModule} 的类，外加一份\textbf{声明}：

\begin{center}
\small
\begin{tabular}{@{}p{0.24\textwidth} p{0.66\textwidth}@{}}
\toprule
\textbf{字段} & \textbf{含义与它在本文档中的作用}\\
\midrule
\akey{module_id} & 点分小写标识，至少两段（如 \pid{physics.quantum.sp1d}）。
  与章首模块卡片逐字符一致\\
\akey{module_version} & 语义化版本。数值行为改变必须升版本（见 \S\ref{sec:arch-fingerprint}）\\
\akey{api_version} & 接口版本；主版本不匹配时\emph{拒绝加载}\\
\akey{category} & \akey{mesh} / \akey{physics} / \akey{transport} / \akey{extract} / \akey{solver}\\
\akey{capabilities} & 能力标签，供槽位契约检查（如 \texttt{quantum}、\texttt{linear}）\\
\akey{requires} & 需要的 artifact 键。\textbf{这是调度边的唯一来源}\\
\akey{produces} & 产出的 artifact 键。服务模块此字段为空\\
\akey{parameters} & 参数规格表：类型、默认值、范围、单位、是否可标定、出处\\
\akey{verification} & \texttt{verified} / \texttt{partially-verified} / \texttt{unverified}\\
\akey{tags}, \akey{author}, \akey{description} & 元数据\\
\bottomrule
\end{tabular}
\end{center}

三个字段值得单独说明，因为它们在本文档中反复出现。

\textbf{（i）\akey{requires}/\akey{produces} 同时是文档结构。}
本文档第 3 章的 artifact 契约矩阵就是从这两列生成的；
它也解释了为什么 \crefchap{ch:dft} 的模块"依赖 \akey{structure} 却从不读它"——
那条依赖声明的唯一作用是\emph{把自己排到极化与量子之前}。

\textbf{（ii）\akey{verification} 是一个\emph{三值}标签，不是一个百分比。}
模块说明书里没有"完成度 80\%"这种说法：要么有闭式解或收敛阶证据
（\texttt{partially-verified}），要么已在独立基准上标定通过
（\texttt{verified}），要么是占位实现（\texttt{unverified}）。
本文档各章的卡片沿用同一取值。

\textbf{（iii）\akey{parameters} 里 \akey{calibratable} 的含义是"允许被拟合"。}
\crefchap{ch:dft} 的模块\textbf{没有任何}可标定参数，且这一条被测试固化为断言；
其理由是方法论上的：数值方法的选择（截断能、$k$ 点密度）不应该被
自动拟合循环挪动，否则标定与收敛性会混为一谈。

\section{注册中心}
\label{sec:arch-registry}

\subsection{发现与加载}

注册中心扫描 \fpath{plugins/} 目录（支持\emph{多文件插件包}，
包内可用相对导入），或按点分导入路径加载。一条重要的工程性质是：
\textbf{单个插件损坏不会中断整轮发现} —— 一个语法错误的插件只会让它
自己缺席，其余模块照常可用。对由自动化层驱动的系统，这条性质是必要的：
一次失败的插件生成不应该让整个工作台不可用。

\subsection{版本：共存、钉住与语义}
\label{sec:arch-version}
\index{1concept@概念!semantic versioning@语义化版本}

同一 \akey{module_id} 可以\textbf{同时注册多个版本}，旧版本不会被覆盖。
这不是可有可无的：\crefchap{ch:sp1d} 的 Jacobian 精确化把模块从
1.1.0 升到 1.1.1，而验证过的 Deck 钉住的是 \texttt{\^{}1.0.0}。
两个版本必须能共存，旧结果才永远可复现。

依赖范围支持五种写法，其语义见 \creftab{tab:arch-spec}。

\begin{veriftable}{版本范围语法}{tab:arch-spec}
\vrow{\texttt{>=1.0,<2.0}}{交集（逗号即 AND）}{两个条件同时满足}
\vrow{\texttt{\^{}1.2.3}}{允许 $1.x$ 的小版本升级，只挡主版本}{
  \texttt{\^{}1.0.0} \emph{会}解析到 $1.1.1$}
\vrow{\texttt{\textasciitilde1.2}}{只允许补丁级}{$1.2.x$}
\vrow{\texttt{a|b}}{并集（OR）}{任一满足}
\vrow{\texttt{*}}{任意版本}{用于非关键的辅助模块}
\vnote{第二行是本文档反复强调的一个语义细节：\textbf{"钉住"}钉住的到底是什么。
      \texttt{\^{}1.0.0} 允许小版本升级，所以它\emph{不能}阻止数值行为的
      变化 —— 那必须靠主版本号表达。\crefchap{ch:sp1d} 记录了这一点。}
\end{veriftable}

\subsection{依赖解析}

解析做三件事：\emph{拓扑排序}、\emph{环检测}、\emph{版本冲突检测}。
三者都是"在加载阶段就把错误暴露出来"，而不是等到运行到一半才发现
某个协作者版本不对。

\subsection{参数校验：给自动化层设的闸门}
\label{sec:arch-params}

未知参数、越界值、类型错误一律拒绝，并且对疑似拼写错误给出提示。
本文档把这条机制称为\textbf{"LLM 生成配置的安全闸门"}：
在"生成 $\to$ 校验 $\to$ 带着错误信息重试"的闭环里，
一个\emph{确定性的}拒绝比一段自然语言的建议有用得多。

\subsection{热插拔}
\index{1concept@概念!hot plug@热插拔}

运行中可以注册新版本，而\textbf{已经创建的实例继续使用它被创建时的版本}。
因此"升级模块"是一个显式动作（改 Deck 绑定），而不是一次静默的全局替换。

\section{服务即插件}
\label{sec:arch-services}
\index{1concept@概念!service slot@服务槽}

\subsection{两类模块}

大多数模块是\emph{数据阶段}：读 artifact、写 artifact，由流水线排序。
另有一类是\emph{库模块} —— 非线性求解器、线性求解器、计算后端 ——
它们不转换任何上下文、不产出任何 artifact，但会被\emph{另一个模块内部}
需要。

后者通过上下文上的 \texttt{ServiceRegistry} 按需解析，而 Deck 仍然配置它们，
位置是专属的\emph{服务槽}：

\begin{equation}
\text{服务槽} \in \{\ \akey{backend},\ \akey{nonlinear},\ \akey{linear}\ \}.
\label{eq:arch-service-slots}
\end{equation}

\begin{center}
\small
\begin{tabular}{@{}p{0.16\textwidth} p{0.20\textwidth} p{0.56\textwidth}@{}}
\toprule
\textbf{槽} & \textbf{要求能力} & \textbf{本文档中的实现}\\
\midrule
\akey{nonlinear} & \texttt{nonlinear} & \pid{nonlinear.newton}（\crefchap{ch:newton}）\\
\akey{linear} & \texttt{linear} & \pid{linear.blocktridiag}（\crefchap{ch:blocktridiag}）\\
\akey{backend} & \texttt{backend} & 未实现（预留给计算后端）\\
\bottomrule
\end{tabular}
\end{center}

\subsection{解析语义}

\begin{itemize}[leftmargin=1.7em,itemsep=0.28em,topsep=0.35em]
  \item \textbf{惰性}：首次 \texttt{get()} 时才建实例；
  \item \textbf{按配置缓存}：缓存键为
        $(\text{模块标识} \mid \text{版本} \mid \text{规范化参数})$，
        因此同一模块的两组参数会得到两个独立实例 —— 否则后一次调用会
        静默继承前一次的配置；
  \item \textbf{调用方可覆盖}：物理模块传入的参数覆盖 Deck 槽里的同名参数。
        这正是物理模块注入\emph{自己量纲}的容差的方式
        （\crefchap{ch:sp1d} 传的是以 \U{eV} 计的容差）；
  \item \textbf{可软解析}：\texttt{required=False} 时缺失返回空值而不是抛错，
        供可选协作者使用；
  \item \textbf{可统一释放}：\texttt{release\_all()} 逐个终止，
        单个失败不阻断其余。
\end{itemize}

\subsection{为什么值得这么做}

\crefeq{eq:arch-service-slots} 的设计后果是：
\textbf{换迭代策略与换物理模型是同一种操作} —— 改 Deck 里的一行绑定。
一个未来的 Krylov 或信赖域求解器只要声明 \akey{nonlinear} 能力、
暴露同样的 \texttt{solve()} 就能顶上，物理模块一行不用改。

这不是设想：本文档记录了三处真实的替换验收，见 \S\ref{sec:arch-acceptance}。

\section{流水线：顺序由数据推导}
\label{sec:arch-pipeline}
\index{1concept@概念!pipeline@流水线}

执行顺序\textbf{不由槽位表决定，而由 manifest 推导}：
若模块 $\mathcal{A}$ 产出的 artifact 是模块 $\mathcal{B}$ 所需要的，
则 $\mathcal{A}$ 必在 $\mathcal{B}$ 之前。槽位顺序只做\emph{平局裁决}
（当两个模块之间没有数据依赖时，按槽位表的固定顺序排列）。

这条规则有一个可以被观察到的后果：把 \akey{materials} 槽从内置库换成
\pid{materials.dft}（\crefchap{ch:dft}）之后，流水线\emph{多出一个阶段}
并排在首位，而其余阶段的相对顺序\emph{一个都没动}。
这不是特例处理的结果，而是"顺序由数据推导"的必然。

\section{Deck：两级校验}
\label{sec:arch-deck}
\index{1concept@概念!deck@配置清单}

Deck 是描述一次仿真活动的声明式文档。它的校验分两级，
这一分法的动机与参数校验相同 —— 服务于自动化闭环：

\begin{enumerate}[leftmargin=1.9em,itemsep=0.35em,topsep=0.4em]
  \item \textbf{结构校验}（构造时抛 \texttt{SchemaError}）：
        类型、范围、必填、跨字段规则；未知字段一律拒绝并给拼写提示。
  \item \textbf{上下文校验}（返回 \texttt{(errors, warnings)} 而\emph{不抛异常}）：
        插件是否存在、版本是否可满足、参数是否为该插件已知、
        材料名是否可解析、\emph{artifact 是否可达}。
\end{enumerate}

第二级返回 issue 列表而不是抛异常，是刻意的：
调用方（尤其是一个自动化的生成器）需要拿到\emph{完整}的错误清单
一次性修正，而不是被迫玩"改一个、再撞下一个"的游戏。

\section{架构假设的验收}
\label{sec:arch-acceptance}

一套架构主张只有在被\emph{替换}检验过之后才算数。本文档记录了三处：

\begin{veriftable}{三处即插即用替换及其验收}{tab:arch-acceptance}
\vrow{\textbf{量子模型}}{解析三角阱 $\to$ 自洽 SP}{
  只改 Deck 一行；下游 \pid{extract.fom} 零改动；
  版本钉住生效（\crefchap{ch:sp1d}）}
\vrow{\textbf{材料库}}{内置库 $\to$ DFT 覆盖}{
  流水线只多一个阶段且排在首位，其余顺序不动；
  同一 Deck 两次运行指纹一致（\crefchap{ch:dft}）}
\vrow{\textbf{线性步}}{SciPy 稀疏（设计稿）$\to$ 纯 NumPy 块三对角}{
  经 \akey{linear} 槽解析；\pid{transport.dd} 的物理代码未因此改动
  （\crefchap{ch:blocktridiag}）}
\vnote{第一处的意义最强：它同时证明了能力契约、artifact 契约与版本钉住
      三件事。第三处则是"零第三方依赖"这一约束下的被迫替换，
      恰好反过来验证了服务槽的存在价值。}
\end{veriftable}

\section{指纹：可复现性的定义}
\label{sec:arch-fingerprint}

一次运行的\emph{指纹}是如下内容的稳定散列（前 16 位十六进制）：

\begin{equation}
\text{fp} \;=\; \mathrm{hash}_{16}\bigl(
\underbrace{\text{参数}}_{\text{Deck 快照}},\
\underbrace{\text{模块版本}}_{\text{解析结果}},\
\underbrace{\text{材料库指纹}}_{\text{内容散列}},\
\underbrace{\text{标记}}_{\text{调用方注解}}\bigr).
\label{eq:arch-fingerprint}
\end{equation}

\textbf{刻意的缺失：artifact 的数值不进指纹。}
这不是疏忽，而是一条重要的语义约定：指纹回答的是
"\emph{算了什么}"，而不是"\emph{算出了什么}"。
因此当数值行为发生变化时，指纹\emph{不会}自动改变 ——
它必须通过\textbf{升版本号}来表达。\crefchap{ch:sp1d} 记录了这条约定的
一次真实应用：Jacobian 精确化使 self-consistent 指纹从
\texttt{c26cb881f33c000c} 变为 \texttt{9d716641ad967853}，
而这只是因为模块版本从 1.1.0 升到了 1.1.1。

把这条约定写在总论里，是因为它容易被误用：若期待指纹能捕获数值漂移，
就会漏掉"忘了升版本"这一类错误。指纹保证的是\emph{同样的输入与同样的
版本给出同样的东西}，不是\emph{代码没变}。

\cleardoublepage
          % 第 2 章   系统架构
% 总论第 3 章 —— 数据流与 artifact 契约

\chapter{数据流与 artifact 契约}
\label{ch:dataflow}
\index{1concept@概念!artifact@数据通道}

模块之间\textbf{从不互相 import}，只通过命名数据通道交换数据。
本章记述这条通道的机制、它的全部内容（契约矩阵），以及三处
"接口取舍的代价"—— 后者是本章最有价值的部分，因为接口设计的好坏
只能在\emph{代价被量化}之后才能判断。

\section{命名数据通道}
\label{sec:data-context}

\subsection{三个阶段式原语}

一个模块从上下文读它需要的东西、写它产出的东西。写入时声明生产者，
读取时若键不存在则报错：

\begin{center}
\small
\begin{tabular}{@{}p{0.30\textwidth} p{0.60\textwidth}@{}}
\toprule
\textbf{原语} & \textbf{语义}\\
\midrule
\akey{require(key)} & 读一个必需的 artifact；缺失即抛异常\\
\akey{get(key)} & 读一个可选的 artifact；缺失返回空值\\
\akey{put(key, value, ...)} & 写 artifact，同时声明生产者、版本与元数据\\
\bottomrule
\end{tabular}
\end{center}

契约的机械性在于：\textbf{下游只认键，不认实现}。
这正是 \crefchap{ch:sp1d} 的替换只需改一行绑定的原因，
也是 \crefchap{ch:dft} 的"零改动替换"验收能够成立的原因。

\subsection{stage 作用域与溯源}

模块执行时被包在一个 \texttt{stage} 作用域里，声明它的输入与输出：

\begin{center}
\small
\begin{tabular}{@{}p{0.28\textwidth} p{0.62\textwidth}@{}}
\toprule
\textbf{记录项} & \textbf{内容}\\
\midrule
模块与版本 & \akey{module_id} 与 \akey{module_version}（即解析结果）\\
状态 & 成功 / 失败\\
耗时 & 该阶段的墙钟时间\\
输入输出 & 声明过的 artifact 键\\
\bottomrule
\end{tabular}
\end{center}

\noindent
\index{1concept@概念!provenance@溯源}
溯源不是装饰：\crefchap{ch:dft} 的 \akey{dft_provenance} 之所以能
作为\emph{独立的 artifact} 存在，前提是上下文本身已经能记录
"哪一步产出了什么"。

\section{artifact 契约矩阵}
\label{sec:data-matrix}

\creftab{tab:data-matrix} 汇总了本文档当前版本的全部 artifact 键，
以及它们的生产者与消费者。它是由九个模块的 \akey{produces}/\akey{requires}
声明生成的，因此与代码同步是机械的。

{\footnotesize
\setlength{\tabcolsep}{4pt}
\begin{longtable}{@{}p{0.225\textwidth} p{0.235\textwidth} p{0.245\textwidth} p{0.215\textwidth}@{}}
\caption{artifact 契约矩阵：每个键的生产者与消费者}\label{tab:data-matrix}\\
\toprule
\textbf{键} & \textbf{生产者} & \textbf{消费者} & \textbf{载体}\\
\midrule
\endfirsthead
\toprule
\textbf{键} & \textbf{生产者} & \textbf{消费者} & \textbf{载体}\\
\midrule
\endhead
\bottomrule
\endfoot
\akey{structure} & Deck（种入） &
  \pid{mesh.structured2d}, \pid{materials.dft},
  \pid{physics.polarization.ambacher}, \pid{transport.dd} & 字典\\
\akey{materials} & Deck 种入 \emph{或} \pid{materials.dft} &
  \pid{physics.polarization.ambacher}, \pid{transport.dd} & \texttt{MaterialDB}\\
\akey{dft_provenance} & \pid{materials.dft} & （审计，无消费者） & 字典\\
\akey{mesh} & \pid{mesh.structured2d} & \pid{transport.dd} & 字典\\
\akey{polarization} & \pid{physics.polarization.ambacher} &
  \pid{physics.quantum.analytic_tw}, \pid{physics.quantum.sp1d},
  \pid{transport.dd}, \pid{extract.fom} & 字典\\
\akey{quantum} & \pid{physics.quantum.analytic_tw} \emph{或} \pid{physics.quantum.sp1d} &
  \pid{extract.fom}, \pid{transport.dd}（可选初值） & 字典\\
\akey{ns} & \pid{physics.quantum.analytic_tw}, \pid{physics.quantum.sp1d},
  \pid{transport.dd} & \pid{extract.fom} & 标量\\
\akey{electrostatics} & \pid{transport.dd} & （未来：\pid{extract.iv}） & 字典\\
\akey{carriers} & \pid{transport.dd} & （未来：\pid{extract.iv}） & 字典\\
\akey{current} & \pid{transport.dd} & （未来：\pid{extract.iv}） & 字典\\
\akey{iv} & \pid{transport.dd} & （未来：\pid{extract.iv}） & 字典\\
\akey{results.fom} & \pid{extract.fom} & 终端输出 & 字典\\
\akey{gate}, \akey{temperature_K} & Deck & 多个模块 & 字典 / 标量\\
\end{longtable}
}

\noindent
读这张表有两种读法。\textbf{横向}看依赖：\akey{polarization} 是扇出最广的键
（四个消费者），因此它的契约稳定性最重要。\textbf{纵向}看层级：
从 Deck 种入的 \akey{structure}/\akey{materials} 出发，
经 \akey{polarization} 分成"量子支"与"输运支"，
最后由 \pid{extract.fom} 收敛到单一终端键。

\section{三处接口取舍及其代价}
\label{sec:data-tradeoffs}

契约设计中最有信息量的部分不是"我们提供了什么"，而是
"\emph{我们明知有代价、仍然这么选}"。本节记录三处，
其中第一处的代价已被量化。

\subsection{界面落位：两个模块的不一致是被接受的}

\crefchap{ch:mesh} 产出的网格把层界放在\textbf{节点}上；
\crefchap{ch:sp1d} 的量子求解器刻意构造网格把异质界面放在\textbf{面}上。
两者对同一条物理界面给出了不同的离散表示。

这不是疏忽，而是接口边界的取舍：网格模块产出的是\emph{几何}，
供 \crefchap{ch:transport_dd} 的有限体积漂移扩散使用，后者按设计
\emph{原样消费}给定网格，因此界面只能在节点上。

代价是量化的：\crefchap{ch:transport_dd} 实测，由该网格解出的
二维电子气面密度比极化给出的上界 $\sigeff/q$ 高约 $1.7$ 倍
（界面附近一个单元的抹平）。同时，设计稿中"与一维 SP 的带边差
$<1$\U{meV}"这条判据在本离散下\emph{原理上不可达}。

\boundary{本条的教训是：\textbf{跨模块的接口取舍必须把代价写进文档，
而不是留给读者去发现两章之间的叙述矛盾}。本文档在
\crefchap{ch:mesh}（声明落位方式）、\crefchap{ch:sp1d}（声明相反方式）、
\crefchap{ch:transport_dd}（量化代价）三处交叉引用，形成一条闭环线索。
}

\subsection{\akey{ns}：一个刻意冗余的标量键}

\akey{ns} 与 \akey{quantum.ns_cm2} 携带同一个数。看似冗余，
但它承担一个明确职责：让\emph{只需要面密度}的消费者不必依赖
量子模型的完整契约。\crefchap{ch:transport_dd} 也产出 \akey{ns}
（最后偏压点的值），因此 \pid{extract.fom} 既可以在平衡态流水线上
工作，也可以在有偏压的流水线上工作，而它对上游的\emph{唯一}要求是
"有人产出 \akey{ns}"。

代价是可能出现两个不一致的 \akey{ns}。本文档的应对是：
\pid{extract.fom} 的校验里有一条断言，要求 \akey{ns} 与
\akey{quantum.ns_cm2} 在 $10^{-6}$ 相对误差内一致
（\crefchap{ch:fom}）。\textbf{冗余可以接受，冗余的静默漂移不可以。}

\subsection{场数组只记最后一个偏压点}

\crefchap{ch:transport_dd} 的 \akey{electrostatics} 与 \akey{carriers}
只携带\emph{最后}求解的偏压点的二维场，而\emph{每一个}偏压点的摘要
都在 \akey{current} 里。理由是容量：二维场数组随偏压点数线性增长，
全存会撑爆上下文；而"哪一点没收敛、为什么"这个诊断需求，
用标量摘要就能满足。

代价是：若使用者想比较两个偏压点的场分布，必须在 Deck 里把扫描
拆成两次运行。这是一个被接受的使用限制，写在 \crefchap{ch:transport_dd}
的接口一节。

\section{契约演化的纪律}
\label{sec:data-evolution}
\index{1concept@概念!contract evolution@契约演化}

由上述取舍可以总结出三条本文档遵循的契约纪律，
它们是后续新增模块时必须遵守的：

\begin{enumerate}[leftmargin=1.9em,itemsep=0.4em,topsep=0.4em]
  \item \textbf{契约键的语义只增不改。} 新增键是兼容的；
        改变一个已有键的\emph{单位或含义}必须升主版本，
        因为下游无法从值本身察觉这种变化。
  \item \textbf{冗余键必须有一致性断言。} 见 \S\ref{sec:data-tradeoffs}
        的 \akey{ns} 一例：允许冗余，但必须有测试锁住两者一致。
  \item \textbf{审计信息另开 artifact，不混进数据契约。}
        \crefchap{ch:dft} 的 \akey{dft_provenance} 是范例：
        \akey{materials} 的键空间对下游是不区分来源的，
        因此溯源必须\emph{另开}一个通道，否则下游模块会被迫理解
        自己本不该关心的东西。
\end{enumerate}

第 3 条还可以反向表述，作为判断一个契约设计是否健康的标准：
\textbf{如果一个下游模块需要知道"上游是谁"，那么契约设计就已经失败了。}

\cleardoublepage
              % 第 3 章   数据流与 artifact 契约
% 总论第 4 章 —— 验证方法学

\chapter{验证方法学}
\label{ch:verification-methodology}
\index{1concept@概念!verification@验证}

本章记述"凭什么相信这些数字"这件事本身。它不描述任何一个模块，
而是描述九个模块章里那些验证表\emph{共同遵循的判据体系}。
把它单列一章的理由很直接：本文档全篇出现的几百个数字，
其可信度完全取决于本章所定义的判据是否恰当。

\section{四层验证框架}
\label{sec:verify-framework}

方案把验证分为四层，本文档沿用其分类：

\begin{center}
\small
\begin{tabular}{@{}p{0.24\textwidth} p{0.34\textwidth} p{0.34\textwidth}@{}}
\toprule
\textbf{层} & \textbf{内容} & \textbf{本文档中的体现}\\
\midrule
代码验证 & 收敛阶、闭式解、制造解 & 各章 §x.5 的主体\\
解验证 & 网格收敛、容差与步长扫描 & 第 \ref{ch:sp1d}、\ref{ch:transport_dd} 章\\
确认（Validation） & 与实测/独立实现对比 & \textbf{仅部分完成}，见下\\
不确定度量化 & 敏感性、蒙特卡洛 & 未实现\\
\bottomrule
\end{tabular}
\end{center}

第三层是本文档最弱的一环，必须明说：\textbf{基准阶梯 L0--L7 目前只走到 L1}。
L2（Hall 对标面密度）及其后的每一级都尚未进行，因此
"绝对数值不可信"这句话在本文档中反复出现，其根据就在这里。

\begin{veriftable}{基准阶梯的当前进度}{tab:verify-ladder}
\vrow{L0\ \ 平衡态、一维 Poisson、极化解析式}{$<10^{-6}$}{\textbf{达成}（第 \ref{ch:sp1d}、\ref{ch:polarization} 章）}
\vrow{L1\ \ 一维 Schrödinger--Poisson 子带}{子带能量差 $<5$\U{meV}}{部分达成（无独立实现对照）}
\vrow{L2\ \ 标准异质结 $\ns$、C--V}{$\ns$ 差 $<5\%$}{\textbf{未做}（需霍尔测量）}
\vrow{L3\ \ 与商用 TCAD 交叉对比 IV/C--V}{$<5\%$}{未做}
\vrow{L4\ \ 直流输出/转移特性}{$I_{d\max}/g_m<10\%$，$V_{\text{th}}<0.15$\U{V}}{未做}
\vrow{L5\ \ S 参数 $\to f_T/f_{\max}$}{$<10$--$15\%$}{未实现（无 S 参数链）}
\vrow{L6\ \ pulsed IV、gate/drain-lag}{坍塌幅度 $<20\%$}{未实现（无陷阱模型）}
\vrow{L7\ \ load-pull}{$P_{\text{out}}$/PAE $<10$--$15\%$}{未实现}
\end{veriftable}

\section{五类判据}
\label{sec:verify-criteria}

本文档 §x.5 的验证表使用且只使用五类判据。每一类都有它的适用场景
与它\emph{不能}证明的东西 —— 知道后者比知道前者更重要。

\subsection{第一类：闭式解}
\label{sec:verify-closed-form}
\index{3method@方法!closed-form benchmark@闭式解基准}

\textbf{能证明}：离散格式在这个特定问题上的精度。
\textbf{不能证明}：格式在别的问题上也是对的。

闭式解是最强的单项证据，因为它是\emph{精确}的，比较不需要容差解释。
本文档中的典型值：

\begin{veriftable}{闭式解基准的代表性结果}{tab:verify-closed}
\vrow{双介质平行板（含界面面电荷）}{精确}{\textbf{$<10^{-10}$\U{V}}（\crefchap{ch:sp1d}）}
\vrow{无限深方阱 $E_n=n^2\pi^2\hb^2/(2\meff L^2)$}{二阶收敛}{二阶收敛（\crefchap{ch:sp1d}）}
\vrow{均匀电阻 $J=qN_D\mumu V/L$}{$<0.1\%$}{\textbf{$<10^{-6}$}（\crefchap{ch:transport_dd}）}
\vrow{块三对角解对稠密解}{$<10^{-10}$}{通过（四组尺寸，\crefchap{ch:blocktridiag}）}
\vrow{Bernoulli 函数 $B(x)=x/(\ee^x-1)$}{$10^{-12}$ 相对}{通过（\crefchap{ch:transport_dd}）}
\end{veriftable}

注意最后两行的性质差别：第一至三行检验的是\emph{离散格式}，
第四行检验的是\emph{求解器}，第五行检验的是\emph{一个特殊函数}。
把它们分开陈述，是因为它们失效时的含义完全不同。

\subsection{第二类：收敛阶}
\label{sec:verify-order}
\index{3method@方法!convergence order@收敛阶}

\textbf{能证明}：格式的\emph{阶}是对的 —— 这比"某个网格上误差小"
强得多，因为它意味着加密时误差会按预期速率下降。
\textbf{不能证明}：误差常数小。一个一阶格式可能在粗网格上比
二阶格式更准。

网格收敛阶由 Richardson 三点公式估计：

\begin{equation}
p \;=\;
\frac{\ln\!\bigl(\abs{u_1-u_2}\,/\,\abs{u_2-u_3}\bigr)}
     {\ln\!\bigl(h_1/h_2\bigr)},
\qquad h_1>h_2>h_3 .
\label{eq:verify-richardson}
\end{equation}

迭代法的\emph{局部收敛阶}由相邻三个残差估计：

\begin{equation}
p \;=\;
\frac{\ln\!\bigl(\norm{r_{k+1}}/\norm{r_k}\bigr)}
     {\ln\!\bigl(\norm{r_k}/\norm{r_{k-1}}\bigr)} .
\label{eq:verify-local-order}
\end{equation}

\crefeq{eq:verify-local-order} 的使用有一个必须说明的细节：
\textbf{只有落在渐近窗口内的三元组才有意义}。实现上取
"残差严格下降且比值本身在缩小"的三元组，并排除最小成员已落到
舍入地板（$\sim10^{-15}$）的样本 —— 地板以下的比值是噪声。
本文档报告的两个关键值都取自这样的窗口：耦合 SP 的
$p\approx1.99$（\crefchap{ch:sp1d}）与制造解的
$2.008/2.002/2.001$（\crefchap{ch:transport_dd}）。

\subsection{第三类：制造解}
\label{sec:verify-mms}
\index{3method@方法!method of manufactured solutions@制造解方法}

\textbf{能证明}：算子在\emph{任意}已知解上的精度，而不只是在恰好有
简单闭式解的那几个问题上。
\textbf{不能证明}：边界条件的处理是否正确（制造解通常搭配周期性或
简单边界）。

做法是：取一个解析函数 $\psi^{\star}$，把它代入微分算子算出源项，
再把该源项作为已知右端求解，比较数值解与 $\psi^{\star}$。
本文档中的范例在 \crefchap{ch:transport_dd}：取
$\psi^{\star}=\sin(\pi z/L)$（与横向无关，从而\emph{隔离}竖直算子），
加密四档得到的观测阶为 $2.008/2.002/2.001$。
标准表述见 \cite{roache2002,salari2000}。

制造解的一个常被忽略的好处是：它使\textbf{"界面以外的格式是二阶的"}
成为一个可断言的命题，即便界面本身只有一阶 —— 只要制造解在界面两侧
都是光滑的。\crefchap{ch:transport_dd} 正是这样把两件事分开陈述的。

\subsection{第四类：不变量与趋势}
\label{sec:verify-invariants}
\index{1concept@概念!invariant assertion@不变量断言}

\textbf{能证明}：某个\emph{必须成立的关系}在实现中确实成立。
\textbf{不能证明}：任何数值精度。

不变量断言的价值在于\emph{耐久}：它断言的是"若把 Al 组分调高，
极化必然增强"这类关系，而不是某个具体数字。
因此它在参数演化后依然有效 —— 这对一个参数尚未标定的项目
尤其重要。本文档中此类断言的例子：

\begin{itemize}[leftmargin=1.7em,itemsep=0.28em,topsep=0.35em]
  \item 平衡态通量恒为零（\crefchap{ch:transport_dd}，实测 $2.7\times10^{-15}$\U{A/m}）；
  \item 势偏导的反对称性 $\partial J/\partial\psi_L=-\partial J/\partial\psi_R$
        （规范不变性的离散形式）；
  \item 密度偏导与密度无关（跨 $17$ 个数量级验证）；
  \item 面密度低于极化上界 $\sigeff/q$；
  \item 屏蔽因子严格按比例缩放被测到 $10^{-6}$ 位。
\end{itemize}

\subsection{第五类：负向基准}
\label{sec:verify-negative}
\index{1concept@概念!negative benchmark@负向基准}

\textbf{能证明}：某个机制是\emph{必要}的。
\textbf{不能证明}：它是最优的。

负向基准记录的是\emph{另一个配置下的失败}。本文档中有两处：

\begin{veriftable}{负向基准（记录失败，以证明机制的必要性）}{tab:verify-negative}
\vrow{关闭线搜索后从远起点出发}{不收敛}{
  通过（$x^3-1$ 自 $x=40$、上限 $4$ 步；\crefchap{ch:newton}）}
\vrow{遗漏态密度质量响应时的误差地板}{不随差分步长下降}{
  $1.3\times10^{-5}$，是"缺项"的特征（\crefchap{ch:sp1d}）}
\vnote{第二行尤其值得注意：它记录的是一个\emph{已修正的}错误实现的症状。
      保留它，是因为"如何识别一个缺项"这件事无法用正面例子说清。}
\end{veriftable}

\section{导数校验：一个必须写进方法学的教训}
\label{sec:verify-jacobian}
\index{3method@方法!Jacobian verification@Jacobian 校验}

Newton 类方法的正确性几乎完全取决于 Jacobian。本文档记录了一条
\emph{代价高昂}的教训，它的普遍性超出了本项目。

\subsection{校验公式}

解析 Jacobian 与中心差分比较：

\begin{equation}
\mat{J}^{\text{FD}}_{:,j} \;=\;
\frac{\mat{R}(\bm{x}+h\bm{e}_j)-\mat{R}(\bm{x}-h\bm{e}_j)}{2h},
\qquad h = 10^{-6}\max\!\bigl(1,\norm{\bm{x}}_\infty\bigr).
\label{eq:verify-fd}
\end{equation}

分歧在于\textbf{如何把 $\mat{J}^{\text{FD}}_{:,j}-\mat{J}_{:,j}$ 归一化}。
两种做法：

\begin{equation}
e_j^{\text{col}} \;=\;
\frac{\bigl\lVert \mat{J}^{\text{FD}}_{:,j}-\mat{J}_{:,j}\bigr\rVert_\infty}
{\max\bigl(\bigl\lVert \mat{J}^{\text{FD}}_{:,j}\bigr\rVert_\infty,\
       \bigl\lVert \mat{J}_{:,j}\bigr\rVert_\infty\bigr)},
\qquad
e_i^{\text{row}} \;=\;
\frac{\bigl\lVert \mat{J}^{\text{FD}}_{i,:}-\mat{J}_{i,:}\bigr\rVert_\infty}
{\max\bigl(\bigl\lVert \mat{J}^{\text{FD}}_{i,:}\bigr\rVert_\infty,\
       \bigl\lVert \mat{J}_{i,:}\bigr\rVert_\infty\bigr)} .
\label{eq:verify-fd-norms}
\end{equation}

\subsection{为什么逐列不够}

\crefchap{ch:transport_dd} 的未知量从 $(\psi,\varphi_n)$ 改为 $(\psi,n)$ 时，
Jacobian 里残留了一行旧代码，把第二个未知量当作准费米势代进了态密度公式。
后果是\emph{内部密度与残差相差 $17$ 个数量级}，
连续性行因此完全失去对 $\psi$ 的耦合，整个耦合矩阵在 Newton 最需要的
方向上近奇异。

而当时的\emph{逐列}归一化校验报告相对误差只有 $10^{-10}$ ——
因为坏掉的那一项在整列中绝对值很小，被子列中量级正常的其他项淹没了。
改成\emph{逐行}相对判据后，它立刻暴露。

\boundary{由此得到两条纪律，它们适用于任何耦合系统的实现：

\textbf{（i）导数校验用逐行相对判据，不用逐列。}
行对应"一个方程"，列对应"一个未知量"；一个方程整体错位，
在行视角下是 $100\%$ 的错误，在列视角下可能微不足道。

\textbf{（ii）校验要在单元级与沿路径做两遍。}
"通量本身对中心差分"（单元级）与"整个 Jacobian 沿 Newton 路径
对中心差分"（系统级）是两道独立的闸门。\crefchap{ch:transport_dd}
两道都做了：前者逐项验四个偏导，后者沿路径验到 $\sim10^{-10}$。

\textbf{（iii）一致性不等于正确性。}
\crefchap{ch:transport_dd} 另记录了一例：Bernoulli 函数的反射恒等式
写反之后，通量仍然\emph{完全自洽} —— 系统级 Jacobian 对中心差分
照样通过 $4\times10^{-10}$ —— 却在每个带边沿电流方向上升的面上
系统性错误。修好后平衡态残差从 $35$ 降到 $4.8\times10^{-6}$（七个数量级）。
\textbf{自洽性只能排除"不自洽"，不能排除"一致地错"。}
}

\subsection{本文档中两个校验器的差异}

一个值得记录的细节：\crefchap{ch:newton} 的通用校验器用的是
\emph{逐列}归一化（\crefeq{eq:verify-fd-norms} 左式），
而 \crefchap{ch:transport_dd} 用自己的逐行判据。
这不是不一致，而是适用范围不同：前者的服务对象是良缩放的稠密系统，
后者是行均衡过的耦合系统。本文档在两处都注明了这一差异，
把它作为\emph{适用边界}而不是藏起来。

\section{可复现性作为一类判据}
\label{sec:verify-reproducibility}
\index{1concept@概念!reproducibility@可复现性}

有三类断言检验的不是精度，而是"同样的输入是否给出同样的东西"。
它们看似平凡，却是上述一切验证的前提。

\begin{veriftable}{三类可复现性断言}{tab:verify-repro}
\vrow{同输入同指纹}{两次运行的指纹相等}{
  通过（\crefchap{ch:sp1d} 的 $2\times$ 运行、\crefchap{ch:dft} 的同 Deck 两次运行）}
\vrow{换实现则指纹不同}{指纹不等}{
  通过（解析量子 Deck 与自洽量子 Deck、\crefchap{ch:dft} 的材料槽替换）}
\vrow{初值/延拓无关性}{解与初值路径无关}{
  \crefchap{ch:sp1d}：热启动 $0.35\to0.95$ 上界，相对差 $<0.5\%$；
  \crefchap{ch:transport_dd}：两条延拓阶梯，相对差 $<5\times10^{-3}$}
\vnote{第三行是本文档中最被低估的一类断言。它证明了"初值与延拓只影响
      收敛路径，不影响收敛解"—— 而这条性质一旦不成立，
      Newton 类方法的全部结果都会变得依赖调参历史，不再是方程的解。}
\end{veriftable}

\section{文档侧的表述纪律}
\label{sec:verify-doc-discipline}

最后，把上述判据体系翻译成一条写作纪律。本文档的每一张验证表
都有三列，且三列的含义是固定的：

\begin{center}
\small
\begin{tabular}{@{}p{0.16\textwidth} p{0.72\textwidth}@{}}
\toprule
\textbf{列} & \textbf{含义与纪律}\\
\midrule
基准 & 检验的对象。措辞要能回答"若它错了，什么会跟着错"\\
判据 & \textbf{测试断言的门槛}。它是回归门禁，不随版本漂移\\
实测 & \textbf{编写时观测到的值}。它带有效数字、可回查，但会随版本演化；
       引用时必须连同模块版本一并引用\\
\bottomrule
\end{tabular}
\end{center}

\noindent
纪律的另一半是\textbf{什么不能进"实测"列}：仅有趋势正确、
仅有一阶吻合、或无独立基准可对照者，一律进 §x.6 的"适用边界"。
本文档第 \ref{ch:transport_dd} 章 T1 基准中"未做"的一项
（与一维 SP 的带边差）就是这条纪律的应用：它在\emph{原理上}
不可达，因此被写成诚实结论，而不是待办事项。

\cleardoublepage
          % 第 4 章   验证方法学
% 总论第 5 章 —— 模块索引与依赖关系

\chapter{模块索引与依赖关系}
\label{ch:module-index}

本章是全文的索引页：把九个模块的契约、依赖、能力与版本集中在一处，
使读者不必翻遍九章才能拼出整体结构。

\section{九个模块一览}
\label{sec:index-overview}

\creftab{tab:index-modules} 是本文档当前版本的完整模块清单。
它与代码中的 \texttt{MANIFEST} 逐字段一致；任何一行变化
都意味着某个模块章需要同步更新。

{\footnotesize
\setlength{\tabcolsep}{4pt}
\begin{longtable}{@{}p{0.05\textwidth} p{0.245\textwidth} p{0.075\textwidth}
                    p{0.165\textwidth} p{0.165\textwidth} p{0.135\textwidth}@{}}
\caption{九个模块的完整清单（与代码 MANIFEST 逐字段一致）}\label{tab:index-modules}\\
\toprule
\textbf{章} & \textbf{模块标识} & \textbf{版本} &
\textbf{依赖 artifact} & \textbf{产出 artifact} & \textbf{验证状态}\\
\midrule
\endfirsthead
\toprule
\textbf{章} & \textbf{模块标识} & \textbf{版本} &
\textbf{依赖 artifact} & \textbf{产出 artifact} & \textbf{验证状态}\\
\midrule
\endhead
\bottomrule
\endfoot
6 & \pid{mesh.structured2d} & 0.1.0 &
  \akey{structure} & \akey{mesh} & partially\\
7 & \pid{materials.dft} & 0.1.0 &
  \akey{structure} & \akey{materials}, \akey{dft_provenance} & partially\\
8 & \pid{physics.polarization.ambacher} & 0.1.0 &
  \akey{structure}, \akey{materials} & \akey{polarization} & partially\\
9 & \pid{physics.quantum.analytic_tw} & 0.1.0 &
  \akey{polarization} & \akey{quantum}, \akey{ns} & partially\\
10 & \pid{physics.quantum.sp1d} & \textbf{1.1.1} &
  \akey{polarization}, \akey{materials} & \akey{quantum}, \akey{ns} & partially\\
11 & \pid{transport.dd} & 0.1.0 &
  \akey{mesh}, \akey{structure}, \akey{materials}, \akey{polarization} &
  \akey{electrostatics}, \akey{carriers}, \akey{current}, \akey{iv} & partially\\
12 & \pid{extract.fom} & 0.1.0 &
  \akey{ns}, \akey{quantum}, \akey{polarization} & \akey{results.fom} &
  \textbf{unverified}\\
13 & \pid{nonlinear.newton} & 1.0.0 &
  （无） & （无） & partially\\
14 & \pid{linear.blocktridiag} & 0.1.0 &
  （无） & （无） & partially\\
\end{longtable}
}

\noindent
表中 "partially" 为 \texttt{partially-verified} 的缩写。
九个模块中只有 \pid{extract.fom} 是 \texttt{unverified}，
且是\emph{按设计如此}：它本身就是一个模型的闭式解，
没有闭式解可以对，只能等 L2 基准（见 \crefchap{ch:fom}）。

\section{依赖图}
\label{sec:index-dag}

依赖关系是一个有向无环图，其分层结构见 \creftab{tab:index-dag}。
分层的依据是 \texttt{requires}/\texttt{produces}：
第 $\ell$ 层的模块只消费第 $<\ell$ 层的产出。

\begin{veriftable}{依赖图的分层结构}{tab:index-dag}
\vgroup{第 0 层\ \ 外部种入（由 Deck 提供，不由任何模块产出）}
\vrow{}{}{\akey{structure}, \akey{materials}（内置库）, \akey{gate}, \akey{temperature_K}}
\vgroup{第 1 层\ \ 几何与材料}
\vrow{\crefchap{ch:mesh}}{\akey{structure}}{$\to$ \akey{mesh}}
\vrow{\crefchap{ch:dft}}{\akey{structure}}{$\to$ \akey{materials}, \akey{dft_provenance}}
\vgroup{第 2 层\ \ 极化}
\vrow{\crefchap{ch:polarization}}{\akey{structure}, \akey{materials}}{$\to$ \akey{polarization}}
\vgroup{第 3 层\ \ 量子（二选一）}
\vrow{\crefchap{ch:analytic_tw}}{\akey{polarization}}{$\to$ \akey{quantum}, \akey{ns}}
\vrow{\crefchap{ch:sp1d}}{\akey{polarization}, \akey{materials}}{$\to$ \akey{quantum}, \akey{ns}}
\vgroup{第 4 层\ \ 输运}
\vrow{\crefchap{ch:transport_dd}}{\akey{mesh}, \akey{structure}, \akey{materials},
  \akey{polarization}（+\akey{quantum} 可选）}{
  $\to$ \akey{electrostatics}, \akey{carriers}, \akey{current}, \akey{iv}, \akey{ns}}
\vgroup{第 5 层\ \ 提取}
\vrow{\crefchap{ch:fom}}{\akey{ns}, \akey{quantum}, \akey{polarization}}{$\to$ \akey{results.fom}}
\vnote{第 3 层的两个模块\emph{互斥}：它们声明同一能力 \texttt{quantum}、
      产出同一组键，Deck 的 \akey{quantum} 槽只能绑其中一个。
      第 4 层对 \akey{quantum} 的依赖是\emph{可选}的
      （\akey{quantum_warm_start="auto"} 时存在即用，不存在则退回
      Poisson 热启动），因此输运层不因量子槽空着而无法运行。}
\end{veriftable}

\subsection{扇出与扇入}

\begin{itemize}[leftmargin=1.7em,itemsep=0.28em,topsep=0.35em]
  \item \textbf{扇出最广的键是 \akey{polarization}}（四个消费者）。
        它的契约稳定性最重要 —— 改它的键名会同时影响四章。
  \item \textbf{扇入最多的是 \pid{transport.dd}}（四个必需依赖加一个可选）。
        它是全流水线中装配最复杂的一步，也是唯一一个需要\emph{协作者}
        （\crefchap{ch:blocktridiag}）的物理模块。
  \item \textbf{两个服务模块与主图无关}：它们没有依赖也没有产出，
        只被调用。因此它们在 \creftab{tab:index-dag} 中不占任何一层。
\end{itemize}

\section{服务槽与能力矩阵}
\label{sec:index-slots}

\subsection{能力标签}

Deck 的槽位绑定由\emph{能力}校验，而不是由模块名硬编码。
\creftab{tab:index-capabilities} 列出全部能力标签及其声明者。

{\small
\begin{longtable}{@{}p{0.20\textwidth} p{0.34\textwidth} p{0.36\textwidth}@{}}
\caption{Deck 槽位、要求能力与声明该能力的模块}\label{tab:index-capabilities}\\
\toprule
\textbf{槽位} & \textbf{要求能力} & \textbf{声明该能力的模块}\\
\midrule
\endfirsthead
\toprule
\textbf{槽位} & \textbf{要求能力} & \textbf{声明该能力的模块}\\
\midrule
\endhead
\bottomrule
\endfoot
\akey{mesh} & \texttt{mesh} & \pid{mesh.structured2d}\\
\akey{materials} & \texttt{materials} & \pid{materials.dft}（另声明 \texttt{materials.dft}、\texttt{materials.validation}）\\
\akey{polarization} & \texttt{polarization} & \pid{physics.polarization.ambacher}\\
\akey{quantum} & \texttt{quantum} & \pid{physics.quantum.analytic_tw}、
  \pid{physics.quantum.sp1d}\\
\akey{transport} & \texttt{transport} & \pid{transport.dd}（另声明 \texttt{transport.dd}、
  \texttt{current}、\texttt{iv}）\\
\akey{mobility} & \texttt{mobility} & \textbf{无}（预留槽）\\
\akey{traps} & \texttt{traps} & \textbf{无}（预留槽）\\
\akey{thermal} & \texttt{thermal} & \textbf{无}（预留槽）\\
\midrule
\akey{nonlinear}\ \ (服务) & \texttt{nonlinear} & \pid{nonlinear.newton}\\
\akey{linear}\ \ (服务) & \texttt{linear} & \pid{linear.blocktridiag}\\
\akey{backend}\ \ (服务) & \texttt{backend} & \textbf{无}（预留槽）\\
\end{longtable}
}

\noindent
三个预留槽（\akey{mobility}、\akey{traps}、\akey{thermal}）的存在
说明了架构的意图：它们对应方案中尚未实现的物理
（场依赖迁移率的独立插件、陷阱、自热），而 Deck 的槽位表已经为它们留好了位置。
\crefchap{ch:transport_dd} 目前把迁移率作为\emph{自身参数}而非独立模块，
并明确记录了这一临时状态。

\subsection{提取器链}

\akey{extractors} 是一个\emph{列表}槽而不是单槽，因此可以串联多个提取器。
当前只有 \pid{extract.fom} 一个实现。流水线为它们分配的槽名形如
\texttt{extract:0}、\texttt{extract:1}，这一点在
\crefchap{ch:fom} 的验收中可见。

\section{版本与指纹}
\label{sec:index-versions}

\begin{veriftable}{当前解析到的模块版本}{tab:index-versions}
\vrow{九个模块的版本}{见 \creftab{tab:index-modules}}{
  最高版本为 \pid{physics.quantum.sp1d} 的 1.1.1}
\vrow{多版本共存}{同一标识可注册多个版本}{
  已验证：\pid{physics.quantum.sp1d} 的 1.1.1 与 2.0.0 并存，
  旧实例继续用旧版本}
\vrow{版本钉住语义}{\texttt{\^{}1.0.0} 允许 $1.x$ 升级}{
  因此它\emph{不}保证数值不变；那必须靠主版本号表达}
\vrow{指纹语义}{参数 + 模块版本 + 材料库 + 标记}{
  \textbf{artifact 数值不进指纹}，见 \crefchap{ch:architecture}}
\vnote{本文档记录的一次真实指纹变化：\pid{physics.quantum.sp1d} 的
      Jacobian 精确化使 self-consistent 指纹从 \texttt{c26cb881f33c000c}
      变为 \texttt{9d716641ad967853}，原因是模块版本从 1.1.0 升到 1.1.1；
      解析模型未受影响，指纹不变。}
\end{veriftable}

\section{章、插件与测试对照}
\label{sec:index-crossref}

\begin{center}
\small
\begin{tabular}{@{}p{0.055\textwidth} p{0.30\textwidth} p{0.30\textwidth} p{0.24\textwidth}@{}}
\toprule
\textbf{章} & \textbf{插件目录} & \textbf{测试文件} & \textbf{设计稿}\\
\midrule
6 & \fpath{huntian/plugins/mesh_structured2d/} & \fpath{tests/test_plugins.py} & 方案 §5.1\\
7 & \fpath{huntian/plugins/materials_dft/} & \fpath{tests/test_materials_dft*.py} &
  DFT 设计稿\\
8 & \fpath{huntian/plugins/polarization_ambacher/} & \fpath{tests/test_plugins.py} & 方案 §5.1\\
9 & \fpath{huntian/plugins/quantum_analytic_tw/} & \fpath{tests/test_plugins.py} & 方案 §5.1\\
10 & \fpath{huntian/plugins/quantum_sp1d/} & \fpath{tests/test_quantum_sp1d.py} & 方案 §5.2\\
11 & \fpath{huntian/plugins/transport_dd/} & \fpath{tests/test_transport_dd.py} &
  DD 设计稿\\
12 & \fpath{huntian/plugins/extract_fom/} & \fpath{tests/test_plugins.py} & 方案 §5.4\\
13 & \fpath{huntian/plugins/nonlinear_newton/} & \fpath{tests/test_nonlinear_newton.py} & 方案 §6\\
14 & \fpath{huntian/plugins/linear_blocktridiag/} & \fpath{tests/test_linear_blocktridiag.py} &
  DD 设计稿 §9\\
\midrule
— & — & \fpath{tests/test_hotswap_*.py} & 即插即用验收\\
— & — & \fpath{tests/test_services.py}、\fpath{test_pipeline.py}、
  \fpath{test_registry.py}、\fpath{test_deck.py}、\fpath{test_context.py} &
  骨架层验收\\
\bottomrule
\end{tabular}
\end{center}

\noindent
最后一组文件不属于任何单个模块章，它们验收的是\crefchap{ch:architecture}
与 \crefchap{ch:dataflow} 所描述的骨架本身：服务槽的装配、
流水线顺序的推导、注册中心的版本与依赖语义、Deck 的两级校验、
上下文的溯源与指纹。本文档当前版本对应的测试规模为
\textbf{426 个用例，全部通过}。

\section{阅读本文档之外的入口}

本文档是\emph{原理与证据}的汇总，不是使用手册。与本仓库的其他材料的分工是：

\begin{center}
\small
\begin{tabular}{@{}p{0.28\textwidth} p{0.62\textwidth}@{}}
\toprule
\textbf{材料} & \textbf{回答的问题}\\
\midrule
本文档 & 每个模块\emph{在算什么}、\emph{凭什么信}、\emph{边界在哪}\\
\fpath{README.md} & 怎么跑、跑出什么样、当前进度\\
\fpath{工作日志.md} & 每一步\emph{为什么}这么做、踩过什么坑、欠什么账\\
三份设计稿 & 实现\emph{之前}的意图，以及实现后的偏差记录\\
\bottomrule
\end{tabular}
\end{center}

\cleardoublepage
          % 第 5 章   模块索引与依赖关系

% ---- 第二部分 -------------------------------------------------------
\part{网格与材料}
% module_id: mesh.structured2d
% module_version: 0.1.0
% plugin_dir: huntian/plugins/mesh_structured2d
% source_files: huntian/plugins/mesh_structured2d/__init__.py
% tests: tests/test_plugins.py
%
% 模块章模板 —— 见 modules/quantum_sp1d.tex 文件头的骨架说明。
% x.3 与 x.4 两槽对非数值型模块使用模块专属标题，但两槽都必须在。

\chapter{结构化二维网格生成器}
\label{ch:mesh}

\begin{quote}
\small
\noindent\textbf{模块标识}\quad\pid{mesh.structured2d}\qquad
\textbf{版本}\quad 0.1.0\qquad
\textbf{类别}\quad mesh
\end{quote}

\noindent
本章记述器件几何到计算网格的映射。这是整条流水线的第一个阶段：它把外延层表
与版图参数转成一张\textbf{张量积}（结构化）网格，供后续的极化、量子与输运
模块共享。此处的决策看似只是"划分网格"，但它决定了后续所有算子能长成什么
形状 —— \crefchap{ch:transport_dd} 的 Jacobian 之所以是块三对角、可以用
$O(n_zm^{3})$ 的直接法求解，前提正是本章产出的结构化网格。

\section{模块契约}
\label{sec:mesh-contract}
\index{4module@模块!mesh.structured2d}

\begin{modulecard}{mesh.structured2d}{0.1.0}
  \cardrow{类别}{mesh（网格生成，流水线阶段）}
  \cardrow{能力}{mesh, mesh.structured2d}
  \cardrow{依赖 artifact}{structure}
  \cardrow{产出 artifact}{mesh}
  \cardrow{验证状态}{partially-verified}
  \cardrow{源码}{\fpath{huntian/plugins/mesh_structured2d/}}
  \cardrow{测试}{\fpath{tests/test_plugins.py}（DemoRunTests 与 ParameterTrendTests 的 mesh 部分）}
  \cardsep
  \cardrow{可标定参数}{无}
\end{modulecard}

本模块\textbf{没有任何可标定参数}。它的四个参数（\akey{nx}、
\akey{target_spacing_nm}、\akey{min_nodes_per_layer}、
\akey{lateral_margin_um}）都是几何与数值分辨率的\emph{选择}，
而不是需要向实验数据拟合的物理量。这一区分在流水线层面是重要的：
自动拟合循环只应触碰可标定量，网格分辨率不在此列。

\section{问题提法与几何假设}
\label{sec:mesh-problem}
\index{1concept@概念!structured grid@结构化网格}

\subsection{输入：器件结构的描述}
\label{sec:mesh-input}

本模块消费 \akey{structure} 这一 artifact，其结构由 Deck 的
\texttt{DeviceSpec} 定义：

\begin{center}
\small
\begin{tabular}{@{}p{0.24\textwidth} p{0.66\textwidth}@{}}
\toprule
\textbf{字段} & \textbf{含义}\\
\midrule
\akey{name}, \akey{substrate} & 器件标识与衬底材料键\\
\akey{layers[]} & 外延层表，\emph{自表面向下}排列；每层给出
  \akey{name}, \akey{material}, \akey{role}, \akey{thickness_nm}, \akey{x_al},
  \akey{doping_cm3}, \akey{doping_kind}\\
\akey{gate} & \akey{length_um}, \akey{kind}, \akey{barrier_eV}, \akey{workfunction_eV}\\
\akey{geometry} & \akey{gate_width_um}, \akey{source_gate_spacing_um},
  \akey{gate_drain_spacing_um}, \akey{fingers}\\
\bottomrule
\end{tabular}
\end{center}

其中 \akey{role} 是层在电学上的功能标签（\akey{barrier}、\akey{channel}、
\akey{buffer} 等\index{1concept@概念!layer role@层功能标签}）。下游模块
\emph{按角色而非按层名}定位自己的作用区：\pid{physics.polarization.ambacher}
（\crefchap{ch:polarization}）通过 \akey{role="barrier"} 与
\akey{role="channel"} 找到势垒与沟道。校验层要求恰好一个势垒层与一个沟道层，
并要求势垒层的 \akey{x_al} 严格为正 —— 因为"势垒不含铝"在物理上是矛盾的。

\subsection{张量积假设}
\label{sec:mesh-assumption}
\index{1concept@概念!tensor-product grid@张量积网格}

网格被假定为\emph{张量积}形式：$\mathcal{G}=\mathcal{X}\times\mathcal{Z}$，
即横向节点集合与竖直节点集合的笛卡尔积。其等价表述是：竖直方向的分层只随
$z$ 变化，横向的加密只随 $x$ 变化，二者互不干涉。这一假设带来两个直接后果：

\begin{itemize}[leftmargin=1.7em,itemsep=0.3em,topsep=0.35em]
  \item \textbf{网格可用两个一维数组完整描述}，节点坐标是
        $(x_i,z_j)$ 的直积，无须存储连接关系；单元是规则矩形；
  \item \textbf{二维算子可以按行分块}：把横向索引排在前、竖直索引排在后，
        则竖直方向上的耦合只发生在相邻的两"行"之间，整个 Jacobian 天然是
        \emph{块三对角}的。
\end{itemize}

第二条是本模块对全流水线最重要的贡献。它把二维隐式求解的线性代数代价从
一般的 $O\big((n_xn_z)^{3}\big)$ 降到 $O\big(n_z(n_xm)^{3}\big)$
（$m$ 为每个横向节点携带的未知量数），见 \crefchap{ch:blocktridiag}。

代价是几何表达能力：张量积网格\emph{无法}表达倾斜台面、凹槽栅、再生长等
非规则结构。本模块接受的正是这一取舍，理由见 \S\ref{sec:mesh-limits}。

\section{网格构造的公式化}
\label{sec:mesh-formulation}
\index{3method@方法!grid generation@网格构造}

\subsection{竖直轴：逐层定节点数}
\label{sec:mesh-vertical}

设层表自表面向下依次为 $l=1,\dots,L$，第 $l$ 层厚度为 $t_l$。竖直轴按层分段构造，
每层的\emph{单元数}取

\begin{equation}
n_l \;=\; \max\!\left(N_{\min},\ \mathrm{round}\!\left(\frac{t_l}{h_{\text{target}}}\right)\right),
\label{eq:mesh-ncells}
\end{equation}

其中 $h_{\text{target}}$ 是目标节点间距（参数 \akey{target_spacing_nm}），
$N_{\min}$ 是每层单元数下限（参数 \akey{min_nodes_per_layer}，默认 $5$）。
层内节点\emph{均匀}：

\begin{equation}
z^{(l)}_j \;=\; z^{(l)}_{\text{top}} \;+\; t_l\,\frac{j}{n_l},
\qquad j=0,1,\dots,n_l ,
\label{eq:mesh-layer-nodes}
\end{equation}

而层顶坐标由游标累加得到：$z^{(1)}_{\text{top}}=0$，
$z^{(l+1)}_{\text{top}}=z^{(l)}_{\text{top}}+t_l$。相邻两层\emph{共享}分界节点，
因此全局节点总数为

\begin{equation}
N_z \;=\; 1+\sum_{l=1}^{L} n_l .
\label{eq:mesh-node-count}
\end{equation}

\crefeq{eq:mesh-ncells} 中的下限 $N_{\min}$ 是为极薄层设的保护：若某层厚度远小于
$h_{\text{target}}$，四舍五入会把该层的单元数压到 $0$ 或 $1$，使该层的材料属性
在离散意义下消失。下限把它强行撑到 $N_{\min}$ 格。这是一把双刃剑，
其副作用见 \S\ref{sec:mesh-limits}。

\subsection{横向轴：以栅足迹为锚}
\label{sec:mesh-lateral}

横向轴是\emph{均匀}的，其范围由三个接触的位置决定。以栅中心为原点，
源、漏接触位于

\begin{equation}
x_{S} \;=\; -\left(\frac{L_g}{2}+L_{sg}\right),
\qquad
x_{D} \;=\; +\left(\frac{L_g}{2}+L_{gd}\right),
\label{eq:mesh-contacts}
\end{equation}

$L_g$ 为栅长（\akey{gate.length_um}），$L_{sg}$、$L_{gd}$ 为源栅与栅漏间距。
再向外各留出边距 $M$（参数 \akey{lateral_margin_um}）：

\begin{equation}
x_{\min} \;=\; x_{S}-M, \qquad x_{\max} \;=\; x_{D}+M,
\qquad
\Delta x \;=\; \frac{x_{\max}-x_{\min}}{n_x-1},
\label{eq:mesh-lateral}
\end{equation}

节点为 $x_i=x_{\min}+i\,\Delta x$，$i=0,\dots,n_x-1$。
把栅中心固定在原点是一种\emph{约定}，好处是栅边缘的位置
$x=\pm L_g/2$ 可以直接读出，无须再做坐标平移。

\begin{equation}
N_{\text{node}} \;=\; n_x\,N_z,
\qquad
N_{\text{cell}} \;=\; n_x\,(N_z-1).
\label{eq:mesh-size}
\end{equation}

\section{构造算法与网格度量}
\label{sec:mesh-algorithm}
\index{2qty@物理量!node spacing@节点间距}

\subsection{算法}

构造过程是层表的单遍扫描：维护一个竖直游标，对每层算出 $n_l$、
追加 $n_l$ 个节点、记录该层的节点切片与几何元数据。全局复杂度为
$O(N_z+n_x)$，且只需线性存储。

每个层条目除几何量（\akey{top_nm}、\akey{bottom_nm}、\akey{thickness_nm}、
\akey{n_cells}、\akey{n_nodes}）之外，还记录 \akey{node_slice} =
$[\text{start},\text{end}]$。这不是冗余信息：它使\emph{反查}成为可能 ——
给定一个节点号即可知它属于哪一层，无须再次扫描层表。下游模块反复需要这一映射
（例如按层取介电常数、按层取有效质量），把层结构显式写进 artifact
比让每个消费者各自重算要可靠。

\subsection{网格度量与不变量}

构造结束后记录两组度量：竖直间距的极值 $\Delta z_{\min}$、$\Delta z_{\max}$
（它们直接决定显式时间步的稳定性上界与空间截断误差的量级），以及
\crefeq{eq:mesh-size} 的规模计数。另有两条结构性不变量，在测试中固化：

\begin{itemize}[leftmargin=1.7em,itemsep=0.3em,topsep=0.35em]
  \item \textbf{层界一定落在节点上。} 因为相邻层共享分界节点，
        界面坐标必然等于某个 $z_j$；
  \item \textbf{节点切片与轴自洽。} 对任意层，
        $z_{\text{start}}=\text{top\_nm}$ 且 $z_{\text{end}}=\text{bottom\_nm}$。
\end{itemize}

\subsection{一处刻意的取舍：界面落在节点上}
\label{sec:mesh-interface}
\index{1concept@概念!interface placement@界面落位}

第一条不变量值得单独讨论，因为它是本模块与
\pid{physics.quantum.sp1d}（\crefchap{ch:sp1d}）之间一处\emph{有意为之的
不一致}，且其后果是可测量的。

\crefchap{ch:sp1d} 的薛定谔--泊松求解器\emph{拒绝}把异质界面放在节点上：
它专门构造网格使界面落在单元之间的\emph{面}上，从而让导带带阶、介电跃变、
有效质量跃变与极化面电荷四件事同时落在单元边界，恢复二阶收敛
（实测 $p\approx2.08$，见 \crefchap{ch:sp1d}）。本模块则相反，界面落在节点上。

这不是疏忽，而是接口边界的取舍：本模块产出的是\emph{几何}，
供\crefchap{ch:transport_dd} 的有限体积漂移扩散使用；后者按设计
（设计稿 §3）\emph{原样消费}给定的网格，因此界面只能在节点上，
材料跃变在界面附近被一个单元宽度的平均所抹平。

抹平的后果是量化的，不是定性的：\crefchap{ch:transport_dd} 实测，
由该网格解出的二维电子气面密度比极化给出的上界
$\sigeff/q$ 高约 $1.7$ 倍。这一偏差在设计稿中已被预见，在
\crefchap{ch:transport_dd} 中首次被定量。对只依赖几何的下游用途
（接触定位、接入区电阻）它没有影响；对界面敏感的带边对齐则是硬伤。

\section{验证与基准}
\label{sec:mesh-verification}

本模块的验证是\emph{结构性的}：断言网格的拓扑与几何不变量，
而不是断言某个坐标的魔数。所有基准在 \fpath{tests/test_plugins.py} 中固化。

\begin{veriftable}{网格的结构不变量（演示器件：Al$_{0.25}$Ga$_{0.75}$N 20\,nm
/ GaN 100\,nm / GaN 1500\,nm，$L_g=0.15$\,\U{\mu m}，\akey{nx}=41）}{tab:mesh-invariants}
\vrow{规模自洽}{$N_{\text{node}}=n_x\,N_z$}{通过（$41\times1621=66461$）}
\vrow{横向数组长度}{$=n_x$}{$41$}
\vrow{竖直轴覆盖整栈}{$z_0=0$，$z_{N-1}=\sum_l t_l$}{$0\to1620$\U{nm}（精确）}
\vrow{竖直轴严格单调}{$z_j<z_{j+1}$}{通过（无重复、无倒序）}
\vrow{层与界面计数}{$L$ 层、$L-1$ 界面}{$3$ 层、$2$ 界面}
\vrow{界面坐标}{等于层厚累加值}{$20.0$ 与 $120.0$\U{nm}（精确）}
\vrow{节点切片与轴自洽}{$z_{\text{start}}=\text{top}$，$z_{\text{end}}=\text{bottom}$}{逐层通过}
\vrow{接触次序}{$x_S<x_{gl}<x_{gr}<x_D$}{通过}
\vrow{横向范围随栅漏间距单调}{$L_{gd}\uparrow\Rightarrow x_{\max}\uparrow$}{$1.0\to3.0$\U{\mu m} 时 $x_{\max}$ 增大}
\vnote{竖轴节点数由 \crefeq{eq:mesh-node-count} 给出：$1+20+100+1500=1621$，
      其中各层单元数由 \crefeq{eq:mesh-ncells} 在 $h_{\text{target}}=1$\U{nm} 下取整而得。}
\end{veriftable}

\section{误差来源与适用边界}
\label{sec:mesh-limits}
\index{1concept@概念!validity limits@适用边界}

\limitsheading

\begin{bounds}
  \item \textbf{界面被一个单元宽度抹平。} 见 \S\ref{sec:mesh-interface}。
        可测量后果：下游二维电子气面密度比极化上界高约 $1.7$ 倍。

  \item \textbf{横向均匀，无栅边缘加密。} 栅边缘的电场集中是本器件最重要的
        二维效应之一，而均匀横向网格无法在其附近加密。模块自身的说明
        也把栅边缘加密列为后续版本的计划。

  \item \textbf{张量积几何的表达能力有限。} 倾斜台面、凹槽栅、再生长结构、
        场板等一概无法表达。这是 \S\ref{sec:mesh-assumption} 所述取舍的另一面。

  \item \textbf{$N_{\min}$ 下限可能引入过密的网格。} 当某层远薄于
        $h_{\text{target}}$ 时，\crefeq{eq:mesh-ncells} 的下限会把它撑到
        $N_{\min}$ 格，使该层的实际间距远小于目标值。局部过密只影响效率，
        但会让 $\Delta z_{\min}$ 与 $\Delta z_{\max}$ 相差很大，
        在按 $\Delta z$ 做稳定性估计的下游算法中需要留意。

  \item \textbf{只覆盖外延栈。} 竖直轴止于层表最后一层的底面，
        衬底未建模，也没有背面接触。因此本网格不适用于自热或衬底耦合分析。

  \item \textbf{\akey{fingers} 只是元数据。} 多指器件的版图字段被读入并记录，
        但未参与几何构造：本模块只生成\emph{单指}截面。
        把指数与栅宽折算成总栅宽的工作属于提取阶段，不属于网格阶段。

  \item \textbf{掺杂只随层记录，不引起网格加密。} 缓冲层的受主掺杂
        （演示器件中 $1\times10^{17}$\U{cm^{-3}}）写在层条目里，
        但不会因此在该区域加密；对耗尽区宽度的分辨率由
        $h_{\text{target}}$ 单独控制。
\end{bounds}

\section{接口契约与小结}
\label{sec:mesh-interface-contract}
\index{5artifact@数据通道!structure}\index{5artifact@数据通道!mesh}

本模块读取一个 artifact、写出一个：

{\small
\begin{longtable}{@{}p{0.06\textwidth} p{0.22\textwidth} p{0.62\textwidth}@{}}
\toprule
\textbf{方向} & \textbf{键} & \textbf{本模块使用/提供的字段}\\
\midrule
\endfirsthead
\toprule
\textbf{方向} & \textbf{键} & \textbf{本模块使用/提供的字段}\\
\midrule
\endhead
\bottomrule
\endfoot
入 & \akey{structure} &
  \akey{layers[].{name, material, role, thickness_nm, x_al, doping_cm3, doping_kind}},
  \akey{gate.length_um}, \akey{geometry.{source_gate_spacing_um, gate_drain_spacing_um}}\\
出 & \akey{mesh} &
  \akey{kind}, \akey{nx}, \akey{nz}, \akey{n_nodes}, \akey{n_cells},
  \akey{x_nm[]}, \akey{z_nm[]}, \akey{layers[]}（含 \akey{node_slice}）,
  \akey{interfaces_nm[]}, \akey{total_depth_nm}, \akey{lateral_extent_nm},
  \akey{vertical_spacing_nm.{min,max}}, \akey{contacts_nm.*}\\
\end{longtable}
}

\noindent
产出物是一个\emph{普通字典}，下游模块无须依赖任何网格库即可消费它。
这是刻意的：网格是流水线上第一个"共享事实"，把它做成惰性数据而非
活对象，可以使任何模块（包括第三方插件）在不引入依赖的前提下读取几何。

\subsection*{小结}

本章的结论可以概括为三点。

其一，\textbf{结构化网格不是一个偷懒的选择，而是一次架构级的取舍}。
它换来的是块三对角 Jacobian 与 $O(n_zm^{3})$ 的线性步，代价是几何表达能力。

其二，\textbf{网格产出的不只是坐标，还有\emph{可反查的结构}}。
层条目携带节点切片，使"某节点属于哪一层"成为 $O(1)$ 查询；
这让下游不必各自重新解释层表。

其三，\textbf{界面落位是模块间接口的产物，不是局部最优}。
本模块把界面放在节点上，与 \crefchap{ch:sp1d} 的面基界面相反；
两者都对各自的求解器是恰当的，而它们之间的失配有一个已被量化的后果
（\crefchap{ch:transport_dd} 中的 $1.7$ 倍面密度偏差）。

\cleardoublepage
      % 第 6 章   mesh.structured2d            v0.1.0
% module_id: materials.dft
% module_version: 0.1.0
% plugin_dir: huntian/plugins/materials_dft
% source_files: huntian/plugins/materials_dft/{__init__,plugin,table_io,qe_driver,slurm}.py
% tests: tests/test_materials_dft.py, tests/test_materials_dft_table.py,
%        tests/test_materials_dft_overrides.py, tests/test_materials_dft_benchmarks.py,
%        tests/test_hotswap_materials.py

\chapter{第一性原理材料参数通道}
\label{ch:dft}

\begin{quote}
\small
\noindent\textbf{模块标识}\quad\pid{materials.dft}\qquad
\textbf{版本}\quad 0.1.0\qquad
\textbf{类别}\quad physics
\end{quote}

\noindent
本章记述材料库的第一性原理覆盖通道。它不引入新的物理模型，而是回答一个
更基础的问题：\emph{器件模型里那些材料常数是从哪里来的，可信到什么程度}。
本模块的整体设计围绕一个判断展开 —— 在器件级仿真中，材料常数的\textbf{系统性
误差}往往比数值格式误差大一个量级以上，因此对材料参数的\emph{来源、
不确定度与适用条件}的显式管理，比把求解器精度再提高一档更有价值。

\section{模块契约}
\label{sec:dft-contract}
\index{4module@模块!materials.dft}

\begin{modulecard}{materials.dft}{0.1.0}
  \cardrow{类别}{physics（材料数据，流水线阶段）}
  \cardrow{能力}{materials, materials.dft, materials.validation}
  \cardrow{依赖 artifact}{structure}
  \cardrow{产出 artifact}{materials, dft\_provenance}
  \cardrow{验证状态}{partially-verified}
  \cardrow{源码}{\fpath{huntian/plugins/materials_dft/}（多文件插件包）}
  \cardrow{测试}{\fpath{tests/test_materials_dft*.py} 等五份（含 D0--D3 基准与替换验收）}
  \cardsep
  \cardrow{可标定参数}{\textbf{无}（见下）}
  \cardrow{后端}{\akey{table}（默认，零第三方依赖）/ \akey{qe_local} / \akey{qe_slurm}（后两者本阶段未实现）}
\end{modulecard}

本模块\textbf{没有任何可标定参数}，且这一条是被测试固化的断言，而不是注释里的
约定。理由是方法论上的：\akey{ecut_wfc_ry}、\akey{kspacing_1_A}、
\akey{pseudopotential_family} 这些量是\emph{数值方法的选择}，
不是器件的物理参数。若允许自动拟合循环去调整它们，就等于\emph{为了对上数据
而调网格} —— 那是把标定与收敛性混为一谈，会系统性地破坏两者的可信度。
可标定量属于器件模型（如 \akey{screening}，见 \crefchap{ch:polarization}）。

\section{物理前提与精度红线}
\label{sec:dft-premise}

\subsection{泛函的系统误差不是"误差"，是偏移}
\label{sec:dft-functional}

密度泛函理论在设备级仿真的语境下，其价值不在于"更准"，而在于
\emph{可追溯}：每一个常数都能追到一次具体的计算设定与一个具体的出处。
但有一个必须显式处理的系统性偏移 —— 交换关联泛函对带隙的估计。

广义梯度近似（PBE）\cite{perdew1996} 对半导体带隙的系统性低估是
众所周知且量级巨大的：本仓库的查找表记录 GaN 的 PBE 原始带隙为
$1.72$\U{eV}，而实验值为 $3.42$\U{eV} —— 偏低约 $1.7$\U{eV}，
相对误差 $50\%$。这\emph{不是实现缺陷}，是泛函的已知行为。

问题在于这个数会被下游怎样使用。带隙在本流水线中流经两条通路：

\begin{itemize}[leftmargin=1.7em,itemsep=0.3em,topsep=0.35em]
  \item 导带带阶 $\dEc=r\,(E_g^{\text{barrier}}-E_g^{\text{channel}})$，
        进入 \pid{physics.quantum.sp1d}（\crefchap{ch:sp1d}）的
        薛定谔方程，直接决定势阱深度；
  \item 极化模型的带阶，进入 \crefchap{ch:polarization}。
\end{itemize}

把偏低 $1.7$\U{eV} 的带隙喂进去，势阱深度会错到使整个二维电子气
计算结果失去意义。因此本模块把"PBE 带隙不得进入器件模型"写成
\textbf{代码里的红线}，而不是文档里的建议。

\subsection{三条精度红线}
\label{sec:dft-redlines}
\index{1concept@概念!precision red line@精度红线}

\begin{enumerate}[leftmargin=1.9em,itemsep=0.4em,topsep=0.4em]
  \item \textbf{\akey{eg_policy="reject_pbe"}（默认）。}
        当 \akey{functional="pbe"} 与 \akey{eg_policy="reject_pbe"} 同时出现时，
        \texttt{validate()} 返回一条硬 issue，阶段拒绝运行。错误信息明确
        指出带隙的消费者（\pid{physics.quantum.sp1d} 与极化带阶）与原因。
        这不是警告，是拒跑。

  \item \textbf{\akey{eg_policy="scissor_to_reference"} 必须记账。}
        允许用 PBE 带隙，但只能以剪刀修正的形式，且平移量、原始 PBE 值
        与"验证等级被降级"三件事都写进溯源。见 \S\ref{sec:dft-scissor}。

  \item \textbf{无收敛证据的条目必须标注。}
        只有同时具备 \akey{ecut_wfc_ry} 与 \akey{kspacing_1_A}
        两条\emph{两点收敛记录}且声明了外推的条目，才被认作
        \akey{converged=true}。其余条目即便被采纳，也必须在覆盖报告中
        标为 \akey{converged=false}，并列入溯源的假设清单 ——
        \textbf{绝不静默使用}。
\end{enumerate}

第 3 条在本仓库的查找表上有一个诚实的后果：\fpath{huntian/plugins/materials_dft/data/algan_gan.json}
中的\emph{全部}条目都是 \akey{convergence: null}，因为逐篇文献的两点收敛
设置大多未以可提取的形式发表。查找表自己的 \akey{honesty_note} 写明：
所引不确定度区间\emph{跨越了已发表计算的散布}，它是文献收敛性的
\emph{代理指标}，不能替代真正的收敛记录。因此该表只有在用户显式选择
\akey{override_policy} 时才会被采纳。

\subsection{三种后端，一个契约}
\label{sec:dft-backends}
\index{1concept@概念!backend contract@后端契约}

\akey{backend} 有三个取值，它们共享同一份 artifact 契约：

{\small
\begin{longtable}{@{}p{0.14\textwidth} p{0.34\textwidth} p{0.42\textwidth}@{}}
\toprule
\textbf{后端} & \textbf{含义} & \textbf{本阶段状态}\\
\midrule
\endfirsthead
\toprule
\textbf{后端} & \textbf{含义} & \textbf{本阶段状态}\\
\midrule
\endhead
\bottomrule
\endfoot
\akey{table} & 服务一份预计算、带版本号的 JSON 查找表 &
  \textbf{已实现}，零第三方依赖\\
\akey{qe_local} & 经 ASE \cite{larsen2017} 驱动本机
  Quantum ESPRESSO \cite{giannozzi2009} &
  仅完成\emph{可用性探测}（\akey{ase} 可导入、\texttt{pw.x} 在 PATH 上）；
  计算流程未实现\\
\akey{qe_slurm} & 同上，提交到集群 &
  未实现\\
\end{longtable}
}

把三者放在同一个契约下的意义在于：一旦 \akey{qe_local} 落地，
下游模块\emph{一行都不用改}。本阶段的验收测试正是这一性质的证明
（\S\ref{sec:dft-verification}）。

\paragraph{惰性导入是硬要求。} 插件发现机制会加载每一个插件，因此加载本模块
\textbf{不得}把 \akey{ase}、\akey{qe_driver}、\akey{slurm} 拉进
\texttt{sys.modules}。测试用一个独立子进程断言这一点：在
\akey{table} 默认后端下，\texttt{sys.modules} 中不得出现这三者。
这一条保证了 JSON Deck 在任何环境（包括无 ASE、无 QE 的纯标准库环境）
都能运行。

\section{公式化}
\label{sec:dft-formulation}

\subsection{合金插值：Vegard 加弯曲项}
\label{sec:dft-vegard}
\index{3method@方法!Vegard 插值}

三元合金 Al$_x$Ga$_{1-x}$N 的常数由两个端点材料线性插值，并叠加一个
\emph{弯曲项}（bowing）修正：

\begin{equation}
V(x) \;=\; (1-x)\,V_{\text{GaN}} \;+\; x\,V_{\text{AlN}}
\;+\; b_V\,x\,(1-x),
\label{eq:vegard}
\end{equation}

$b_V$ 为属性 $V$ 的弯曲系数。弯出项 $x(1-x)$ 在两端点为零、在中点最大，
恰好刻画"合金不是端点性质的线性混合"这一物理事实。本仓库为 AlGaN
记录了两个非零弯曲系数：

\begin{equation}
b_{p_{\text{sp}}} = 0.021\ \U{C/m^2},
\qquad
b_{E_g} = 1.0\ \U{eV},
\label{eq:bowing}
\end{equation}

其余属性取 $b_V=0$，退化为狭义 Vegard 定律。带隙的弯曲系数在本器件族中
历来显著（文献值散布较大），$1.0$\U{eV} 属常用取值；它直接影响
$\dEc$，因此也直接影响势阱深度。

\crefeq{eq:vegard} 中 $x$ 是 AlN 摩尔分数，定义域 $[0,1]$ —— 越界即报错，
不做静默钳制。

\subsection{剪刀修正}
\label{sec:dft-scissor}
\index{3method@方法!scissor 修正}

当 \akey{eg_policy="scissor_to_reference"} 时，PBE 带隙被\emph{刚性平移}
到参考带隙（内置库的、来自实验的值），平移量为

\begin{equation}
\Delta_{\text{scissor}} \;=\; E_g^{\text{PBE}} \;-\; E_g^{\text{ref}} .
\label{eq:scissor}
\end{equation}

采纳的带隙即参考值 $E_g^{\text{ref}}$，而两个附带量被写入材料库：

\begin{equation}
\mathtt{dft\_eg\_uncorrected\_eV} \;=\; E_g^{\text{PBE}},
\qquad
\mathtt{dft\_eg\_shift\_eV} \;=\; \Delta_{\text{scissor}} .
\label{eq:scissor-record}
\end{equation}

两者都带 \akey{dft_} 前缀。这个前缀是契约的一部分：内置物理模块
不消费任何 \akey{dft_} 开头的键，因此原始 PBE 值与平移量只能被审计，
不可能被误用。

同时，验证等级被\textbf{降级}：若条目原为 \akey{verified}，
经剪刀修正后降为 \akey{partially-verified}。理由直白：
剪刀修正是一个\emph{已知不物理}的补救，它使带隙对上实验的同时，
必然破坏了带边与有效质量之间的相对关系。把它标成"已验证"是自欺。

\subsection{陈旧判据}
\label{sec:dft-stale}
\index{3method@方法!stale 判据}

查找表记录了一组\emph{名义生产设置}。若调用者请求的数值设置与表中记录不符，
则表被判定为\emph{陈旧}（stale）。判据是一个集合：

\begin{equation}
\mathcal{M} \;=\;
\Bigl\{\, k \in \mathcal{K}_{\text{stale}} \;:\; p_k \neq t_k \,\Bigr\},
\qquad
\mathcal{K}_{\text{stale}} = \{
\mathtt{pp}, \mathtt{ecut}, \mathtt{kspacing}, \mathtt{sqs\_seed}, \mathtt{spin\_orbit}
\},
\label{eq:stale}
\end{equation}

$p_k$ 为参数请求值、$t_k$ 为表中记录值；浮点比较用 $10^{-12}$ 的绝对容差。
$\mathcal{M}\neq\varnothing$ 且 \akey{allow_stale_cache=False}（默认）时，
校验返回硬 issue。

注意 \akey{functional} \emph{不在} $\mathcal{K}_{\text{stale}}$ 中。这不是遗漏：
泛函由带隙红线管辖，而不由陈旧判据管辖，且一份表完全可能合法地
同时持有 PBE 的结构常数与 HSE 的带隙（本仓库的表正是如此）。

$\mathcal{K}_{\text{stale}}$ 中的 \akey{kspacing_1_A} 是布里渊区积分的
$k$ 点密度（Monkhorst--Pack 间距 \cite{monkhorstPack1976}）。
把它纳入陈旧判据是必要的：$k$ 点密度与平面波截断是\emph{两个独立的}
收敛参数，只检查其中一个会让另一维度的欠收敛悄悄溜过去。

同一集合中的 \akey{pseudopotential_family} 默认取
\texttt{SSSP/PBEsol/efficiency@1.3.0}，即固体态赝势库的"效率"版本
\cite{prandini2018}。赝势族必须随结果一起记录：同一结构用不同赝势族
算出的晶格常数与弹性常数会有可观差异，而这种差异\emph{不会}在
任何收敛记录中显示出来。

\subsection{基准判据：相对误差与容差带}
\label{sec:dft-tolerance}

材料侧的自验证基准（D 阶梯）以\emph{实验值}为参考，相对误差定义为

\begin{equation}
e(v) \;=\; \frac{\abs{v - v_{\text{ref}}}}{\abs{v_{\text{ref}}}} ,
\label{eq:rel-error}
\end{equation}

并按属性类别给出容差带：

\begin{equation}
\text{D0}:\ e_a,e_c \le 0.01, \quad
\text{D1}:\ e_{C_{13}},e_{C_{33}} \le 0.10, \quad
\text{D2}:\ e_{p_{\text{sp}}} \le 0.15, \quad
\text{D3}:\ e_{E_g} \le 0.05 .
\label{eq:tolerance}
\end{equation}

容差带的宽窄反映了各属性的可计算程度：晶格常数是 DFT 最擅长的量（$1\%$），
弹性常数次之（$10\%$），自发极化依赖 Berry 相位方法且文献间分歧大（$15\%$），
而带隙必须用杂化泛函才可比于实验（$5\%$ 是针对 HSE 采纳值而言，
不是对 PBE 值的宽容）\cite{heyd2003}。

参考值在测试中\emph{独立于内置材料库}给出，因此对
\fpath{materials.json} 的任何错误编辑都不可能悄悄移动靶子。

\section{合并流程与实现}
\label{sec:dft-merge}
\index{3method@方法!material overlay@材料覆盖合并}

\subsection{覆盖策略三态}
\label{sec:dft-override}

\akey{override_policy} 决定每个属性最终由谁供给：

{\small
\begin{longtable}{@{}p{0.22\textwidth} p{0.30\textwidth} p{0.38\textwidth}@{}}
\toprule
\textbf{取值} & \textbf{非带隙属性} & \textbf{带隙}\\
\midrule
\endfirsthead
\toprule
\textbf{取值} & \textbf{非带隙属性} & \textbf{带隙}\\
\midrule
\endhead
\bottomrule
\endfoot
\akey{dft_wins}（默认） & 采纳 DFT 值 & 按泛函与 \akey{eg_policy} 判定\\
\akey{reference_wins} & 保留内置值（不覆盖） & 保留内置值\\
\akey{dft_only_gapless} & 采纳 DFT 值 & \emph{不}采纳（保留内置值）\\
\end{longtable}
}

第三态的存在理由值得说明：它表达了"我接受 DFT 的结构与弹性常数，
但我不想让 DFT 的带隙以任何形式（包括剪刀修正）影响我的器件模型"
这一立场。这是一个合理的保守选择，而把它表达为一个显式的策略值，
比让用户去手工编辑材料库要安全得多。

\subsection{合并规则}

合并遵循"\emph{叠加而非替换}"的原则：以内置库为底，把 DFT 层叠上去，
未被覆盖的属性原样保留。因此合并后的库仍然是一个完整的材料库，
而不是一份"只有 DFT 关心过的那几个键"的残缺品。

每条覆盖都被记入覆盖报告的一行，含材料名、属性名、内置值、DFT 值、
采纳方向（\akey{"dft"} 或 \akey{"reference"}）、相对偏差百分比、
条目的收敛标志与出处字符串。\textbf{没有出处字符串的覆盖行是不允许存在的}，
该性质被测试直接断言。

\subsection{溯源 artifact}
\label{sec:dft-provenance}
\index{5artifact@数据通道!dft\_provenance}

除 \akey{materials} 之外，本模块还写出 \akey{dft_provenance}，
其结构固定包含：

\begin{itemize}[leftmargin=1.7em,itemsep=0.28em,topsep=0.35em]
  \item \akey{backend}、\akey{functional}、\akey{pseudopotential_family}；
  \item \akey{numerics}：\akey{ecut_wfc_ry}、\akey{kspacing_1_A}、
        \akey{sqs_seed}、\akey{spin_orbit}，以及
        \akey{table_version} 与 \akey{table_content_hash}；
  \item \akey{entries}：逐材料的出处、验证等级、收敛标志、不确定度与假设清单；
  \item \akey{override_report}：逐属性的采纳决策（见上）；
  \item \akey{assumptions}：模块级的诚实边界清单（见 \S\ref{sec:dft-limits}）。
\end{itemize}

其中 \akey{table_content_hash} 是查找表\emph{文件内容}的散列
（前 16 位十六进制）。它使"这次仿真用的是哪一份表"成为可核对的量，
而不是靠文件路径与记忆。

\subsection{缓存键}

对 \akey{qe_*} 后端，结果缓存键被设计为

\begin{equation}
\begin{split}
k_{\text{cache}} \;=\; \mathrm{hash}\bigl(&\texttt{structure},\
  \texttt{functional},\ \texttt{pseudopotential},\\
  &\texttt{ecut},\ \texttt{kspacing},\ \texttt{sqs\_seed}\bigr),
\end{split}
\label{eq:cache-key}
\end{equation}

\noindent
\akey{sqs_seed} 进入键是刻意的：特殊准随机结构（SQS）的生成依赖随机种子，
若种子不进键，同一合金组分在不同运行中会命中不同的超胞，
可复现性随之失效。

\section{验证与基准}
\label{sec:dft-verification}

验证分三层：红线作为\emph{行为}被断言、材料数值对\emph{独立实验值}、
以及\emph{即插即用}层面的替换验收。

\subsection{红线与契约}

\begin{veriftable}{红线、契约与惰性导入}{tab:dft-redlines}
\vrow{PBE 带隙默认被拒}{返回 \akey{eg-policy} 硬 issue}{通过（阶段拒绝运行）}
\vrow{HSE06 默认通过校验}{无 issue}{通过}
\vrow{剪刀模式记录平移量（\akey{dft_eg_uncorrected_eV} 与
  \akey{dft_eg_shift_eV}）}{三项同时落库}{
  $E_g=3.42$\U{eV}、原始 $1.72$\U{eV}、平移 $-1.70$\U{eV}（精确到 $10^{-6}$）}
\vrow{剪刀模式降级验证：不高于 \akey{partially-verified}}{降级一档}{通过（由 \akey{verified} 起算）}
\vrow{陈旧数值默认被拒}{返回 \akey{stale-cache} 硬 issue}{通过（$90$ vs $60$\U{Ry}）}
\vrow{陈旧数值需显式接受}{无 issue 且假设清单含 stale}{通过（列出 \akey{ecut_wfc_ry}）}
\vrow{\textbf{无可标定参数}（\akey{calibratable_params()}）}{返回空}{空元组}
\vrow{\textbf{默认后端零依赖}}{\akey{table} 为默认；\akey{reject_pbe} 为默认；
  \akey{hse06} 为默认}{通过}
\vrow{\textbf{惰性导入无泄漏}（\texttt{sys.modules} 不含
  \akey{ase}/\akey{qe_driver}/\akey{slurm}）}{无泄漏}{通过（\texttt{LEAKED []}）}
\vnote{最后三行是本模块最具方法论意义的三条：它们把"不依赖第三方"
      "不允许拟合"与"不静默使用未收敛数据"从注释变成了可执行的断言。}
\end{veriftable}

\subsection{材料数值：D0--D3 对实验}

参考值为实验常数，与内置材料库无关，因此对 \fpath{materials.json}
的误编辑无法移动靶子。实测值为 \fpath{tests/golden/dft_benchmarks_d0_d3.json}
中固化的相对误差（金样本，任何漂移都会使测试失败）。

\begin{veriftable}{D0--D3 基准：DFT 覆盖后的相对误差（\crefeq{eq:rel-error}）}{tab:dft-benchmarks}
\vgroup{D0\ \ 晶格常数（容差 $1\%$）}
\vrow{GaN $a$}{$\le 0.01$}{$0.3449\%$}
\vrow{GaN $c$}{$\le 0.01$}{$0.3279\%$}
\vrow{AlN $a$}{$\le 0.01$}{$0.0321\%$}
\vrow{AlN $c$}{$\le 0.01$}{$0.0602\%$}
\vgroup{D1\ \ 弹性常数（容差 $10\%$）}
\vrow{GaN $C_{13}$}{$\le 0.10$}{$0.0000\%$}
\vrow{GaN $C_{33}$}{$\le 0.10$}{$1.7284\%$}
\vrow{AlN $C_{13}$}{$\le 0.10$}{$2.7778\%$}
\vrow{AlN $C_{33}$}{$\le 0.10$}{$1.8767\%$}
\vgroup{D2\ \ 自发极化（容差 $15\%$；须记录方法差异）}
\vrow{GaN $p_{\text{sp}}$}{$\le 0.15$}{$0.0000\%$}
\vrow{AlN $p_{\text{sp}}$}{$\le 0.15$}{$0.0000\%$}
\vgroup{D3\ \ 带隙（容差 $5\%$，仅对已采纳的 HSE 值）}
\vrow{GaN $E_g$}{$\le 0.05$}{$0.8772\%$}
\vrow{AlN $E_g$}{$\le 0.05$}{$3.0645\%$}
\vnote{GaN 的 $C_{13}$ 与两者的 $p_{\text{sp}}$ 相对误差恰为 $0$，
      因为所采纳的 DFT 值与参考实验值在表中取同一数值 —— 这是
      \emph{数值上的吻合}，不构成独立验证；真正独立的是那些非零项。}
\end{veriftable}

带隙红线在基准层面还有一条更强的断言：合并后的 $E_g$ \emph{不得}等于
PBE 原始值，同时原始值必须\emph{在库中可查}。实测：GaN 合并带隙
$3.45$\U{eV} $\neq 1.72$\U{eV}，而 \akey{dft_eg_uncorrected_eV}
精确等于 $1.72$\U{eV}；AlN 对应 $4.07$\U{eV}。此外，所引不确定度区间
必须覆盖实验值（或与之相距不超过 $0.6$\U{eV}），以保证引用的散布
没有把参考值排除在外。

\subsection{即插即用：替换验收}

本模块的设计目标是\emph{零改动替换材料槽}。验收方式与
\crefchap{ch:sp1d} 相同：同一器件、同一批下游模块，只改 Deck 里
\akey{materials} 槽的绑定，断言其余的一切不变。

\begin{veriftable}{材料槽替换验收（\fpath{decks/demo_hemt_sp1d.json} vs
\fpath{decks/demo_hemt_dft.json}）}{tab:dft-hotswap}
\vrow{新阶段排在首位}{stages[0].slot $=$ \akey{materials}}{通过}
\vrow{\textbf{只增加一个阶段，其余顺序不动}}{其余槽位序列完全相同}{
  \texttt{dft\_slots[1:] == sp\_slots}}
\vrow{调度来自 manifest 依赖而非槽位表}{\akey{requires} 仅含 \akey{structure}}{
  通过（\akey{produces} 含 \akey{materials}）}
\vrow{两次运行均成功}{status $=$ success}{通过}
\vrow{\textbf{下游模块完全一致}}{模块标识与版本逐一相同}{
  \akey{polarization}、\akey{quantum} 逐字段相同；\akey{extract:0}
  仍为 \pid{extract.fom}，\akey{results.fom} 仍然存在}
\vrow{下游确实消费了合并库}{GaN $a=3.178$\U{\AA}（内置为 $3.189$）}{
  通过（合并库被改动，内置库不受影响）}
\vrow{\textbf{替换确实改变了物理}}{
  $0.005 < \abs{\Delta\ns}/\ns < 0.20$}{落在带内（既非毫无影响，也非荒谬偏离）}
\vrow{两次运行的指纹不同}{两个 Deck 的指纹不等}{通过}
\vrow{同一 Deck 两次指纹相同}{指纹相等}{通过}
\vrow{溯源可核对}{每条覆盖都有出处；\akey{adopted} $\in$ \{dft, reference\}}{
  通过}
\end{veriftable}

\section{误差来源与适用边界}
\label{sec:dft-limits}
\index{1concept@概念!validity limits@适用边界}

\limitsheading

\begin{bounds}
  \item \textbf{本阶段的表是文献转录，不是本所计算。} 查找表的
        \akey{provenance} 明确声明：数值来自已发表的 DFT 文献，
        逐值标注出处与不确定度，\textbf{并非自行计算}。
        以 GaN/AlN 端点为例，晶格常数与带隙取自
        Kyrtsos 等的 HSE/PBE 平面波计算 \cite{kyrtsos2019}，
        自发极化取自 Berry 相位计算 \cite{bernardini1997,dreyer2016}，
        弹性常数取自第一性原理计算与综述的汇合区间 \cite{vurgaftman2003}。
        它的作用是"比凭记忆抄常数更硬的先验"，不是终点。

  \item \textbf{全部条目缺少两点收敛记录。} 每一篇文献的收敛设置
        大多未以可提取的形式发表，因此表中所有条目均为
        \akey{convergence: null}、\akey{converged=false}。
        所引不确定度是\emph{文献散布}这一代理指标，不是收敛证据。

  \item \textbf{0\,K、理想晶体、无缺陷。} 真实 MOCVD 材料的
        位错、杂质与组分涨落被吸收进所引不确定度，而未被建模。

  \item \textbf{无温度修正。} 热导率、热指数等热学键保持内置库取值，
        不随 DFT 覆盖改变。

  \item \textbf{合金值来自端点插值。} \crefeq{eq:vegard} 的 Vegard 加弯曲形式
        是对端点值的插值，不是合金的超胞计算。SQS 超胞统计属于
        \akey{qe_slurm} 阶段，本阶段未实现。

  \item \textbf{PBE 带隙仍然是"可被记录"的。} 红线禁止的是它\emph{进入
        器件模型}，不是禁止它出现在库中；\akey{dft_eg_uncorrected_eV}
        正是为审计而保留的。读者不应把它当作可用值。

  \item \textbf{\akey{qe_local} 与 \akey{qe_slurm} 未实现。} 仅完成
        可用性探测与清晰的失败信息。本阶段不存在"真正重算一遍"的能力，
        这也是表中"文献转录"声明必须保留的原因。

  \item \textbf{不确定度区间不是置信区间。} 区间 $[v_{\text{lo}},v_{\text{hi}}]$
        由已发表结果的散布给出，其宽度反映的是\emph{文献间的分歧}，
        不含任何统计意义。把它当作误差棒做不确定度传播是误用。
\end{bounds}

\section{接口契约与小结}
\label{sec:dft-interface}
\index{5artifact@数据通道!materials}

{\small
\begin{longtable}{@{}p{0.06\textwidth} p{0.22\textwidth} p{0.62\textwidth}@{}}
\toprule
\textbf{方向} & \textbf{键} & \textbf{本模块使用/提供的字段}\\
\midrule
\endfirsthead
\toprule
\textbf{方向} & \textbf{键} & \textbf{本模块使用/提供的字段}\\
\midrule
\endhead
\bottomrule
\endfoot
入 & \akey{structure} &
  仅需存在（本模块不读几何；依赖声明使调度把它排在极化与量子之前）\\
入 & \akey{materials}（种入） &
  内置材料库，作为叠加的底\\
出 & \akey{materials} &
  叠加后的完整材料库；被覆盖的键为 \akey{a_lattice_A}, \akey{c_lattice_A},
  \akey{C13_GPa}, \akey{C33_GPa}, \akey{e31_C_m2}, \akey{e33_C_m2},
  \akey{p_sp_C_m2}, \akey{m_eff}, \akey{Eg_eV}；
  另附 \akey{dft_uncertainty}、\akey{dft_eg_uncorrected_eV}、
  \akey{dft_eg_shift_eV}\\
出 & \akey{dft_provenance} &
  后端、泛函、赝势族、数值设置、表版本与内容散列、逐材料条目、
  覆盖报告、假设清单\\
\end{longtable}
}

\noindent
契约上有一处细节值得点明：\textbf{被覆盖的键是"物理量名"，不是"DFT 量名"}。
\akey{materials} 的键空间对下游是\emph{不区分来源}的 ——
\pid{physics.polarization.ambacher} 读 \akey{p_sp_C_m2} 时
并不知道也不关心它是内置值还是 DFT 值。这正是零改动替换得以成立的机制，
也是为什么审计信息必须\emph{另开}一个 artifact（\akey{dft_provenance}）
而不能混进材料库的键空间。

\subsection*{小结}

本章的结论可以概括为三点。

其一，\textbf{本模块管理的不是常数，而是常数的可信度}。
不确定度、收敛标志、出处字符串、表内容散列与覆盖报告，
构成了一个可审计的先验，而"文献转录"这一身份被主动声明而非掩盖。

其二，\textbf{红线必须在代码里。} "PBE 带隙不得进入器件模型"若只写在
文档里，迟早会被一次便利的调用绕过；写成 \texttt{validate()} 的硬 issue
之后，它成了一个不可绕过的门。

其三，\textbf{零改动替换是可以被机械验收的。} 替换材料槽后，
流水线只多一个阶段、下游模块的标识与版本逐一不变、同一 Deck 的指纹
可复现、而物理量确实发生了落在合理带内的移动 —— 这四条合起来，
才是"即插即用"的完整含义。

\cleardoublepage
          % 第 7 章   materials.dft                 v0.1.0

% ---- 第三部分 -------------------------------------------------------
\part{极化与量子}
% module_id: physics.polarization.ambacher
% module_version: 0.1.0
% plugin_dir: plugins/polarization_ambacher
% source_files: plugins/polarization_ambacher/__init__.py
% tests: tests/test_plugins.py, tests/test_materials_dft_benchmarks.py (D2)

\chapter{Ambacher 极化模型}
\label{ch:polarization}

\begin{quote}
\small
\noindent\textbf{模块标识}\quad\pid{physics.polarization.ambacher}\qquad
\textbf{版本}\quad 0.1.0\qquad
\textbf{类别}\quad physics
\end{quote}

\noindent
本章记述二维电子气的\emph{来源}。在 AlGaN/GaN 异质结构中，导电沟道不是靠
调制掺杂做出来的，而是靠\emph{极化}做出来的：自发极化与压电极化在异质界面
上的不连续，等效为一层束缚面电荷，正是这层面电荷把导带拉出一个势阱。
因此本模块是整条物理链的起点 —— \crefchap{ch:analytic_tw} 与
\crefchap{ch:sp1d} 的量子求解器消费的"面密度上界"，就是这里算出的
$\sigeff/q$。

\section{模块契约}
\label{sec:pol-contract}
\index{4module@模块!physics.polarization.ambacher}

\begin{modulecard}{physics.polarization.ambacher}{0.1.0}
  \cardrow{类别}{physics（物理模型，流水线阶段）}
  \cardrow{能力}{polarization, polarization.ambacher}
  \cardrow{依赖 artifact}{structure, materials}
  \cardrow{产出 artifact}{polarization}
  \cardrow{验证状态}{partially-verified}
  \cardrow{源码}{\fpath{plugins/polarization_ambacher/}}
  \cardrow{测试}{\fpath{tests/test_plugins.py}（解析值、趋势、屏蔽比例）}
  \cardsep
  \cardrow{可标定参数}{\akey{screening}\calib、\akey{relaxed_fraction}\calib}
  \cardrow{极性}{仅支持 Ga 面（\akey{polarity="ga_face"}），N 面显式拒绝}
\end{modulecard}

本模块有两个可标定量，其中 \akey{screening} 是整条极化链上
\textbf{最大的不确定度来源}（见 \S\ref{sec:pol-screening}）。
它的默认值未经论证，模块自己的文档也直言："没有公认取值，
必须用 Hall 数据标定"。

\section{极化物理与控制方程}
\label{sec:pol-physics}

\subsection{几何与极性约定}
\label{sec:pol-geometry}
\index{1concept@概念!Ga-face polarity@Ga 面极性}

纤锌矿 III--氮化物沿 $c$ 轴（$[0001]$）缺少反演对称，因此即使无应变也存在
\emph{自发极化}。Ga 面生长（Ga 原子在最外层的 (0001) 面）意味着
$c$ 轴自衬底指向表面，此时的极化方向与面电荷符号是确定的 —— 本模块
只实现这一极性，并在校验阶段对其他取值返回 \akey{unsupported} issue。

沿 $z$ 轴（自表面指向体内）的坐标约定与 \crefchap{ch:sp1d} 一致。
界面上的束缚面电荷密度由两侧极化强度的\emph{不连续}给出，
其符号约定决定了势阱的朝向；若符号取反，导带会被抬高而不是压低，
二维电子气随之消失。

\subsection{双轴应变}
\label{sec:pol-strain}
\index{1concept@概念!biaxial strain@双轴应变}

势垒层（Al$_x$Ga$_{1-x}$N）在沟道层（GaN）上赝晶生长时，面内晶格常数被
强制与沟道一致，产生\emph{双轴}应变。由于 $a(\text{AlGaN}) < a(\text{GaN})$，
势垒处于面内\emph{张}应变：

\begin{equation}
\varepsilon \;=\; f_{\text{relax}}\;
\frac{a_{\text{ch}} - a_{\text{bar}}}{a_{\text{bar}}} ,
\label{eq:strain}
\end{equation}

$a_{\text{bar}}$ 与 $a_{\text{ch}}$ 为两侧的面内晶格常数
（由 \crefeq{eq:vegard} 对合金插值得到），
$f_{\text{relax}}\in[0,1]$ 为\emph{弛豫度}（参数 \akey{relaxed_fraction}）：
$f_{\text{relax}}=1$ 表示完全赝晶（应变全部保留），
$0$ 表示完全弛豫（无应变）。以分母取势垒侧的 $a_{\text{bar}}$ 而非沟道侧，
是因为应变定义在\emph{被应变的那一层}上。

\subsection{压电极化}
\label{sec:pol-piezo}
\index{1concept@概念!piezoelectric polarization@压电极化}

沿 $c$ 轴的双轴应变诱发的压电极化为

\begin{equation}
\boxed{\;
P_{\text{pz}} \;=\; 2\,\varepsilon\,
\left(e_{31} \;-\; e_{33}\,\frac{C_{13}}{C_{33}}\right)
\;}
\label{eq:piezo}
\end{equation}

所有材料常数取\emph{势垒层}的合金值。\crefeq{eq:piezo} 的结构值得解读：
因子 $2$ 来自双轴应变在两个面内方向上等价作用；
括号内的组合 $e_{31}-e_{33}C_{13}/C_{33}$ 是\emph{受夹}压电响应 ——
面内应变同时通过泊松效应在 $c$ 轴上产生一个反向的应变分量，
其比例由弹性常数比 $C_{13}/C_{33}$ 给出。忽略这一项（即只用 $e_{31}$）
会高估压电贡献。

对 AlGaN 而言 $e_{31}<0$、$e_{33}>0$、$C_{13}/C_{33}\approx0.26$，
因此括号内为负，而 $\varepsilon>0$，故 $P_{\text{pz}}<0$ ——
压电极化与自发极化\emph{同向}叠加。这正是 Ga 面 AlGaN/GaN 能做出
高面密度二维电子气的原因。

\subsection{极化不连续与束缚面电荷}
\label{sec:pol-sheet}
\index{1concept@概念!bound sheet charge@束缚面电荷}

界面上的束缚面电荷等于两侧极化强度之差：

\begin{equation}
\Delta P \;=\; \left(P_{\text{sp}}^{\text{bar}} + P_{\text{pz}}^{\text{bar}}\right)
\;-\; P_{\text{sp}}^{\text{ch}},
\qquad
\sigpol \;=\; \abs{\Delta P},
\label{eq:sigma-pol}
\end{equation}

再乘上屏蔽因子得到\emph{有效}面电荷：

\begin{equation}
\boxed{\;\sigeff \;=\; s\,\sigpol\;}
\label{eq:sigma-eff}
\end{equation}

$\sigeff$ 是本模块交给下游的\textbf{唯一}电荷量。由
$\sigeff$ 立即得到面密度的\emph{硬上界}

\begin{equation}
\nsu \;=\; \frac{\sigeff}{q},
\label{eq:ns-bound}
\end{equation}

即"若势垒层内电场为零、全部极化电荷都被二维电子气补偿"时的极限。
实际面密度总是低于它，因为势垒层内必须存在电场以容纳有限厚度的势降。
\crefchap{ch:sp1d} 报告中给出的占据率（演示器件下 $0.83$ 左右）
正是 $\ns/\nsu$。

\subsection{屏蔽因子：一个诚实的未知}
\label{sec:pol-screening}
\index{2qty@物理量!screening factor@屏蔽因子}

\crefeq{eq:sigma-eff} 中的 $s$（参数 \akey{screening}，默认 $0.8$）
并不是一个从第一性原理推出的量，它\emph{吸收了表面态屏蔽的全部未知物理}：
真实器件表面存在施主型表面态，它们会俘获一部分极化电荷，
使到达沟道的有效电荷小于理论极化值。本流水线不含表面态模型
（那属于陷阱模块），于是这部分物理被压缩成一个可标定系数。

模块文档对此的处理是坦率的：\emph{没有公认取值}，
未经屏蔽的理论值在实践中会高估面密度，因此必须用 Hall 测量标定。
把不确定性集中到一个显式命名、显式标记为可标定的参数上，
比把它分散到若干个"经验修正系数"里要好 —— 至少它可被追责。

\subsection{导带带阶}
\label{sec:pol-offset}

异质界面的导带带阶由两侧带隙差给出：

\begin{equation}
\Delta E_g \;=\; E_g^{\text{bar}} - E_g^{\text{ch}},
\qquad
\dEc \;=\; r\,\Delta E_g ,
\label{eq:pol-offset}
\end{equation}

比例 $r$ 不在本模块中决定：本模块只交出 $\Delta E_g$（artifact 键
\akey{band_offset_eV}），由量子模块（\crefchap{ch:sp1d}）用参数
\akey{delta_ec_ratio}（默认 $0.7$）折算为 $\dEc$。
这样分工的理由是：$\Delta E_g$ 是材料属性，$r$ 是模型选择，
两者的可信度与标定方式都不同。

\section{公式化与合金插值}
\label{sec:pol-formulation}

本模块的全部材料常数来自材料库，合金值按 \crefeq{eq:vegard}
（Vegard 加弯曲项）插值。用于本模块的键为
\akey{p_sp_C_m2}、\akey{a_lattice_A}、\akey{e31_C_m2}、
\akey{e33_C_m2}、\akey{C13_GPa}、\akey{C33_GPa}、\akey{eps_r}、
\akey{m_eff}、\akey{Eg_eV}。其中只有 $p_{\text{sp}}$ 与 $E_g$
带非零弯曲系数（\crefeq{eq:bowing}）。

若把材料槽换成 \pid{materials.dft}（\crefchap{ch:dft}），
这些常数会被 DFT 覆盖值替换，而\emph{本模块一行不改} ——
这正是 \crefchap{ch:dft} 中"零改动替换"验收的另一半。

\subsection{演示器件的算例}
\label{sec:pol-worked}

下表给出演示器件（Al$_{0.25}$Ga$_{0.75}$N，$20$\U{nm} 势垒，
$f_{\text{relax}}=1$，$s=0.8$）逐步代入的结果。所有中间量都可在
代码中逐一复现，因此这张表同时是一个\emph{可核对的算例}。

\begin{veriftable}{演示器件的极化链逐步计算（$x=0.25$，材料常数取自内置库）}{tab:pol-worked}
\vrow{合金晶格常数 $a_{\text{bar}}$}{文献插值}{$3.16975$\U{\AA}}
\vrow{双轴应变 $\varepsilon$，\crefeq{eq:strain}}{$>0$ 且 $<0.02$}{$6.073\times10^{-3}$}
\vrow{合金自发极化 $P_{\text{sp}}^{\text{bar}}$}{文献插值}{$-0.04406$\U{C/m^2}}
\vrow{合金压电极化 $P_{\text{pz}}$，\crefeq{eq:piezo}}{$<0$}{$-9.196\times10^{-3}$\U{C/m^2}}
\vrow{沟道自发极化 $P_{\text{sp}}^{\text{ch}}$}{GaN 值}{$-0.034$\U{C/m^2}}
\vrow{极化不连续 $\Delta P$，\crefeq{eq:sigma-pol}}{---}{$-1.9258\times10^{-2}$\U{C/m^2}}
\vrow{极化面电荷 $\sigpol$}{$\in(1.10,1.30)\times10^{13}$\U{cm^{-2}}}{
  $1.2020\times10^{13}$\U{cm^{-2}}}
\vrow{有效面电荷 $\sigeff$（$s=0.8$）}{$=0.8\,\sigpol$（$10^{-6}$ 位）}{
  $1.5407\times10^{-2}$\U{C/m^2}}
\vrow{面密度上界 $\nsu=\sigeff/q$}{---}{$9.616\times10^{12}$\U{cm^{-2}}}
\vrow{合金带隙 $E_g^{\text{bar}}$（$b_{E_g}=1.0$\U{eV}）}{---}{$4.285$\U{eV}}
\vrow{导带带阶 $\Delta E_g$}{---}{$0.865$\U{eV}}
\vnote{表中最后两行交出的是 \(\Delta E_g\) 而非 \(\dEc\)：
      \(r=0.7\) 的折算属于\crefchap{ch:sp1d}，在此处折算会掩盖分工。
      面电荷一行的判据即 \fpath{tests/test_plugins.py} 中的断言区间。}
\end{veriftable}

\section{计算流程与实现}
\label{sec:pol-implementation}
\index{3method@方法!alloy interpolation@合金插值}

\subsection{流程}

计算是无迭代的纯代数流程：取势垒层与沟道层 $\to$ 插值合金常数 $\to$
算应变与压电极化 $\to$ 求极化不连续与有效面电荷 $\to$ 求带阶 $\to$
打包 artifact。没有离散化，没有求解器，因此本模块的"离散化"一节
（\crefsec{sec:pol-formulation}）实际上是\emph{公式化}一节：
它的可验证内容在于公式与常数，而不在于格式。

\subsection{校验}

校验阶段做三件事：确认层表中同时存在 \akey{barrier} 与 \akey{channel}
两个角色；确认极性为 Ga 面；以及——在\emph{下游}模块中——
确认 $\sigeff>0$。最后一条之所以放在下游，
是因为"极化电荷为零"这一情形下量子求解器无解，
应当由消费者在自己的校验里给出与场景相关的提示
（\crefchap{ch:sp1d} 与 \crefchap{ch:analytic_tw} 都这样做，
并建议检查屏蔽因子与 Al 组分）。

\section{验证与基准}
\label{sec:pol-verification}

极化模型的验证采取"断言解析值 + 断言趋势符号"的组合。
数值解析值来自文献 \cite{ambacher1999,ambacher2000,ambacher2002,vurgaftman2003}；
趋势断言则保证参数变化时\emph{方向}正确，因而在参数演化后依然有效。

\begin{veriftable}{极化量的解析值与趋势}{tab:pol-verify}
\vgroup{解析值与内部自洽}
\vrow{面电荷 $\sigpol$（$x=0.25$）}{$\in(1.10,1.30)\times10^{13}$\U{cm^{-2}}}{
  $1.2020\times10^{13}$\U{cm^{-2}}}
\vrow{应变为张应变且量级合理}{$0<\varepsilon<0.02$}{两者同时成立}
\vrow{屏蔽严格按比例缩放}{$\sigeff=0.8\,\sigpol$（$10^{-6}$ 位）}{通过}
\vrow{压电极化为负}{$P_{\text{pz}}<0$}{通过}
\vrow{势垒侧自发极化低于沟道侧}{$P_{\text{sp}}^{\text{bar}}<P_{\text{sp}}^{\text{ch}}$}{
  $-0.04406 < -0.034$\U{C/m^2}}
\vgroup{趋势符号}
\vrow{Al 组分升高 $\Rightarrow$ 极化增强}{$\sigpol$ 单调增}{
  $x:0.20\to0.35$ 时 $\sigpol$ 增大}
\vrow{屏蔽增强 $\Rightarrow$ 面密度升高}{$\ns$ 单调增}{
  $s:0.4\to0.95$ 时 $\ns$ 增大}
\vrow{势垒增厚 $\Rightarrow$ 面密度升高}{$\ns$ 单调增}{
  $d:15\to30$\U{nm} 时 $\ns$ 增大}
\vrow{弛豫度降低 $\Rightarrow$ 极化减弱}{$\sigpol$ 单调增}{
  $f_{\text{relax}}:0.5\to1.0$ 时 $\sigpol$ 增大}
\vrow{屏蔽增强 $\Rightarrow$ 阈值更负}{$V_{\text{th}}$ 单调减}{
  $s:0.6\to0.9$ 时 $V_{\text{th}}$ 更低}
\vnote{最后一行涉及 \pid{extract.fom}（\crefchap{ch:fom}）的阈值估计，
      在此列出是因为它检验的正是 $\sigeff$ 的符号与量级。}
\end{veriftable}

\begin{veriftable}{自发极化的独立对照（\crefchap{ch:dft} 的 D2 基准）}{tab:pol-d2}
\vrow{GaN $p_{\text{sp}}$ 相对实验}{\crefchap{ch:dft} 的 D2 容差 $15\%$}{$0.0000\%$}
\vrow{AlN $p_{\text{sp}}$ 相对实验}{同上}{$0.0000\%$}
\vnote{两个零并非"完美吻合"的证明，而是数值上取同一参考值的必然结果；
      真正独立的信息是表中引用的 Berry 相位方法被记录下来
      （设计稿 §8 对 D2 的方法记录要求）。其代价与局限见 \crefchap{ch:dft}。}
\end{veriftable}

\section{误差来源与适用边界}
\label{sec:pol-limits}
\index{1concept@概念!validity limits@适用边界}

\limitsheading

\begin{bounds}
  \item \textbf{屏蔽因子没有公认取值。} 默认 $s=0.8$ 只是起点，
        必须用 Hall 数据标定。它吸收了表面态物理的全部未知，
        是极化链上最大的不确定度来源（\S\ref{sec:pol-screening}）。

  \item \textbf{材料常数未经独立验证。} 极化常数、弹性常数、带隙与
        晶格常数转录自文献 \cite{ambacher2000,vurgaftman2003}，
        在材料库中标记 \texttt{unverified}；
        \pid{materials.dft}（\crefchap{ch:dft}）提供了一个带不确定度
        与溯源的覆盖通道，但其自身也是文献转录。
        跑通 L2（Hall 对比）之前，\emph{绝对数值不可信}。

  \item \textbf{只支持 Ga 面。} N 面极性在校验阶段被显式拒绝，
        而不是被静默地按 Ga 面处理 —— 后者的错误不会以异常的形式出现，
        只会以"面密度符号不对"的形式出现，更难发现。

  \item \textbf{无表面态模型。} 表面施主态的俘获、费米能级钉扎、
        以及由此产生的栅滞后（gate lag）均不在模型内。
        这部分物理被 \akey{screening} 一个标量所替代，
        因此它无法描述偏压依赖的屏蔽变化。

  \item \textbf{无应变弛豫的动力学。} $f_{\text{relax}}$ 是输入，
        不是求解结果；超过临界厚度的弛豫、位错引入的应变释放
        不在模型内。

  \item \textbf{合金插值是端点插值。} 见 \crefeq{eq:vegard} 与
        \crefchap{ch:dft} 的边界说明；簇团效应与短程有序不在其中。

  \item \textbf{极性依赖的压电符号未做交叉验证。} 本模块的符号约定
        来自 Ga 面 \cite{ambacher2000}，并通过"应变与面密度上界均为正"
        的断言间接验证；但没有与独立来源的符号约定做过逐项比对。
\end{bounds}

\section{接口契约与小结}
\label{sec:pol-interface}
\index{5artifact@数据通道!polarization}

{\small
\begin{longtable}{@{}p{0.06\textwidth} p{0.22\textwidth} p{0.62\textwidth}@{}}
\toprule
\textbf{方向} & \textbf{键} & \textbf{本模块使用/提供的字段}\\
\midrule
\endfirsthead
\toprule
\textbf{方向} & \textbf{键} & \textbf{本模块使用/提供的字段}\\
\midrule
\endhead
\bottomrule
\endfoot
入 & \akey{structure} &
  \akey{layers[].{role, material, thickness_nm, x_al}}（取 \akey{barrier}
  与 \akey{channel} 两层）\\
入 & \akey{materials} &
  上列九个材料键（含合金插值所需的端点）\\
出 & \akey{polarization} &
  \akey{x_al}, \akey{barrier_thickness_nm}, \akey{channel_material},
  \akey{eps_r_barrier}, \akey{m_eff_barrier},
  \akey{p_sp_barrier_C_m2}, \akey{p_sp_channel_C_m2},
  \akey{p_pz_barrier_C_m2}, \akey{delta_p_C_m2},
  \akey{sigma_pol_C_m2}, \akey{sigma_eff_C_m2},
  \akey{sigma_pol_cm2}, \akey{sigma_eff_cm2},
  \akey{strain}, \akey{relaxed_fraction}, \akey{screening},
  \akey{band_offset_eV}, \akey{lattice_a_barrier_A},
  \akey{lattice_a_channel_A}, \akey{valley_degeneracy}\\
\end{longtable}
}

\noindent
下游有两个消费者，各自只需要其中一部分：\pid{physics.quantum.sp1d}
（\crefchap{ch:sp1d}）取 $\sigeff$、势垒厚度、合金介电常数与有效质量
以及 $\Delta E_g$；\pid{extract.fom}（\crefchap{ch:fom}）另外取
$\sigpol$ 与 $\sigeff$ 的 cm$^{-2}$ 表示。把它们都放在\emph{一个}
artifact 里而不是拆成两个，是因为它们来自同一次计算、共享同一组常数，
拆开只会制造不一致的机会。

\subsection*{小结}

本章的结论可以概括为三点。

其一，\textbf{二维电子气的来源是极化，不是掺杂。} 由
\crefeq{eq:strain}--\crefeq{eq:sigma-eff} 得到的面电荷上界
$\nsu=\sigeff/q$ 是整条物理链上第一个可计算的硬约束，
也是后续所有量子模型的物理上限。

其二，\textbf{公式中的每一项都有物理职责。}
\crefeq{eq:piezo} 中受夹压电响应 $e_{33}C_{13}/C_{33}$ 不可省略；
\crefeq{eq:sigma-pol} 中两侧自发极化必须都算；
$f_{\text{relax}}$ 与 $s$ 虽都是标量，却代表完全不同的物理
（应变保留比例 vs.\ 表面态屏蔽），把它们混为一谈会失去可标定性。

其三，\textbf{不确定度应当被集中而非分散。} 本模块把一个真正的未知
（表面态屏蔽）显式命名、显式标记为可标定，而不是把它摊派成
若干个互相纠缠的"修正系数"。这样 Hall 标定才有明确的目标。

\cleardoublepage
  % 第 8 章   physics.polarization.ambacher v0.1.0
% module_id: physics.quantum.analytic_tw
% module_version: 0.1.0
% plugin_dir: huntian/plugins/quantum_analytic_tw
% source_files: huntian/plugins/quantum_analytic_tw/__init__.py
% tests: tests/test_plugins.py, tests/test_hotswap_quantum.py

\chapter{解析三角阱基线}
\label{ch:analytic_tw}

\begin{quote}
\small
\noindent\textbf{模块标识}\quad\pid{physics.quantum.analytic_tw}\qquad
\textbf{版本}\quad 0.1.0\qquad
\textbf{类别}\quad physics
\end{quote}

\noindent
本章记述一个\emph{解析}模型：用三角阱近似与电荷控制关系，闭式地估计
二维电子气的面密度与两个特征能量。它在流水线中扮演两个角色 ——
既是一个够用的快速估计，也是 \crefchap{ch:sp1d} 的自洽求解器所要替换的
\textbf{基线}。后一个角色是本项目架构主张的第一次实战验收：
两个模块声明同一能力、产出同一组 artifact 键，替换只改 Deck 里的一行绑定。

\section{模块契约}
\label{sec:tw-contract}
\index{4module@模块!physics.quantum.analytic\_tw}

\begin{modulecard}{physics.quantum.analytic\_tw}{0.1.0}
  \cardrow{类别}{physics（物理模型，流水线阶段）}
  \cardrow{能力}{quantum, quantum.analytic, quantum.triangular\_well}
  \cardrow{依赖 artifact}{polarization}
  \cardrow{产出 artifact}{quantum, ns}
  \cardrow{验证状态}{partially-verified}
  \cardrow{源码}{\fpath{huntian/plugins/quantum_analytic_tw/}}
  \cardrow{测试}{\fpath{tests/test_plugins.py}、\fpath{tests/test_hotswap_quantum.py}}
  \cardsep
  \cardrow{可标定参数}{\akey{phi_b_0_eV}\calib、\akey{phi_b_x_eV}\calib、
    \akey{delta_ec_ratio}\calib、\akey{gamma0_eV_m43}\calib}
  \cardrow{算法}{阻尼不动点迭代（无需求解器协作者）}
\end{modulecard}

本模块的说明书里有一句加粗的边界声明，值得原样保留：
\emph{"这是解析三角阱估计，\textbf{不是}自洽的薛定谔--泊松求解"}。
把这句话写进模块契约而不是文档脚注，是因为这两者的输出格式完全相同 ——
下游无法从数据上区分它们，只能从版本号与 artifact 里的 \akey{model} 字段区分。

\section{电荷控制模型}
\label{sec:tw-model}
\index{1concept@概念!charge control@电荷控制}

\subsection{三角阱近似}
\label{sec:tw-triangular}
\index{1concept@概念!triangular well@三角阱}

在强极化面电荷的作用下，沟道侧的导带被拉出一个深而窄的势阱。
若假定势阱内的电场\emph{均匀}（等价于忽略二维电子气自身的屏蔽对势形的
反馈），则势能呈\emph{线性}—— 这就是三角阱近似。在此近似下基态本征能量
有熟知的标度形式

\begin{equation}
E_0 \;=\; \gamma_0\,\ns^{2/3},
\label{eq:e0}
\end{equation}

$\gamma_0$ 是与有效质量、介电常数与电场标度有关的常数，
在本仓库中取 GaN 的常用值 $\gamma_0 = 2.123\times10^{-12}$\U{eV\cdot m^{4/3}}
\cite{ambacher2000}（参数 \akey{gamma0_eV_m43}）。$\ns^{2/3}$ 的标度
来自 Airy 函数的零点：三角阱中 $E\propto F^{2/3}$，而电场 $F$ 又正比于
$\ns$。

\subsection{电荷控制关系}
\label{sec:tw-charge-control}

器件结构自栅极到沟道是一串串联的电压降。设势垒层厚度 $d$、
相对介电常数 $\epsr$，势垒层内的电场由"未被二维电子气补偿的那部分极化电荷"
决定。在完全耗尽的假定下，栅压需承担肖特基势垒、势垒层压降与
费米能级抬升三项，于是势垒层内的面电荷为
$\sigeff - q\ns$，压降为
$(\sigeff-q\ns)\,d/(\epsr\epsO)$。令它等于 $\phib + \Ef - \dEc$，解出

\begin{equation}
\boxed{\;
\ns \;=\; \frac{\sigeff}{q}
\;-\; \frac{\epsr\,\epsO}{q\,d}\,
\Bigl(\phib + \Ef - \dEc\Bigr)
\;}
\label{eq:charge-control}
\end{equation}

全部能量以\U{eV}计、全部电势以\U{V}计，二者数值相等，
故 \crefeq{eq:charge-control} 中不出现额外的换算因子。
第一项是\emph{无势垒压降}的极限

\begin{equation}
\nsu \;=\; \frac{\sigeff}{q},
\label{eq:tw-upper}
\end{equation}

即 \crefeq{eq:ns-bound} 给出的硬上界；第二项是有限势垒厚度必须付出的代价。

\subsection{费米能级与简并}
\label{sec:tw-fermi}

二维电子气的面密度与费米能级之间是二维态密度下的 Fermi--Dirac 关系。
在"基态占据 + 简并抬升"的近似下（参数 \akey{include_fermi}，
默认开启），

\begin{equation}
\Ef \;=\; E_0 \;+\; \frac{\pi\hb^{2}}{q\,\meff\,g_v}\,\ns ,
\label{eq:tw-ef}
\end{equation}

$\meff$ 为势垒材料的有效质量\footnote{这是一个应当指出的不一致：
      在 \crefeq{eq:tw-ef} 中，二维电子气位于\emph{沟道}，
      因此态密度质量原则上应取沟道材料的 $\meff$，而实现取的是
      极化模块交出的势垒侧 $\meff$。对 AlGaN/GaN，
      两者相差 $0.30$ 与 $0.20$，因此这一项被高估约 $50\%$。
      由于该项本身只占 $\Ef$ 的约五分之一（见 \creftab{tab:tw-vs-sp}），
      对 $\ns$ 的影响在千分位量级，故未修改；此处如实记录。}，
$g_v$ 为谷简并度。$g_v$ 的引入是把二维态密度按能谷数成倍放大，
是谷简并度在本模型中的\emph{唯一}作用。

\subsection{自洽方程组}

\crefeq{eq:e0}、\crefeq{eq:tw-ef} 与 \crefeq{eq:charge-control} 构成一个关于
$\ns$ 的标量非线性方程：

\begin{equation}
\ns \;=\; F(\ns),
\qquad
F(\ns) \;=\; \nsu - \frac{\epsr\epsO}{q\,d}
\Bigl(\phib + E_0(\ns) - \dEc\Bigr)
- \frac{\epsr\epsO}{d}\cdot\frac{\pi\hb^{2}}{q^{2}\meff g_v}\,\ns ,
\label{eq:tw-fixedpoint}
\end{equation}

其中第二、三项分别来自 $E_0$ 与 \crefeq{eq:tw-ef} 的简并项。
$F$ 关于 $\ns$ 单调\emph{下降}：$\ns$ 增大使势阱变浅（$E_0$ 上升）
且费米能级抬升，两者都要求更大的栅压来维持，故可容纳的电荷减少。
单调性保证不动点唯一，也保证迭代只要收敛，收敛到的就是唯一的物理根。

\section{不动点映射与阻尼迭代}
\label{sec:tw-formulation}

\subsection{裸映射的增益}

把 \crefeq{eq:tw-fixedpoint} 写成不动点形式后，其收缩性由
$\abs{\mathrm{d}F/\mathrm{d}\ns}$ 决定。由 $E_0=\gamma_0\ns^{2/3}$，

\begin{equation}
\der{E_0}{\ns} \;=\; \tfrac{2}{3}\,\gamma_0\,\ns^{-1/3},
\label{eq:tw-de0}
\end{equation}

在演示器件的量级（$\ns\sim8\times10^{12}$\U{cm^{-2}}）下，
$\mathrm{d}E_0/\mathrm{d}\ns \sim 10^{-14}$\U{eV\cdot m^{2}}，
乘上 $\epsr\epsO/(q d)\sim10^{13}$\U{m^{-2}/eV} 后约为 $0.1$ 量级；
简并项给出另一个量级相近的正贡献。二者之和决定了映射的收缩率 ——
它是一个\emph{收缩}映射，但收缩率不算小，因此裸迭代收敛慢且对初值敏感。

\subsection{阻尼迭代}

实现采用带阻尼的不动点迭代（参数 \akey{relaxation}，默认 $0.5$）：

\begin{equation}
\ns^{(k+1)} \;=\; \ns^{(k)} + \beta\,\bigl(F(\ns^{(k)}) - \ns^{(k)}\bigr),
\label{eq:tw-iteration}
\end{equation}

初值取上界的一半 $\ns^{(0)}=\tfrac12\nsu$。
收敛判据是\emph{相对}判据

\begin{equation}
\abs{\ns^{(k+1)}-\ns^{(k)}} \;\le\; \varepsilon\cdot\max(\nsu,1),
\qquad \varepsilon = 10^{-10}\ \text{（参数 \akey{tolerance}）},
\label{eq:tw-convergence}
\end{equation}

取 $\max(\nsu,1)$ 而非 $\nsu$ 是为了在 $\sigeff\to0$（无势阱）时
避免零除。迭代上限 \akey{max_iterations} 默认 $500$，
超出即报告未收敛，而不是静默返回最后一个值。

收敛后，$E_0$ 与 $\Ef$ 在最终的 $\ns$ 处被\emph{重新计算一遍}。
这一步不是形式：迭代过程中最后一次更新可能已经跳出容差，
用旧的 $E_0$、$\Ef$ 打包会与 $\ns$ 不自洽。

\subsection{肖特基势垒高度}

栅势垒高度有两个来源，由参数 \akey{barrier_source} 控制：
Deck 的 \akey{device.gate.barrier_eV}（默认优先），
或解析拟合 $\phib = \phib^{(0)} + \phib^{(x)}\,x$
（默认 $0.84 + 0.39x$\U{eV}）。哪一路被采用会写进 artifact 的
\akey{phi_b_source} 字段 —— 演示器件取 Deck 值 $0.94$\U{eV}。

\section{验证与基准}
\label{sec:tw-verification}

\begin{veriftable}{解析基线的不变量、契约与趋势}{tab:tw-verify}
\vgroup{物理不变量}
\vrow{面密度落在物理区间}{$4.0\times10^{12}<\ns<1.2\times10^{13}$\U{cm^{-2}}}{通过}
\vrow{面密度低于极化上界}{$\ns<\nsu$}{通过（占据率 $>0$）}
\vrow{费米能级不低于基态}{$\Ef\ge E_0$}{通过}
\vrow{不动点收敛}{converged 且迭代数 $<500$}{通过}
\vrow{势垒高度取自 Deck}{\akey{phi_b_source} $=$ \akey{deck}；
  $\phib=0.94$\U{eV}}{精确到 $10^{-9}$}
\vgroup{趋势符号}
\vrow{屏蔽增强 $\Rightarrow$ $\ns$ 升高}{单调增}{通过}
\vrow{势垒增厚 $\Rightarrow$ $\ns$ 升高}{单调增}{通过}
\end{veriftable}

\subsection{与自洽解的定量对照}

解析模型的价值只有在与自洽解对照时才看得清。同一器件
（Al$_{0.25}$Ga$_{0.75}$N/GaN，$20$\U{nm} 势垒，$L_g=150$\U{nm}，
$s=0.8$）下两者的差别具有明确的\emph{物理形态}：

\begin{veriftable}{解析三角阱与自洽解的对照（同器件）}{tab:tw-vs-sp}
\vrow{量}{解析三角阱}{自洽 SP（差）}
\vrow{$\ns$ (\U{cm^{-2}})}{$7.58287\times10^{12}$}{$7.95848\times10^{12}$（$+4.95\%$）}
\vrow{$\sigeff/q$ 占据率}{$0.7886$}{$0.8276$（$+4.95\%$）}
\vrow{$\Ef$ 距界面带边 (\U{eV})}{$0.461018$}{$0.308840$（$-33.0\%$）}
\vrow{$E_0$ 距界面带边 (\U{eV})}{$0.380340$}{$0.229336$（$-39.7\%$）}
\vrow{$\Ef-E_0$ (\U{eV})}{$0.0806778$}{$0.0795039$（$-1.46\%$）}
\vnote{对照数据与 \crefchap{ch:sp1d} 的 \creftab{tab:sp1d-vs-analytic} 为同一组，
      此处从解析模型一侧重述。}
\end{veriftable}

\noindent
读法很关键：两个模型对\emph{势}的看法差别很大（$\Ef$ 与 $E_0$ 都掉了约
三分之一），但对\emph{统计}的看法收敛到同一处 —— 决定 $\ns$ 的亚带间距
$\Ef-E_0$ 只差 $1.46\%$。这解释了为什么一个没有解任何微分方程的
解析式能给出仅偏 $5\%$ 的面密度：它恰好把"势阱深度"与"统计抬升"
这两个误差抵消在了同一个方向上。

\subsection{即插即用：被替换的那一侧}

本模块最重要的验收不在自身，而在于它\emph{能被替换}
（\fpath{tests/test_hotswap_quantum.py}）：

\begin{veriftable}{能力与 artifact 契约的替换验收}{tab:tw-hotswap}
\vrow{两个模块声明同一能力}{均含 \texttt{quantum}}{通过}
\vrow{产出同一组契约键}{前七个键逐一相同}{通过}
\vrow{下游零改动}{\pid{extract.fom} 与流水线均不变}{通过}
\vrow{版本钉住}{解析 Deck 的 \texttt{\^{}0.1.0} \emph{不会}被解析到 sp1d}{通过}
\vrow{第三方实现可热插拔}{自定义实现同样可挂上}{通过}
\vnote{第三个是本项目"模块即插件"主张的核心证据：
      替换一个物理模型所需的工作量恰好是 Deck 里的一行绑定。}
\end{veriftable}

\section{误差来源与适用边界}
\label{sec:tw-limits}
\index{1concept@概念!validity limits@适用边界}

\limitsheading

\begin{bounds}
  \item \textbf{本模型不含任何微分方程。} 它不解析波函数、不解泊松方程、
        不做自洽。因此它给出的不是"近似解"，而是\emph{另一个模型的解}。
        凡涉及势形、波函数或有多个亚带的问题，本模型在原理上无能。

  \item \textbf{忽略波函数向势垒的穿透。} 这是三角阱近似的直接后果；
        \crefchap{ch:sp1d} 的自洽解显示势阱实际更浅、波函数有可观的
        势垒穿透，使 $E_0$ 与 $\Ef$ 都低约三分之一。

  \item \textbf{忽略多亚带占据。} 只有一个基态被占据，
        而自洽解给出的多亚带阶梯与真实的有效质量都不在本模型内。

  \item \textbf{态密度质量取的是势垒侧参数。} 见 \crefeq{eq:tw-ef}
        的脚注：对 AlGaN/GaN 该质量被高估约 $50\%$，
        影响落在千分位量级，但这是一处已知的不一致。

  \item \textbf{$\gamma_0$ 是一个拟合常数。} $2.123\times10^{-12}$
        \U{eV\cdot m^{4/3}} 取自文献对 GaN 的取值
        \cite{ambacher2000}，它把三角阱的电场标度关系吸收成一个数；
        换成别的材料体系必须重新标定。

  \item \textbf{不产生势剖面与波函数。} artifact 中没有
        \akey{profiles}、没有亚带能量列表、没有束缚态数目 ——
        下游若需要这些，只能换用 \crefchap{ch:sp1d}。

  \item \textbf{材料常数与屏蔽因子的不确定度原样继承。} 见
        \crefchap{ch:polarization} 的边界一节；本模型不引入新的
        不确定度，但也不消解任何已有的。

  \item \textbf{绝对数值未经标定。} 与全流水线一致：
        L2（Hall 对比）之前，绝对面密度不可信，可用的是趋势与量级。
\end{bounds}

\section{接口契约与小结}
\label{sec:tw-interface}
\index{5artifact@数据通道!ns}

{\small
\begin{longtable}{@{}p{0.06\textwidth} p{0.22\textwidth} p{0.62\textwidth}@{}}
\toprule
\textbf{方向} & \textbf{键} & \textbf{本模块使用/提供的字段}\\
\midrule
\endfirsthead
\toprule
\textbf{方向} & \textbf{键} & \textbf{本模块使用/提供的字段}\\
\midrule
\endhead
\bottomrule
\endfoot
入 & \akey{polarization} &
  \akey{barrier_thickness_nm}, \akey{eps_r_barrier}, \akey{m_eff_barrier},
  \akey{sigma_eff_C_m2}, \akey{x_al}, \akey{band_offset_eV}\\
入 & \akey{gate}（可选） & \akey{barrier_eV}（\akey{barrier_source="deck"} 时）\\
出 & \akey{quantum} &
  契约键 \akey{ns_m2}, \akey{ns_cm2}, \akey{ns_upper_bound_cm2},
  \akey{e0_eV}, \akey{ef_eV}, \akey{phi_b_eV}, \akey{phi_b_source},
  \akey{delta_ec_eV}；另加 \akey{dos_coefficient_eV_m2}、
  \akey{capacitance_factor_per_V_m2}、\akey{occupancy_fraction}、
  \akey{iterations}, \akey{converged}, \akey{residual_m2}\\
出 & \akey{ns} & 标量面密度（\U{cm^{-2}}）\\
\end{longtable}
}

\noindent
与 \crefchap{ch:sp1d} 相比，本模块多交出两个\emph{模型系数}
（\akey{dos_coefficient_eV_m2} 与 \akey{capacitance_factor_per_V_m2}）。
它们不是契约键，而是诊断量：把 \crefeq{eq:charge-control} 与 \crefeq{eq:tw-ef}
中的两个组合系数显式暴露出来，使读者可以在不读代码的情况下复算整个不动点映射。

\subsection*{小结}

本章的结论可以概括为三点。

其一，\textbf{"解析近似"与"自洽求解"是两个模型，不是精度等级。}
二者输出同样的 artifact，因此\emph{必须在契约里可区分}
（\akey{model} 字段与模块版本），也必须在文档里可区分（本章与
\crefchap{ch:sp1d}）。

其二，\textbf{它的偏差形态是可解释的，因而是有用的。}
$+4.95\%$ 的面密度偏差来自势阱深度与统计抬升的相互抵消，
而 $-33\%$ 与 $-39.7\%$ 的能级偏差来自势形的错误。
理解这一点，比记住一个"误差 $5\%$"的结论更有价值。

其三，\textbf{它的历史角色是基线。} 本模块存在的最大意义，
是被 \crefchap{ch:sp1d} 在不改动任何下游模块的前提下替换掉 ——
这是本项目对"模块即插件"的第一次实战验收。

\cleardoublepage
    % 第 9 章   physics.quantum.analytic_tw   v0.1.0
% module_id: physics.quantum.sp1d
% module_version: 1.1.1
% plugin_dir: plugins/quantum_sp1d
% source_files: plugins/quantum_sp1d/__init__.py, plugins/quantum_sp1d/solver.py
% tests: tests/test_quantum_sp1d.py
%
% 本文件是"模块章模板"的样板：后续 8 章必须复刻其节次骨架
%   x.1 模块契约 / x.2 控制方程与物理假设 / x.3 离散化与数值算法 /
%   x.4 耦合求解（仅耦合型模块有此节）/ x.5 验证与基准 /
%   x.6 误差来源与适用边界 / x.7 接口契约与小结
% 任何偏离都视为模板违规。

\chapter{自洽 Schrödinger--Poisson 求解器}
\label{ch:sp1d}

\begin{quote}
\small
\noindent\textbf{模块标识}\quad\pid{physics.quantum.sp1d}\qquad
\textbf{版本}\quad 1.1.1\qquad
\textbf{类别}\quad physics
\end{quote}

\noindent
本章记述二维电子气\index{1concept@概念!2DEG@二维电子气}电荷状态的\emph{自洽}求解：
联立求解静电 Poisson 方程与量子力学 Schrödinger 方程，得到亚带能级、波函数、
占据数以及二维面密度。该模块是 \pid{physics.quantum.analytic_tw}
（\crefchap{ch:analytic_tw}，解析三角阱估计）的即插即用替代实现，二者声明同一
\texttt{quantum} 能力、产出同一组 artifact 键，因此下游模块无需任何改动
\cite{ambacher2000,vurgaftman2003}。

\section{模块契约}
\label{sec:sp1d-contract}\index{4module@模块!physics.quantum.sp1d}\index{6param@参数!algorithm}\index{1concept@概念!plugin architecture@插件架构}

\begin{modulecard}{physics.quantum.sp1d}{1.1.1}
  \cardrow{类别}{physics（物理模型，流水线阶段）}
  \cardrow{能力}{quantum, quantum.sp, quantum.schrodinger\_poisson}
  \cardrow{依赖 artifact}{polarization, materials}
  \cardrow{产出 artifact}{quantum, ns}
  \cardrow{验证状态}{partially-verified}
  \cardrow{源码}{\fpath{plugins/quantum_sp1d/}（多文件插件包）}
  \cardrow{测试}{\fpath{tests/test_quantum_sp1d.py}}
  \cardsep
  \cardrow{算法参数}{algorithm $\in$ \{\akey{newton}, \akey{anderson}\}，默认 \akey{newton}}
  \cardrow{协作者}{经 Deck 的 \akey{nonlinear} 槽解析 \pid{nonlinear.newton}（\crefchap{ch:newton}）}
\end{modulecard}

本模块声明为 \texttt{partially-verified}：其\emph{离散格式}与\emph{导数}均有闭式解
或收敛阶证据支撑（见 \S\ref{sec:sp1d-verification}），但其\emph{绝对数值}依赖
材料常数与屏蔽因子，二者均未经 L2 基准标定（见 \S\ref{sec:sp1d-limits}）。
这一区分贯穿全章：格式正确性可断言，器件数值不可断言。

\section{控制方程与物理假设}
\label{sec:sp1d-model}\index{1concept@概念!conduction band edge@导带边}

\subsection{几何与坐标约定}
\label{sec:sp1d-geometry}\index{1concept@概念!coordinate convention@坐标约定}

求解域为沿外延方向的一维切片，坐标 $z$ 自金属/势垒界面指向体内为正：

\begin{center}
\small
\begin{tabular}{@{}l l l@{}}
\toprule
\textbf{位置} & \textbf{物理含义} & \textbf{边界条件}\\
\midrule
$z=0$      & 金属/AlGaN 界面（肖特基） & Dirichlet：$\Ec(0)=q(\phib-V_g)$\\
$0<z<d$    & AlGaN 势垒（非故意掺杂）  & $\epsr=\epsr^{b}$，$\meff=\meff_{b}$\\
$z=d$      & 异质界面：极化面电荷 + 带阶台阶 & \emph{落在网格面上}（见 \S\ref{sec:sp1d-face}）\\
$d<z<L$    & GaN 沟道与缓冲             & $\epsr=\epsr^{c}$，$\meff=\meff_{c}$\\
$z=L$      & 背面（绝缘缓冲）           & Neumann：零电场\\
\bottomrule
\end{tabular}
\end{center}

竖直方向被解析为量子化方向，横向（栅长方向）不进模型。这是一维近似，
其代价在 \S\ref{sec:sp1d-limits} 中说明。

\subsection{静电方程}
\label{sec:sp1d-poisson}\index{1concept@概念!Poisson equation@Poisson 方程}\index{3method@方法!finite volume@有限体积}

未知量取\emph{导带底} $\Ec(z)$ 而非静电势 $\varphi$。二者相差一个符号
（$\Ec=-q\varphi+\mathrm{const}$），但以 $\Ec$ 为未知量使量子子问题所需的
势剖面\emph{直接可得}，无需在薛定谔求解前再做一次符号转换。定义算子

\begin{equation}
\mathcal{L}[\Ec] \;\equiv\; -\der{}{z}\!\left(\epsr(z)\,\epsO\,\der{\Ec}{z}\right),
\label{eq:poisson-operator}
\end{equation}

则静电方程为

\begin{equation}
\boxed{\;\mathcal{L}[\Ec](z) \;=\; -q\,n(z) \;+\; \sigeff\,\delta(z-d)\;}
\label{eq:poisson}
\end{equation}

其中 $n(z)$ 为体电子密度，$\sigeff$ 为屏蔽后的有效极化面电荷。
\crefeq{eq:poisson} 的符号结构值得强调：源项 $-qn$ 使 $\Ec$ 在电子富集处
\emph{上凸}，即电荷把势阱底部\emph{抹平} —— 这正是自洽求解中屏蔽效应的体现；
界面处 $+\sigeff\delta$ 使 $\Ec$ \emph{向下尖凹}，该凹陷即二维电子气势阱。
若把这两项的符号取反，势阱将变成势垒，2DEG 完全消失。

\crefeq{eq:poisson} 未包含势垒层内的带阶台阶 $\dEc$。带阶是\emph{材料台阶}
而非\emph{电荷}，若将其并入未知量并对它做差分，会在界面处产生一个虚假的
界面偶极子。因此本模块把 $\dEc$ 单独携带，仅在薛定谔方程需要真实势剖面时加上：

\begin{equation}
\Ec^{\mathrm{full}}(z) \;=\; \Ec(z) \;-\; \dEc\,\Theta(z-d),
\label{eq:offset}
\end{equation}

$\Theta$ 为 Heaviside 阶跃。迭代未知量始终是 $\Ec$，即 \crefeq{eq:poisson} 的
\textbf{静电部分}。

\subsection{量子方程}
\label{sec:sp1d-schrodinger}\index{1concept@概念!Schrödinger equation@Schrödinger 方程}\index{3method@方法!BenDaniel--Duke 界面条件}

包络函数近似下，竖直方向的单电子哈密顿量为 BenDaniel--Duke 形式
\cite{benDaniel1966}：

\begin{equation}
\boxed{\;
-\frac{\hb^{2}}{2}\,\der{}{z}\!\left(\frac{1}{\meff(z)}\,\der{\psi_i}{z}\right)
\;+\; \Ec^{\mathrm{full}}(z)\,\psi_i(z)
\;=\; \Ei\,\psi_i(z)
\;}
\label{eq:schrodinger}
\end{equation}

与常见的常质量形式不同，\crefeq{eq:schrodinger} 保留 $1/\meff$ 在微分算子内部。
这是异质界面的正确处理：它给出的界面条件是 $\psi$ 连续、
$(1/\meff)\,\dd\psi/\dd z$ 连续 \cite{benDaniel1966}，而常质量形式会
错误地要求 $\dd\psi/\dd z$ 连续。在 AlGaN/GaN 界面上 $\meff$ 由 $0.30\,\mO$
变到 $0.20\,\mO$，差异不可忽略。

在有效质量近似下，导带被假定为抛物型，只解析 $\Gamma$ 谷；谷简并度 $g_v$
仅作为态密度的倍数进入，不产生额外的能谷结构。

\subsection{占据与面密度}
\label{sec:sp1d-occupation}\index{2qty@物理量!sheet density@面密度}\index{2qty@物理量!Fermi level@费米能级}\index{1concept@概念!subband@亚带}

平衡态下费米能级取为能量零点，$\Ef=0$。第 $i$ 个亚带的二维占据密度由
Fermi--Dirac 统计给出：

\begin{equation}
n_i \;=\; A_i\,\ln\!\left(1+\ee^{(\Ef-\Ei)/\kT}\right),
\qquad
A_i \;=\; \frac{g_v\,\meff_i\,\kT}{\pi\hb^{2}},
\label{eq:occupation}
\end{equation}

其中 $\meff_i$ 为该亚带的态密度有效质量，其定义见 \crefeq{eq:subband-mass}。
总面密度与体密度为

\begin{equation}
\ns \;=\; \sum_i n_i,
\qquad
n(z) \;=\; \sum_i n_i\,\abs{\psi_i(z)}^{2}.
\label{eq:ns}
\end{equation}

束缚态满足 $\Ei<0$（即位于费米能级以下），\crefeq{eq:occupation} 中的
$\ln(1+\ee^{\cdot})$ 形式在 $\Ei$ 远低于 $\Ef$ 时自动退化为线性（简并）区，
无需分段处理；实现中使用 \texttt{logaddexp} 以保证 $\Ei\gg\Ef$ 时不溢出。

由 \crefeq{eq:poisson}，$\sigeff$ 对应的面密度上界为

\begin{equation}
\nsu \;=\; \frac{\sigeff}{q}.
\label{eq:ns-upper}
\end{equation}

$\nsu$ 是物理上界：超过它意味着势垒层内的电场反向，2DEG 不可能存在。
本模块默认\emph{不}施加该钳制（见 \S\ref{sec:sp1d-coupling}）。

\subsection{边界条件与能量零点}
\label{sec:sp1d-bc}\index{1concept@概念!Schottky contact@肖特基接触}

金属/势垒界面为肖特基接触，导带底被钉在

\begin{equation}
\Ec(0) \;=\; q\,(\phib - V_g),
\label{eq:bc-front}
\end{equation}

数值上可用伏特表示。背面 $z=L$ 取零电场（绝缘缓冲）：

\begin{equation}
\der{\Ec}{z}\bigg|_{z=L} \;=\; 0.
\label{eq:bc-back}
\end{equation}

界面的带阶台阶按 \crefeq{eq:offset} 施加。此处有一个易错点：若把 $\dEc$ 的符号
取反，沟道侧导带会被\emph{抬高}而非压低，势阱随之消失。本模块的实现
以势垒侧为参考，沟道侧低 $\dEc$。

\section{离散化}
\label{sec:sp1d-discretisation}\index{3method@方法!discretisation@离散化}

\subsection{面基界面：异质界面是一个"面"，不是一个节点}
\label{sec:sp1d-face}\index{3method@方法!face-based interface@面基界面}

本模块最关键的离散化决策是：\textbf{让异质界面精确落在网格单元之间的面上，
而不是落在某个节点上}。网格节点位于 $d-h/2$ 与 $d+h/2$ 两处，由此

\begin{itemize}[leftmargin=1.7em,itemsep=0.25em,topsep=0.3em]
  \item 导带带阶台阶 $\dEc$、
  \item 介电常数跃变 $\epsr^{b}\to\epsr^{c}$、
  \item 有效质量跃变 $\meff_{b}\to\meff_{c}$、
  \item 极化面电荷 $\sigeff$
\end{itemize}

\noindent
四件事\emph{同时}落在单元边界上。若改为把界面压在节点上，这四件事都会被
以 $O(h)$ 的分辨率抹平，\emph{并且直接吃掉耦合求解的一阶收敛精度} ——
实测结果是从二阶（$p\approx2.08$）退化到一阶（$p\approx1.0$），
详见 \S\ref{sec:sp1d-verification}。因此网格构造函数在必要时会主动拒绝
不合法的输入：当 $d\le h_{\text{fine}}/2$ 时抛出异常，因为此时无法把界面
放到面上。

该决策的另一个后果是：每个单元都完全位于\emph{单一材料}之内，于是中点求积
是二阶的；若界面压在节点上，界面两侧的单元会跨材料，中点求积降为一阶。

\subsection{非均匀网格上的有限体积 Poisson 离散}
\label{sec:sp1d-fv}\index{3method@方法!control volume@控制体}

设节点 $z_1<\dots<z_N$，面间距 $h_i=z_{i+1}-z_i$，单元控制体宽度取

\begin{equation}
w_i \;=\;
\begin{cases}
\tfrac12 h_1, & i=1,\\[2pt]
\tfrac12 (h_{i-1}+h_i), & 1<i<N,\\[2pt]
\tfrac12 h_{N-1}, & i=N,
\end{cases}
\label{eq:cv-width}
\end{equation}

面上的介电系数取两侧的\emph{长度加权串联平均}

\begin{equation}
c_i \;=\; \frac{\epsface_i\,\epsO}{h_i},
\qquad
\epsface_i \;=\; \frac{2}{\dfrac{1}{\epsri{i}}+\dfrac{1}{\epsri{i+1}}}.
\label{eq:face-coeff}
\end{equation}

当材料边界恰好落在该面上时，\crefeq{eq:face-coeff} 正是两个半单元电容的
精确串联组合；在材料均匀处它退化为普通平均。算子 $\mat{L}$ 的第 $i$ 行为

\begin{equation}
(\mat{L}\Ec)_i \;=\;
-c_{i-1}\Eci{i-1} \;+\; (c_{i-1}+c_i)\Eci{i} \;-\; c_i\Eci{i+1},
\qquad 1<i<N,
\label{eq:L-row}
\end{equation}

首行与末行分别为 Dirichlet 与零电场条件：

\begin{equation}
(\mat{L}\Ec)_1 = \Eci{1},
\qquad
(\mat{L}\Ec)_N = -c_{N-1}\Eci{N-1} + c_{N-1}\Eci{N} .
\label{eq:L-bc}
\end{equation}

对分片线性解，该格式是\emph{精确的}（见 \S\ref{sec:sp1d-verification} 中
双介质平行板基准）。

\subsection{界面面电荷的精确消去}
\label{sec:sp1d-sheet}\index{1concept@概念!sheet charge@面电荷}\index{3method@方法!series permittivity average@介电常数串联平均}

面电荷落在面上，因此它不进入 $\mat{L}$，而是被\emph{消去}进右端常数项。
记界面所在面的两侧为节点 $k$（势垒侧）与 $k+1$（沟道侧），界面电位为 $E_I$，
则两个单侧通量为

\begin{equation}
G^{-} = \epsr^{b}\epsO\,\frac{E_I-\Ec^{(k)}}{h/2},
\qquad
G^{+} = \epsr^{c}\epsO\,\frac{\Ec^{(k+1)}-E_I}{h/2},
\qquad
G^{+}-G^{-} = \sigeff .
\label{eq:sheet-flux}
\end{equation}

由 \crefeq{eq:sheet-flux} 解出 $E_I$ 并回代，可验证线性算子 $\mat{L}$ 的形式
\emph{完全不变}，面电荷则被分配为两个节点的源项：

\begin{equation}
b_k \mathrel{+}= \sigeff\,\frac{\epsr^{b}}{\epsr^{b}+\epsr^{c}},
\qquad
b_{k+1} \mathrel{+}= \sigeff\,\frac{\epsr^{c}}{\epsr^{b}+\epsr^{c}} .
\label{eq:sheet-split}
\end{equation}

两个权重之和恒为 $\sigeff$，这是格式自洽的必要条件。常见的"50/50 对分"
只是 $\epsr^{b}=\epsr^{c}$ 的极限情形，在 AlGaN/GaN 上并不成立
（$\epsr^{b}$ 随 Al 组分在 $\sim8.5$–$9.5$ 之间，而 $\epsr^{c}=9.5$），
故本模块采用 \crefeq{eq:sheet-split} 的精确权重。

\subsection{BenDaniel--Duke 算子的对称化}
\label{sec:sp1d-bdd}\index{3method@方法!symmetrisation@对称化}

将 \crefeq{eq:schrodinger} 在控制体上积分，并用 \crefeq{eq:cv-width} 的宽度离散，
得广义本征问题。为把它化为\emph{标准对称}本征问题，作变量替换

\begin{equation}
v_i \;\equiv\; \sqrt{w_i}\,\psi_i ,
\label{eq:symmetrise}
\end{equation}

并定义面系数

\begin{equation}
t_i \;=\; \frac{\hb^{2}}{2\mO}\cdot\frac{\mO}{\bar m_i}\cdot\frac{1}{h_i},
\qquad
\frac{1}{\bar m_i} \;=\; \frac12\!\left(\frac{1}{\meff_i}+\frac{1}{\meff_{i+1}}\right),
\label{eq:t-coeff}
\end{equation}

则对称三对角矩阵 $\mat{A}$ 的两条带为

\begin{equation}
A_{ii} \;=\; \frac{t_{i-1}+t_i}{w_i} \;+\; \Ecifl{i},
\qquad
A_{i,i+1} \;=\; A_{i+1,i} \;=\; -\frac{t_i}{\sqrt{w_i\,w_{i+1}}} .
\label{eq:A-bands}
\end{equation}

于是 \crefeq{eq:schrodinger} 化为 $\mat{A}\mat{v}=\Ei\mat{v}$，可直接调用对称
三对角本征求解器。对称性不是美学选择，而是\emph{后续精确 Jacobian 的前提}：
波函数响应需要对三对角 Hamilton 的正交 resolvent 做一次消元
（\S\ref{sec:sp1d-resolvent}），只有对称形式才能保证该消元的良态与规范选择的
自由度。此外，离散化只在\emph{一个函数}中写定，薛定谔求解与两条响应路径
共用它，因此\emph{导数不可能漂移到它所微分的算子之外}。

界面在有效质量上的处理已由 \crefeq{eq:t-coeff} 的调和平均隐含：
$\bar m_i$ 取两侧的调和平均，等价于 $(1/\meff)\dd\psi/\dd z$ 连续这一
BenDaniel--Duke 界面条件 \cite{benDaniel1966}。

\subsection{子带态密度质量}
\label{sec:sp1d-mass}\index{2qty@物理量!density-of-states mass@态密度有效质量}

\crefeq{eq:occupation} 中的 $\meff_i$ 不是材料的 $\meff(z)$，而是该亚带自身的
态密度有效质量，定义为材料质量在波函数概率密度下的调和平均：

\begin{equation}
\frac{1}{\meff_i} \;=\; \sum_{z} \frac{\abs{\psi_i(z)}^{2}\,w(z)}{\meff(z)}
\qquad\Longleftrightarrow\qquad
\meff_i \;=\; \frac{1}{\sum_z \abs{\psi_i(z)}^{2}w(z)/\meff(z)} .
\label{eq:subband-mass}
\end{equation}

对势垒穿透显著的亚带，\crefeq{eq:subband-mass} 给出的质量明显高于沟道材料值，
这是波函数部分位于势垒层（$\meff_{b}>\meff_{c}$）的直接结果。

\crefeq{eq:subband-mass} 在形式上很不起眼，但它是精确 Jacobian 的一个\emph{必需}
组成项：$\meff_i$ 依赖 $\psi_i$，而 $\psi_i$ 依赖 $\Ec$，故
$\partial n_i/\partial\Ec$ 中含有一条通过 $\meff_i$ 的链。遗漏该项不是
"小近似"，而是一个\emph{缺项}，其误差特征恰好是中心差分最能暴露的那种 ——
\emph{不随差分步长下降的误差地板}，实测约 $1.3\times10^{-5}$
（见 \S\ref{sec:sp1d-jacobian}）。

\section{耦合求解与精确 Jacobian}
\label{sec:sp1d-coupling}\index{3method@方法!self-consistent iteration@自洽迭代}

\subsection{耦合残差与 Newton 系统}
\label{sec:sp1d-residual}\index{3method@方法!Newton--Raphson}\index{6param@参数!tolerance\_eV}

联立 \crefeq{eq:poisson} 与 \crefeq{eq:schrodinger}--\crefeq{eq:ns} 得到一个
关于 $\Ec$ 的非线性代数系统。定义残差

\begin{equation}
\mat{R}(\Ec) \;=\; \mat{L}\Ec \;-\; q\,\mat{W}\,n[\Ec] \;+\; \mat{b},
\label{eq:residual}
\end{equation}

其中 $\mat{W}=\diag(w_i)$（Dirichlet 行置零，因为该行不携带空间电荷），
$\mat{b}$ 汇集了界面面电荷 \crefeq{eq:sheet-split} 与栅极值。
自洽解即 $\mat{R}(\Ec^{\star})=\bm{0}$。

求解策略有二，由参数 \akey{algorithm} 选择：

\begin{itemize}[leftmargin=1.7em,itemsep=0.3em,topsep=0.35em]
  \item \textbf{Newton 法（默认）}：以 $\mat{J}=\partial\mat{R}/\partial\Ec$
        为 Jacobian，交由此 Deck 的 \akey{nonlinear} 槽解析出的求解服务
        （\crefchap{ch:newton}）执行。收敛为二次，实测局部阶 $1.99$。
  \item \textbf{Anderson 加速不动点（回退）}：把 $\Ec\mapsto\mat{L}^{-1}(q\mat{W}n[\Ec]-\mat{b})$
        视作不动点映射，加 Anderson (Pulay) 加速 \cite{anderson1965}。
        不需要协作者模块，用作交叉验证。
\end{itemize}

为何必须用 Newton：裸不动点映射的线性增益约为 $-30$ —— 面密度变化 $1\%$
会使势阱深到让面密度再变化约 $30\%$。纯阻尼迭代要求阻尼因子小于 $0.06$，
并需要数百次迭代；Newton 在个位数迭代内到达同一解。

\subsection{响应算子的三项分解}
\label{sec:sp1d-jacobian}\index{3method@方法!exact Jacobian@精确 Jacobian}\index{3method@方法!perturbation theory@一阶微扰论}

由 \crefeq{eq:residual}，

\begin{equation}
\mat{J} \;=\; \mat{L} \;-\; q\,\mat{W}\,\der{n}{\Ec},
\label{eq:jacobian}
\end{equation}

全部困难集中在 $\dd n/\dd\Ec$。它是\emph{稠密}矩阵（每个节点的势扰动都会
影响每个节点的密度），且由三项构成，三项\emph{全部精确}：

\paragraph{（i）亚带能量响应（主项）。}
势扰动使每个亚带整体平移 $\avg{\psi_i|\delta\Ec|\psi_i}$，从而改变
Fermi--Dirac 占据。由 \crefeq{eq:occupation}，

\begin{equation}
\der{n_i}{\Ei} \;=\; -\frac{A_i}{\kT}\,f_i \;=\; -\beta_i,
\qquad
f_i \;=\; \frac{1}{1+\ee^{\Ei/\kT}},
\label{eq:beta}
\end{equation}

于是（$\avg{\psi_i|\delta\Ec|\psi_i}=\sum_z\abs{\psi_i(z)}^{2}w(z)\,\delta\Ec(z)$）

\begin{equation}
\boxed{\;
\der{n}{\Ec}\bigg|^{(E)}_{z,z'}
\;=\; -\sum_i \beta_i\,\abs{\psi_i(z)}^{2}\,\abs{\psi_i(z')}^{2}\,w(z')
\;}
\label{eq:resp-energy}
\end{equation}

这是一个秩 $n$ 更新（$n$ 为携带的亚带数），承载了主导物理：势阱深度与
电荷之间的负反馈。

\paragraph{（ii）波函数响应。}
势扰动不仅平移能级，还\emph{混合}本征态。一阶微扰论给出

\begin{equation}
\delta\psi_i \;=\; \sum_{j\neq i} \psi_j\,
\frac{\avg{\psi_j|\delta\Ec|\psi_i}}{\Ei-\Ej},
\label{eq:psi-response}
\end{equation}

其对密度的贡献为 $\delta n_i^{(\psi)} = 2 n_i\,\psi_i\,\delta\psi_i$。
\emph{关键之处在于 \crefeq{eq:psi-response} 的求和必须遍历整个谱，不能截断}：
被携带的若干亚带之外，盒态连续谱同样对 $\delta\psi_i$ 有贡献。
旧实现曾把该求和截断到携带的亚带集合内，结果 Jacobian 只能收敛到
局部阶约 $1.5$，残差停在约 $5\times10^{-13}$ 而无法到达机器精度。

\paragraph{（iii）态密度质量响应。}
由 \crefeq{eq:subband-mass}，$\meff_i$ 亦依赖 $\psi_i$，因此

\begin{equation}
\delta\meff_i \;=\; -\,(\meff_i)^{2}\sum_z \frac{2\,\psi_i\,\delta\psi_i\,w(z)}{\meff(z)},
\qquad
\delta n_i^{(m)} \;=\; \frac{n_i}{\meff_i}\,\delta\meff_i .
\label{eq:mass-response}
\end{equation}

该项数量级远小于前两项，却\emph{不可省略}：漏掉它会留下一个不随差分步长
下降的误差地板（实测 $1.3\times10^{-5}$ 相对误差），并使得 Newton 的
末位残差无法压缩到机器精度。

\subsection{正交 resolvent 与秩-1 亏格}
\label{sec:sp1d-resolvent}\index{3method@方法!rank-1 deflation@秩-1 亏格}\index{3method@方法!Sherman--Morrison 公式}\index{3method@方法!pseudo-inverse@伪逆}

\crefeq{eq:psi-response} 的全谱求和在数值上不可直接实施（$O(N^{3})$ 且需要
完整本征分解）。把它写进对称化变量 $v=\sqrt{w}\,\psi$ 后，该求和恰好是
三对角 Hamilton 在携带态之外的正交 resolvent：

\begin{equation}
\sum_{j\neq i}\frac{v_j v_j\T}{\Ei-\Ej}
\;=\; -\sum_{k\neq i}\frac{v_k v_k\T}{E_k-\Ei}
\;=\; -\!\left(\mat{A}-\Ei\mat{I}\right)^{\dagger},
\label{eq:resolvent}
\end{equation}

$\dg$ 记 Moore--Penrose 伪逆。求伪逆的障碍是 $\mat{A}-\Ei\mat{I}$ 奇异
（$v_i$ 张成其零空间）。秩-1 的 Sherman--Morrison 公式
\cite{shermanMorrison1950} 给出

\begin{equation}
\mat{B} = \mat{A}-\Ei\mat{I}+\sigma\,v_i v_i\T
\qquad\Longrightarrow\qquad
\mat{B}^{-1} = \frac{1}{\sigma}v_i v_i\T - \left(\mat{A}-\Ei\mat{I}\right)^{\dagger},
\label{eq:woodbury}
\end{equation}

即伪逆与一个可逆矩阵的逆只差一个秩-1 项。\footnote{秩-$k$ 的推广通常归于
Woodbury 1950 年的普林斯顿技术报告（Memorandum Report 42）；本模块只用到
秩-1 形式，故仅引用 \cite{shermanMorrison1950}。}

本模块\emph{不}选取正则化参数 $\sigma$，而是改用\textbf{固定规范}，从而
完全避开 $\sigma$ 的选择问题：在 $\abs{v_i}$ 最大的节点处摘掉矩阵的行与列
（束缚态零空间向量的每个分量都非零，故任一节点皆可），并把该分量钉为零。
缩减后的三对角块非奇异，其解 $x_g$ 精确满足

\begin{equation}
\left(\mat{A}-\Ei\mat{I}\right)x_g = \bm{b} - (v_i\T\bm{b})\,v_i ,
\label{eq:gauge-system}
\end{equation}

再减去零空间分量即得伪逆作用：

\begin{equation}
\boxed{\;
\left(\mat{A}-\Ei\mat{I}\right)^{\dagger}\bm{b}
\;=\; x_g - (v_i\T x_g)\,v_i
\;}
\label{eq:pseudoinverse}
\end{equation}

\subsection{块右端 Thomas 与复杂度}
\label{sec:sp1d-thomas}\index{3method@方法!Thomas algorithm@Thomas 算法}\index{3method@方法!block right-hand side@块右端}

\crefeq{eq:pseudoinverse} 要求伪逆的\emph{全部 $N$ 列}。关键在于：三对角消元
（Thomas 算法）的消元系数只依赖矩阵本身，与右端无关，因此一次分解即可服务
全部 $N$ 列。本模块把 Thomas 求解器扩展为接受\emph{块右端} $(N\times N)$，
一次前消、$N$ 列同时回代，于是

\begin{equation}
\text{代价} \;=\; O(N^{2})\ \text{每亚带},
\qquad\text{对比全谱本征矢求和的}\ O(N^{3}).
\label{eq:complexity}
\end{equation}

此外，伪逆的右端具有特殊结构：第 $j$ 列为 $\bm{e}_j - v_{i,j}v_i$（规范行置零），
因此无需显式构造单位阵的每一列。

为保证正确性，本模块同时保留一条\emph{独立}实现路径
（\texttt{quantum\_response\_jacobian\_reference}），用 $O(N^{3})$ 的全谱
本征矢求和计算同一个 $\dd n/\dd\Ec$。两条路径共享其余全部装配代码，
\emph{只}在 resolvent 的求法上不同，因此二者的任何分歧都孤立地指向秩-1 亏格
这一处。测试断言二者在舍入误差内一致。

另有一道防御：亏格的前提是 $v_i$ 确实零化 $\mat{A}-\Ei\mat{I}$。若调用者
用手工拼装的过期状态（势剖面与算符不匹配）传入，会得到一个\emph{看起来合理
但毫无意义}的 Jacobian。因此装配前有一道本征对一致性守卫，直接抛出异常
而不是静默继续。

\subsection{热启动与初值}
\label{sec:sp1d-warmstart}\index{3method@方法!warm start@热启动}\index{3method@方法!Fang--Howard 变分形状}\index{3method@方法!bisection@二分法}\index{6param@参数!gamma0\_eV\_m43}

初值只影响收敛速度而不影响收敛结果（有测试证明，见 \S\ref{sec:sp1d-verification}），
但一个接近解的初值对 Newton 的稳健性至关重要。热启动分两步。

\paragraph{第一步：耗尽近似给出面密度估计。}
用二分法求解标量电荷控制关系

\begin{equation}
\begin{aligned}
\ns &= \mathrm{DOS}\cdot\bigl(W(\ns) - E_0(\ns)\bigr),\\[2pt]
W(\ns) &= q\!\left(\sigeff-q\ns\right)\frac{d}{\epsr^{b}\epsO} - q(\phib-\dEc),
\qquad
E_0(\ns)=\gamma_0\,\ns^{2/3},
\end{aligned}
\label{eq:depletion-seed}
\end{equation}

其中三角阱基态项 $E_0=\gamma_0\ns^{2/3}$ 取自 Ambacher 等的电荷控制模型
\cite{ambacher2000}。右端关于 $\ns$ 单调下降，故二分法总能夹住根。
需强调：\crefeq{eq:depletion-seed} 给出的是\emph{解析模型}的答案，
它与自洽解\emph{不同}（同器件下面密度相差 $+4.95\%$，见
\S\ref{sec:sp1d-verification}），此处仅用作起点。

\paragraph{第二步：Fang--Howard 变分形状给出一致剖面。}
以 Fang--Howard 变分波函数形状 \cite{fangHoward1966}

\begin{equation}
\abs{\psi_{\mathrm{FH}}(z)}^{2} \;\propto\; \tfrac12 b^{3}(z-d)^{2}\ee^{-b(z-d)},
\qquad
b \;=\; \left(\frac{33\,\meff_{c}q^{2}\ns}{8\,\epsr^{c}\epsO\hb^{2}}\right)^{1/3},
\label{eq:fang-howard}
\end{equation}

构造初始体密度，做\emph{一次线性} Poisson 求解得到初始 $\Ec$ 剖面。
由于 \crefeq{eq:poisson} 此时已包含界面面电荷，初值本身就带有一个成形的势阱，
Newton 只需修正其形状。

\subsection{Anderson 回退与密度钳制}
\label{sec:sp1d-anderson}\index{3method@方法!Anderson acceleration@Anderson 加速}\index{3method@方法!fixed-point iteration@不动点迭代}\index{6param@参数!density\_clamp}

Anderson 加速器作用在不动点映射 $g$ 上：令残差 $r_k = g(x_k)-x_k$，
以历史增量 $\Delta\mat{X}$、$\Delta\mat{R}$ 求解最小二乘问题
$\gamma = \argmin\norm{\Delta\mat{R}\,\gamma - r_k}_2$，再做

\begin{equation}
x_{k+1} \;=\; x_k + \beta\,r_k - \left(\Delta\mat{X}+\beta\,\Delta\mat{R}\right)\gamma .
\label{eq:anderson}
\end{equation}

当最小二乘解病态（$\norm{\gamma}_\infty$ 超过阈值或非有限）时，清空历史并
退回纯阻尼步。自适应阻尼在残差连续三次变差时把 $\beta$ 减半，下限为
\akey{mixing_min}。

\textbf{密度钳制默认关闭}，这是一个刻意的决定。把 $\ns$ 钳制在
$[0,\nsu]$ 上看似是"物理保险"，但把密度\emph{正好钉在}上界会制造一个
\emph{伪不动点}：两种迭代都可能卡死其上（实测确认）。因此它只作为
最后手段的守卫保留，且\emph{不进入 Jacobian} —— 钳制是不连续的，
对它在边界处求导没有意义。

\section{验证与基准}
\label{sec:sp1d-verification}\index{3method@方法!verification@验证}\index{1concept@概念!convergence order@收敛阶}

本模块的验证策略是分层的：子求解器对\emph{闭式解}，耦合系统对\emph{收敛阶}，
导数对\emph{中心差分}与\emph{独立实现}。所有基准都在
\fpath{tests/test_quantum_sp1d.py} 中以断言形式固化。

\subsection{子求解器对闭式解}

\begin{veriftable}{Poisson 与 Schrödinger 子求解器的闭式解基准}{tab:sp1d-closed-form}
\vgroup{Poisson 算子（\S\ref{sec:sp1d-fv}--\S\ref{sec:sp1d-sheet}）}
\vrow{双介质平行板（无面电荷）}{$<10^{-10}$\U{V}}{$<10^{-10}$\U{V}（分片线性精确）}
\vrow{双介质平行板 $+$ 界面面电荷}{$<10^{-10}$\U{V}}{$<10^{-10}$\U{V}（精确，验证 \crefeq{eq:sheet-split} 的权重）}
\vrow{$\Ec$ 算子 vs.\ $\varphi$ 形式独立离散化}{$<10^{-10}$\U{eV}}{
  $2.9\times10^{-13}$\U{eV}}
\vgroup{Schrödinger 算子（\S\ref{sec:sp1d-bdd}）}
\vrow{无限深方阱 $E_n=n^{2}\pi^{2}\hb^{2}/(2\meff L^{2})$}{二阶收敛}{二阶收敛（相邻误差比 $3$--$5$）}
\vrow{谐振子 $E_n=\hb\omega(n+\tfrac12)$}{$<10^{-3}$\U{eV}}{$<10^{-3}$\U{eV}}
\vrow{三角阱 Airy 零点 $E_n=\abs{\zeta_n}(\hb^{2}F^{2}/2\meff)^{1/3}$}{二阶收敛}{
  $5.5\times10^{-5}\to3.4\times10^{-6}$（加密一档）}
\end{veriftable}

\subsection{导数：精确性而非近似性}

\begin{veriftable}{耦合 Jacobian 的精确性证据}{tab:sp1d-jacobian}
\vrow{三项完整 Jacobian 对中心差分}{相对误差小}{$\sim10^{-11}$（即差分公式本身的噪声水平）}
\vrow{秩-1 亏格 vs.\ 全谱本征矢求和}{$<10^{-9}$}{$4.0\times10^{-13}$}
\vrow{亏格伪逆的定义检验
     $\left(\mat{A}-\Ei\right)\mat{G}=\mat{I}-v_iv_i\T$ 且 $v_i\T\mat{G}=\bm{0}$}{$<10^{-9}$}{通过（逐亚带）}
\vrow{\emph{截断}波函数响应时的误差地板}{---}{
  $1.3\times10^{-5}$，\textbf{不随差分步长下降}（即"缺项"的特征）}
\vnote{最后一行是反例基准：它记录的是\emph{错误}实现的症状，用于说明为何
       \crefeq{eq:mass-response} 与全谱求和不可省略。}
\end{veriftable}

\subsection{耦合系统的收敛阶}

\begin{veriftable}{耦合系统的收敛性与不变量}{tab:sp1d-convergence}
\vrow{\textbf{竖直间距收敛阶}（Richardson 三档）}{$p>1.7$}{\textbf{$p\approx2.08$}
     （面基界面之前为 $\approx1.0$）}
\vrow{\textbf{耦合 Newton 局部收敛阶}}{$>1.9$}{\textbf{$1.99$}
     （三项完整；截断 Jacobian 时仅 $\approx1.5$）}
\vrow{Newton 末位残差}{$<10^{-13}$\U{eV}}{$1.3\times10^{-16}$\U{eV}（机器精度）}
\vrow{Newton vs.\ Anderson 同一答案}{相对差 $<5\times10^{-3}$}{通过（钳制使命中值略偏内）}
\vrow{Newton 迭代数优势}{$\text{it}_{\mathrm{N}}\times5<\text{it}_{\mathrm{A}}$}{
  Newton $10$ 次 vs.\ Anderson $60$ 次，落到同一位数}
\vrow{热启动无关性（$0.35\sigeff/q\to0.95\sigeff/q$）}{相对差 $<5\times10^{-3}$}{$<0.5\%$}
\vrow{求解域深度无关性（$30\to120$\U{nm}）}{相对差 $<10^{-3}$}{$8\times10^{-6}$}
\vrow{竖直加密不变性}{相对差 $<10^{-2}$}{通过}
\vrow{面密度低于极化上界 $\nsu$}{$\ns\le\nsu$}{占据率 $0.83$}
\vrow{温度趋势（$77$ vs.\ $300$\U{K}）}{$0.9\ns^{300}<\ns^{77}\le\ns^{300}$}{通过}
\vrow{负栅压耗尽沟道}{$\ns(-4\,\mathrm{V})<\ns(0\,\mathrm{V})$}{通过}
\end{veriftable}

\noindent
\textbf{域深度无关性}一项值得单独说明：盒态连续谱是有限求解域的产物，
其\emph{单个}波函数并非物理的束缚态。该基准证明它们作为三维连续态的
\emph{求积}是合法的 —— 域深度扩大四倍，面密度只变 $8\times10^{-6}$。

\subsection{与解析模型的对照}

同一个器件（Al$_{0.25}$Ga$_{0.75}$N/GaN，$20$\U{nm} 势垒，
$L_g=150$\U{nm}，屏蔽因子 $0.8$）下，解析三角阱与自洽解的差别具有
明确的\emph{物理形态}：

\begin{veriftable}{自洽解与解析三角阱的对照（同器件）}{tab:sp1d-vs-analytic}
\vrow{量}{解析三角阱}{自洽 SP（差）}
\vrow{$\ns$ (\U{cm^{-2}})}{$7.58287\times10^{12}$}{$7.95848\times10^{12}$（$+4.95\%$）}
\vrow{$\sigeff/q$ 占据率}{$0.7886$}{$0.8276$（$+4.95\%$）}
\vrow{$\Ef$ 距界面带边 (\U{eV})}{$0.461018$}{$0.308840$（$-33.0\%$）}
\vrow{$E_0$ 距界面带边 (\U{eV})}{$0.380340$}{$0.229336$（$-39.7\%$）}
\vrow{$\Ef-E_0$ (\U{eV})}{$0.0806778$}{$0.0795039$（$-1.46\%$）}
\end{veriftable}

\noindent
两个模型对\emph{势}的看法差别很大（$\Ef$ 与 $E_0$ 都掉了约三分之一），
但对\emph{统计}的看法收敛到同一处：决定 $\ns$ 的亚带间距 $\Ef-E_0$
只差 $1.46\%$。这一形态本身就是结论 —— 自洽解给出的势阱更浅，
而其占据由几乎相同的亚带间距支配。

自洽解还给出解析模型根本不具备的量：多亚带阶梯、波函数向势垒的穿透、
由 \crefeq{eq:subband-mass} 决定的真实态密度质量、束缚态数目、
波函数质心与均方根展宽。

\subsection{代价的如实记录}

精确 Jacobian 不是免费的。装配代价比旧的截断版本高约一个数量级：

\begin{veriftable}{精确 Jacobian 的装配代价}{tab:sp1d-cost}\index{1concept@概念!assembly cost@装配代价}
\vrow{网格节点数}{装配（截断 $\to$ 精确）}{整解耗时比}
\vrow{$119$}{$0.6\to5.1$\U{ms}}{$85\to145$\U{ms}（$1.7\times$）}
\vrow{$208$}{$1.3\to16.5$\U{ms}}{$300\to460$\U{ms}（$1.5\times$）}
\vrow{迭代数与线性求解次数}{不变}{不变}
\vnote{计时口径：同机交错计时取 $7$ 次最小值。买到的是一致性
      （\crefchap{ch:transport_dd} 误差预算的前提），\emph{不是}更少的迭代。
      两条后续提速路线见 \S\ref{sec:sp1d-limits}。}
\end{veriftable}

\section{误差来源与适用边界}
\label{sec:sp1d-limits}\index{1concept@概念!validity limits@适用边界}

\limitsheading

\begin{bounds}
  \item \textbf{材料常数未经独立验证。} 极化常数、弹性常数、带隙与有效质量
        转录自文献 \cite{ambacher2000,vurgaftman2003,bernardini1997,dreyer2016}，
        在材料库中标记为 \texttt{unverified}。在通过 Hall 基准（L2）之前，
        \emph{绝对数值不可信}，可断言的是格式与趋势。

  \item \textbf{屏蔽因子没有公认取值。} 参数 \akey{screening}（默认 $0.8$）\index{6param@参数!screening}
        吸收了表面态屏蔽的全部未知物理，是极化链上最大的不确定度来源，
        必须用 Hall 数据标定。

  \item \textbf{有效质量近似，忽略非抛物性。} 导带被假定为抛物型；
        AlGaN/GaN 的强极化使非抛物性不可忽略，尤其是在高面密度下。
        应变对有效质量的修正亦未计入。

  \item \textbf{只有 $\Gamma$ 谷。} 参数 \akey{valley_degeneracy} 只缩放
        态密度，不产生额外的能谷结构。

  \item \textbf{势垒以上的态是有限域盒本征态。} 它们是三维连续态的合法求积
        （域深度依赖 $<10^{-5}$），但\emph{单个}波函数不是物理的束缚态。
        亚带数 \akey{n_subbands} 过少时\index{6param@参数!n\_subbands}，密度将成为势的粗糙泛函，
        拖慢耦合迭代。

  \item \textbf{一维近似。} 横向（短栅效应）、栅边缘场集中、源漏接入区
        均不在模型内。这些属于 \crefchap{ch:transport_dd} 的二维域。

  \item \textbf{Anderson 回退比 Newton 脆弱。} 更慢，且在部分网格上会因
        密度钳制而卡住。默认算法是 Newton。

  \item \textbf{精确 Jacobian 的装配较贵}（见 \creftab{tab:sp1d-cost}）。
        两条正路已明确：(a) 三对角逆的半可分（semiseparable）Green 函数
        两步扫描可把装配降到 $O(N)$ 次向量化运算；(b) 线性步改为
        \emph{矩阵无关}的 Krylov 方法后，响应只需对向量作用
        （每亚带一次三对角求解 $O(N)$），根本不必组装稠密矩阵。

  \item \textbf{无陷阱、无自热、无量热耦合。} 因此电流坍塌、gate-lag、
        drain-lag 一概不在本模块的适用范围。

  \item \textbf{密度钳制默认关闭。} 其理由是数值的而非物理的：钳制会
        制造伪不动点，见 \S\ref{sec:sp1d-anderson}。
\end{bounds}

\section{接口契约与小结}
\label{sec:sp1d-interface}\index{5artifact@数据通道!quantum}\index{5artifact@数据通道!ns}\index{5artifact@数据通道!polarization}\index{1concept@概念!artifact contract@数据通道契约}

\subsection{artifact 契约}

本模块从上下文读取两个 artifact，写回两个。表中列出的字段名与代码中
\texttt{ctx.put}/\texttt{ctx.require} 的实参逐字符一致。

{\small
\begin{longtable}{@{}p{0.06\textwidth} p{0.22\textwidth} p{0.62\textwidth}@{}}
\toprule
\textbf{方向} & \textbf{键} & \textbf{本模块使用/提供的字段}\\
\midrule
\endfirsthead
\toprule
\textbf{方向} & \textbf{键} & \textbf{本模块使用/提供的字段}\\
\midrule
\endhead
\bottomrule
\endfoot
入 & \akey{polarization} &
  \akey{barrier_thickness_nm}, \akey{eps_r_barrier}, \akey{m_eff_barrier},
  \akey{sigma_eff_C_m2}, \akey{x_al}, \akey{band_offset_eV},
  \akey{channel_material}\\
入 & \akey{materials} &
  沟道材料的 \akey{eps_r}, \akey{m_eff}\\
出 & \akey{quantum} &
  契约键 \akey{ns_m2}, \akey{ns_cm2}, \akey{ef_eV}, \akey{e0_eV},
  \akey{phi_b_eV}, \akey{delta_ec_eV}, \akey{converged}, \akey{iterations}；
  另加亚带能级/质量/分密度、束缚态数、质心与展宽、残差历史、剖面等自洽细节\\
出 & \akey{ns} & 标量面密度（\U{cm^{-2}}），供提取器直接消费\\
\end{longtable}
}

契约键中的前七项与 \pid{physics.quantum.analytic_tw} 完全一致，这是
\emph{即插即用}的机械保证：整个替换只需改动 Deck 中的一行模块绑定，
\pid{extract.fom}（\crefchap{ch:fom}）与流水线均无须改动，因为下游只认
artifact 键而不认实现。

\subsection{小结}

本章的核心结论可以概括为三点。

其一，\textbf{异质界面必须作为"面"处理}。把带阶、介电跃变、质量跃变与面电荷
四件事同时放到单元边界上，耦合求解才恢复二阶收敛（$p\approx2.08$）；
压在节点上会同时损失精度与一个收敛阶。

其二，\textbf{Jacobian 的精确性是可达的，但需要把三项都写对}。
亚带能量响应、波函数响应（全谱、经秩-1 亏格降为 $O(N^{2})$）与态密度质量响应
三者缺一不可。付清这笔账之后，Newton 的局部阶回到 $1.99$，
末位残差落到机器精度 $1.3\times10^{-16}$。

其三，\textbf{格式可信与数值可信是两件事}。本模块的离散格式有闭式解与收敛阶
背书，但绝对数值受制于未标定的材料常数与屏蔽因子。这一区分在本文档中
始终显式给出，不作合并陈述。

\cleardoublepage
           % 第 10 章  physics.quantum.sp1d          v1.1.1

% ---- 第四部分 -------------------------------------------------------
\part{输运与提取}
% module_id: transport.dd
% module_version: 0.1.0
% plugin_dir: plugins/transport_dd
% source_files: plugins/transport_dd/{__init__,plugin,physics,assemble,newton}.py
% tests: tests/test_transport_dd.py

\chapter{稳态漂移扩散输运}
\label{ch:transport_dd}

\begin{quote}
\small
\noindent\textbf{模块标识}\quad\pid{transport.dd}\qquad
\textbf{版本}\quad 0.1.0\qquad
\textbf{类别}\quad transport
\end{quote}

\noindent
本章记述电荷\emph{如何流动}。前面的极化与量子模块回答了"电荷从哪来"，
给出一个平衡态的电荷分布；本模块把它扩展成\emph{带偏压}的问题：
在结构化二维网格上联立求解静电 Poisson 方程与电子连续性方程，
扫过栅、漏偏压，产出电流--电压曲线。

这是全流水线中数值上最重的一章。它的未知量数是前几章的
$O(10^{4})$ 倍，Jacobian 是块三对角的，求解过程需要三层递进的
保险机制才不至于发散。相应的，它的验证也最系统：六个级别的基准
（T0--T5）全部先于任何"像不像器件"的检查。

\section{模块契约}
\label{sec:dd-contract}
\index{4module@模块!transport.dd}

\begin{modulecard}{transport.dd}{0.1.0}
  \cardrow{类别}{transport（输运求解，流水线阶段）}
  \cardrow{能力}{transport, transport.dd, current, iv}
  \cardrow{依赖 artifact}{mesh, structure, materials, polarization}
  \cardrow{产出 artifact}{electrostatics, carriers, current, iv}
  \cardrow{验证状态}{partially-verified}
  \cardrow{源码}{\fpath{plugins/transport_dd/}（多文件插件包）}
  \cardrow{测试}{\fpath{tests/test_transport_dd.py}}
  \cardrow{可标定参数}{\akey{temperature_K}\calib、\akey{mu0_cm2_Vs}\calib、
    \akey{vsat_cm_s}\calib、\akey{delta_ec_ratio}\calib}
  \cardrow{协作者}{经 Deck 的 \akey{linear} 槽解析 \pid{linear.blocktridiag}
    （\crefchap{ch:blocktridiag}）}
  \cardrow{模型范围}{稳态、等温、单极（仅电子）}
\end{modulecard}

本模块为何是 \texttt{partially-verified} 而不是 \texttt{unverified}，
以及为何还不是 \texttt{verified}，模块说明书自己给出了判据：
设计稿原本的门槛是"T0--T4 全部通过之前保持 unverified"，它们确实通过了
（均匀电阻闭式解、制造解二阶收敛、端点电流连续性、延拓无关性）；
而阻止它升级为 \texttt{verified} 的，是\emph{数值}所依赖的迁移率未经标定、
\texttt{statistics} 只有 Boltzmann，以及两条设计判据在本离散下
\emph{原理上不可达}（见 \S\ref{sec:dd-limits}）。

\section{控制方程与物理假设}
\label{sec:dd-model}

\subsection{方程组}
\label{sec:dd-equations}
\index{1concept@概念!drift-diffusion@漂移扩散}

在稳态、等温、单极（只考虑电子）的假定下，器件由两个方程描述：

\begin{equation}
\nabla\cdot\bigl(\epsr\epsO\nabla\psi\bigr) \;=\; -q\,\bigl(N_{\text{net}} - n\bigr),
\label{eq:dd-poisson}
\end{equation}

\begin{equation}
\nabla\cdot\Jnn \;=\; 0,
\qquad
\Jnn \;=\; q\,\mumu\,n\,\mathbf{E} \;+\; q\,D_n\nabla n ,
\label{eq:dd-continuity}
\end{equation}

$\psi$ 为静电势，$N_{\text{net}}=N_D^{+}-N_A^{-}$ 为净电离杂质浓度，
$\mathbf{E}=-\nabla\psi$。第二式（电流密度）在本模块中并不直接使用 ——
实际离散的是它的\emph{通量形式}，见 \S\ref{sec:dd-sg}。

"单极"这一限定是有代价的，且代价不小：电子单极模型里没有空穴去耗尽受主。
把缓冲层的 $1\times10^{17}$\U{cm^{-3}} 受主当作固定电离电荷，会在缓冲层内
造出非物理的巨大电场，而 Poisson 解隐含的电子密度又会让"半绝缘"缓冲层
\emph{变成导体} —— 实测表现为端点电流比沟道电流大一个数量级，
电流直接漏进衬底。因此本模块采用单极 HEMT 的标准抽象：
\textbf{缓冲层 = 中性介质 + 载流子钉零 + 对连续性方程的零通量壁}。

\subsection{为什么未知量取 $(\psi,n)$ 而不是 $(\psi,\varphi_n)$}
\label{sec:dd-unknowns}
\index{1concept@概念!unknown choice@未知量选择}

教科书通常以 $(\psi,\varphi_n)$（$\varphi_n$ 为准费米势）为未知量，
因为此时
$n=N_C\exp\!\bigl(q(\psi-\varphi_n)/kT\bigr)$ 使密度随 $\varphi_n$ 指数变化，
数值动态范围被压缩。本模块反其道而行，以\emph{密度本身}为未知量。理由在
\S\ref{sec:dd-sg} 的偏导结构里，此处先给出结论：

宽禁带势垒层与耗尽缓冲层中，密度会下溢到零。若以 $\varphi_n$ 为未知量，
连续性方程在这些区域会变成\emph{恒等于零}的空方程 —— 方程没坏，
但它不含任何信息，Newton 系统在这些行上是奇异的。而以 $n$ 为未知量时，
Scharfetter--Gummel 通量对密度的偏导\emph{与密度无关}
（\crefeq{eq:dd-flux-partials}），连续性方程在任何密度下都保持良态。

代价是未知量的动态范围极大，这由行缩放（\S\ref{sec:dd-scaling}）处理。

\subsection{统计}
\label{sec:dd-statistics}

载流子密度取 Boltzmann 近似：

\begin{equation}
n \;=\; \NC\,\exp\!\left(\frac{\psi-\varphi_n}{kT/q}\right),
\qquad
\NC \;=\; 2\left(\frac{2\pi \meff kT}{h^{2}}\right)^{3/2},
\label{eq:dd-boltzmann}
\end{equation}

$kT/q$ 以伏特计，使指数无量纲。指数在 $\pm100$ 处被截断；
这不是数值上的偷懒，而是一个\emph{可微}的钳制：截断之上实现残差
严格是"指数取截断值"的形式，因此 Jacobian 仍是它的精确导数，
Newton 不会因钳制而失去一致性。

参数 \akey{statistics} \textbf{只接受} \akey{boltzmann}，其余取值
显式报错。这是一个刻意留下的缺口：简并修正（Joyce--Dixon 或
Fermi--Dirac 积分）\emph{必须同时修改 Scharfetter--Gummel 通量}
（广义 Einstein 关系），只改密度公式会得到一个静默不一致的格式。
宁可拒绝，不要半搬。

\subsection{边界条件}
\label{sec:dd-boundary}

\begin{center}
\small
\begin{tabular}{@{}p{0.24\textwidth} p{0.66\textwidth}@{}}
\toprule
\textbf{边界} & \textbf{条件}\\
\midrule
源、漏接触 & 对 $\psi$ \emph{与} $n$ 同时施加 Dirichlet；两者随端电压同步平移，
  使接触面不携带伪电流\\
栅极 & 对 $\psi$ 的 Schottky Dirichlet（$v_{0.1.0}$ 无栅电流），
  对 $n$ 不加约束\\
背面 & 接地（$\psi=0$），衬底视为理想接地\\
无栅表面 & 钉在 Schottky 参考电位（工程近似，见 \S\ref{sec:dd-limits}）\\
缓冲层 & 载流子钉零，且对连续性方程是零通量壁\\
\bottomrule
\end{tabular}
\end{center}

"接触同时钉 $\psi$ 与 $n$"这一点值得强调。若只钉 $\psi$，接触面的
SG 通量会因两侧密度不匹配而产生一个虚假的电流，端点电流随即失去守恒性；
本模块的做法是把平衡态解在接触列上的取值\emph{整个}作为储层，
漏极在此基础上叠加 $V_{ds}$，因此接触面天然无电流。

\section{离散化}
\label{sec:dd-discretisation}

\subsection{有限体积五点格式}
\label{sec:dd-fv}
\index{3method@方法!finite volume@有限体积}

两个方程都在\emph{控制体}上积分。以 $V_i$ 记节点 $i$ 的控制体体积，
$A_{ij}$、$h_{ij}$ 记节点 $i$ 与其邻居 $j$ 之间面的面积与间距，则
Poisson 方程为

\begin{equation}
\sum_{j\in\mathcal{N}(i)}
\frac{\epsface_{ij}\epsO\,A_{ij}}{h_{ij}}
\bigl(\psi_i-\psi_j\bigr)
\;=\;
q\bigl(N_{\text{net},i}-n_i\bigr)V_i \;+\; \sigma_{\text{sheet}}\,\delta_{i,\text{int}},
\label{eq:dd-poisson-fv}
\end{equation}

面上的介电常数取两侧的\emph{调和}（串联）平均：

\begin{equation}
\epsface_{ij} \;=\; \frac{2\,\epsri{i}\,\epsri{j}}{\epsri{i}+\epsri{j}}.
\label{eq:dd-harmonic}
\end{equation}

调和平均是介电跃变与节点\emph{之间}时的正确处理；若界面恰在节点上，
它退化为算术平均。**本模块的界面恰在节点上**（\crefchap{ch:mesh}），
因此这一处在界面上被抹平一个格子宽度 —— 这是设计稿明示接受的代价，
其量化后果见 \S\ref{sec:dd-limits}。

连续性方程按通量守恒写：

\begin{equation}
\sum_{j\in\mathcal{N}(i)} J_{ij}\,A_{ij} \;=\; 0,
\label{eq:dd-continuity-fv}
\end{equation}

$J_{ij}$ 为由 Scharfetter--Gummel 公式给出的面电流密度。

\subsection{Scharfetter--Gummel 通量}
\label{sec:dd-sg}
\index{3method@方法!Scharfetter--Gummel 通量}

面电流不取 $\mathbf{E}$ 与 $\nabla n$ 的直接差分离散（那在大 Peclet 数下
会振荡），而取指数拟合形式 \cite{scharfetter1969}。以 $L$、$R$ 记面两侧：

\begin{equation}
\boxed{\;
J_{LR} \;=\; \frac{q\,\mumu\,kT}{h}
\Bigl[\,n_R\,B(\delta) \;-\; n_L\,B(-\delta)\,\Bigr],
\qquad
\delta \;=\; \frac{\Ec^{L}-\Ec^{R}}{kT},
\;}
\label{eq:sg-flux}
\end{equation}

\begin{equation}
B(x) \;=\; \frac{x}{\ee^{x}-1},
\qquad
B(-x) \;=\; B(x) + x .
\label{eq:bernoulli}
\end{equation}

\crefeq{eq:bernoulli} 中的\emph{反射恒等式}是 SG 权重的基本性质，也是本模块
踩过的一个坑：把反射关系写反（写成 $-B(x)+x$）之后，通量仍然
\emph{完全自洽} —— 系统级 Jacobian 对中心差分照样通过 $4\times10^{-10}$ ——
却在"带边沿电流方向上升"的每个面上系统性错误。修正后平衡态残差从 $35$
降到 $4.8\times10^{-6}$（七个数量级）。教训是：
\textbf{一致性不等于正确性}；通量类子程序必须对闭式解单独测试，
不能只靠"系统自洽"。

实现同时使用两个分支的解析导数 $B'(x)$，并且只对\emph{衰减}的指数求值：

\begin{equation}
B'(x) = \frac{\ee^{x}-1-x\ee^{x}}{(\ee^{x}-1)^{2}},
\qquad
B'(-x) = -B'(x)-1 \quad (x>0),
\label{eq:bernoulli-slope}
\end{equation}

$x\to0$ 处退回 Taylor 展开，避开 $0/0$。

\subsection{通量的解析偏导}
\label{sec:dd-flux-partials}

\crefeq{eq:dd-flux-partials} 列出通量对四个自变量的解析偏导，
它们是 Jacobian 的构件，也都被单独对中心差分验证：

\begin{equation}
\begin{aligned}
\pder{J_{LR}}{n_R} &= \frac{q\mumu kT}{h}\,B(\delta), &
\pder{J_{LR}}{n_L} &= -\frac{q\mumu kT}{h}\,B(-\delta),\\[2pt]
\pder{J_{LR}}{\psi_L} &= \frac{q\mumu kT}{h}\cdot\frac{1}{kT}
  \Bigl[n_R B'(\delta) + n_L B'(-\delta)\Bigr], &
\pder{J_{LR}}{\psi_R} &= -\,\pder{J_{LR}}{\psi_L}.
\end{aligned}
\label{eq:dd-flux-partials}
\end{equation}

三条性质值得逐条点明。

\textbf{（i）密度偏导与密度无关。} $\partial J/\partial n$ 只由几何、
迁移率与带边差决定，不含 $n$。这正是 \S\ref{sec:dd-unknowns} 选择 $n$
作未知量的理由，也解释了为什么宽禁带势垒或耗尽缓冲层里的连续性方程
仍然良态。

\textbf{（ii）势偏导反对称。} 把两个电位同时平移一个常数不改变 $\delta$，
也就不改变电流，因此 $\partial J/\partial\psi_L = -\partial J/\partial\psi_R$。
这个不变量是一个强判据：实现里只写了"带 $n_L$ 的那一项"的写法
是 $100\%$ 错误的，而它的密度偏导仍然精确 —— 正是那种"只在参考状态
做一次系统级 Jacobian 检查"会漏掉的错误。

\textbf{（iii）带边包含异质结台阶。} $\delta$ 中的 $\Ec$ 是
$\Ec^0-\psi$（数值上），因此 SG 通量\emph{同时}携带静电势差与
异质结带阶；平衡态下 $\delta$ 处处为零，通量恒等于零，
这是 T1 基准断言的不变量。

\subsection{场依赖迁移率}
\label{sec:dd-mobility}
\index{3method@方法!Caughey--Thomas 迁移率}

低场迁移率之外，本模块默认启用 Caughey--Thomas 速度饱和模型
\cite{caugheyThomas1967}：

\begin{equation}
\mumu(E) \;=\; \mumu_0\left[1+\left(\frac{\mumu_0 E}{v_{\text{sat}}}\right)^{\beta}\right]^{-1/\beta},
\label{eq:dd-mobility}
\end{equation}

低场下趋于 $\mumu_0$，高场下趋于 $v_{\text{sat}}/E$。
面电场取 $E=\abs{\Delta\Ec}/h$，即\emph{该面自己的}场。

关键在于它的导数必须进 Jacobian：SG 前因子带 $\mumu$，而面电场依赖带边，
所以 \crefeq{eq:dd-flux-partials} 的势偏导要多出一项

\begin{equation}
\pder{J}{\psi_L}\bigg|_{\mumu} \;=\;
\frac{q\,kT}{h}\Bigl[n_R B(\delta)-n_L B(-\delta)\Bigr]\pder{\mumu}{E}\pder{E}{\psi_L},
\qquad
\pder{E}{\psi_L} = -\frac{\operatorname{sgn}(\delta)}{h},
\label{eq:dd-mobility-chain}
\end{equation}

$\dd\mumu/\dd E$ 有闭式（在 $E=0$ 处保持有限，$\beta\ge1$）。
这一项是\emph{反对称}的，因此性质（ii）在启用场依赖后依然成立 ——
测试对此有单独断言，且**分别在启用与关闭场依赖两种情形下**做了
通量偏导的中心差分校验。

\subsection{块三对角结构与未知量排序}
\label{sec:dd-block}
\index{1concept@概念!block tridiagonal@块三对角}

每个节点携带两个未知量，按 $(\psi,\varphi_n)$ 交错存储；
竖直行 $j$ 展平后是 $2n_x$ 个未知量。张量积网格上，行 $j$ 只与
$j\pm1$ 耦合，因此 Jacobian 是\emph{块三对角}的，块大小 $m=2n_x$：

\begin{equation}
\mat{J} \;=\; \operatorname{tridiag}\bigl(\mat{L}_j,\ \mat{D}_j,\ \mat{U}_j\bigr),
\qquad \mat{L}_j,\mat{D}_j,\mat{U}_j \in \mathbb{R}^{m\times m}.
\label{eq:dd-block-structure}
\end{equation}

对演示网格（$n_x=41$，$n_z=1621$），未知量总数 $2n_xn_z = 132\,922$ ——
稠密分解不可行，而块三对角直接解的代价是 $O(n_zm^{3})=O(n_z(2n_x)^3)$，
这正是 \crefchap{ch:blocktridiag} 存在的理由。线性步经 Deck 的
\akey{linear} 槽解析，因此本模块并不知道、也不依赖具体用什么分解。

\subsection{行缩放：量纲化而非对角均衡}
\label{sec:dd-scaling}
\index{3method@方法!row scaling@行缩放}

两个方程块的单位不同（每单位宽度的电荷 vs.\ 每单位宽度的电流），
量级相差许多个数量级，直接取残差范数会失去意义。本模块对每个
非 Dirichlet 行除以\emph{固定的物理参考量}：

\begin{equation}
\text{Poisson 行} \div C_{\text{local}},
\qquad
\text{连续性行} \div I_{\text{ref}},
\label{eq:dd-row-scaling}
\end{equation}

前者把 Poisson 残差变成\emph{伏特}，后者把连续性残差变成\emph{相对电流}。
这是一次左乘正对角矩阵，不改变解，但使残差范数成为两个块共同的、
有物理意义的收敛度量。

有一条被\emph{否决}的做法值得记录：用 Jacobian 对角做行均衡
（即自适应缩放）。它会让宽禁带行与耗尽行获得巨大权重，
残差范数随即被相消噪声主导，线搜索于是接受"物理残差变差十二个数量级"的步。
实测把平衡态残差从 $4.2$ 改到 $1.3\times10^{-3}$ 才回到正轨。
未知量的\emph{列}缩放也被试过并否决，理由写在代码里：
在 $(\psi,n)$ 表述下矩阵元素与未知量本身无关
（$\dd J/\dd n$ 不含 $n$，Poisson 的 Laplacian 不含 $\psi$），
列缩放无法改善条件数，盲目施加反而会让块消元的枢轴溢出。

\section{求解算法}
\label{sec:dd-solver}

\subsection{三层保险}
\label{sec:dd-three-layers}
\index{3method@方法!pseudo-transient continuation@伪瞬态延拓}

方案把"Newton 收敛困难"列为最高风险，落地为三层递进的保险：

\begin{enumerate}[leftmargin=1.9em,itemsep=0.4em,topsep=0.4em]
  \item \textbf{偏置扫描延拓}：漏压（外层）与栅压（内层）逐点推进，
        前一点的收敛解作为下一点的初值；
  \item \textbf{阻尼线搜索}：残差不降则步长减半（Armijo 式接受），
        最多 \akey{damping_max_halvings} 次；
  \item \textbf{伪瞬态延拓（PTC）}：当线搜索连续失败时，在\emph{迭代矩阵}
        上注入物理时间项
        \begin{equation}
        R^{\text{cont}}_i \;\to\; R^{\text{cont}}_i
        + \bigl(n_i-n_i^{\text{prev}}\bigr)\frac{V_i}{\tau},
        \label{eq:dd-ptc}
        \end{equation}
        并让 $\tau$ 每成功一步 $\times10$（上限 $1$\U{s}）、每失败一步
        $\times0.1$；$\tau$ 足够大、Newton 进入收敛域后自动撤除，
        回到纯 Newton。
\end{enumerate}

\crefeq{eq:dd-ptc} 的关键在于\textbf{稳态残差本身不被改动}。时间项
只改变\emph{迭代方向}，Newton 循环始终在监视并判定\emph{物理}残差
$\mat{R}^{\text{phys}}$（实现上由 \texttt{aux['physical']} 单独携带）。
否则收敛判据会把"沿延拓路径走得怎么样"当成"稳态方程解得怎么样"，
两者不是一回事。

\subsection{平衡态锚点}
\label{sec:dd-equilibrium}

零偏压下的解不是"随便一个初值"，而是\emph{精确解}，
因此可以当作整个扫描的锚：先解一次非线性 Poisson
（$\psi$ 与 $n$ 在 Boltzmann 关系下自洽），把结果采用为接触储层参考：

\begin{equation}
\psi^{\text{src}}_j = \psi_{\text{eq}}(z_j, x_{\min}),
\qquad
\psi^{\text{drn}}_j = \psi_{\text{eq}}(z_j, x_{\max}).
\label{eq:dd-reservoirs}
\end{equation}

随后每个偏压点的接触条件由 \crefeq{eq:dd-reservoirs} 加端电压得到，
于是"接触面携带零电流"成为构造性的性质，而不是需要靠收敛去逼近的目标。

\subsection{偏压延拓的电位预移位}
\label{sec:dd-bias-shift}

把上一点的解直接当新初值有一个隐蔽的失败模式：漏极接触的电位跳了
$\Delta V_{ds}$，而内部还停在旧值上 —— 这个\emph{接触局部冲击}
足以让线搜索停滞，即便 Jacobian 是精确的。修法是按接触数据所强加的
同一个线性斜坡预移位内部电位：

\begin{equation}
\psi(x) \;\leftarrow\; \psi(x) + \Delta V_{ds}\cdot
\frac{x-x_{\min}}{x_{\max}-x_{\min}} .
\label{eq:dd-ramp}
\end{equation}

实测效果是把"停滞的解"变成"五步迭代的解"。此外，扫描本身按
\akey{max_step_V} 自动插入中间步（\emph{不}上报，只负责把解带过去）：
从 $0$ 直接跳到 $5$\U{V} 不是延拓，是带错初值的冷启动。

\subsection{端点电流与连续性门禁}
\label{sec:dd-continuity-gate}
\index{3method@方法!continuity gate@连续性门禁}

端点电流由接触面通量求和得到（进入器件为正）：

\begin{equation}
I_s \;=\; \sum_j J_{j,0}\,w_{z,j},
\qquad
I_d \;=\; -\sum_j J_{j,n_x-1}\,w_{z,j}.
\label{eq:dd-terminal-current}
\end{equation}

求和\emph{只遍历真正导通的面}（与残差使用同一掩码）。这一点不是细节：
若把势垒层的载流子面也算进去，端点电流会与求解器刚刚满足的
连续性平衡相差若干个数量级。

守恒性由门禁判定：

\begin{equation}
\frac{\abs{I_s+I_d}}{\max\bigl(\abs{I_d},\abs{I_s}\bigr)} \;\le\; \varepsilon_{\text{cont}}
\quad\text{当}\quad \abs{I_s+I_d} > I_{\text{floor}},
\label{eq:dd-gate}
\end{equation}

绝对地板 $I_{\text{floor}}$ 的存在是必要的：在平衡态附近端点电流是
机器零（实测 $2.7\times10^{-15}$\U{A/m}），此时\emph{相对}误差由浮点噪声
主导，对分母极小的比值设门槛只会制造假警报。低于地板者不设门槛。

\subsection{局部平衡闭合与辅助掩码}
\label{sec:dd-masks}

宽禁带势垒层内没有电流，其连续性方程对 $n$ 是线性齐次方程，
密度因此近乎\emph{自由}（Jacobian 的最小奇异值比最大低十个数量级）。
本模块用一个在物理上精确的替换消除这个自由模：在无电流行上以
\emph{局域平衡}闭合

\begin{equation}
n_i \;=\; \NCeff{i}\,\exp\!\left(\frac{\psi_i}{kT/q}\right)
\qquad\text{（取 } \varphi_n=0\text{）},
\label{eq:dd-local-equilibrium}
\end{equation}

取代连续性残差；行缩放随后把它变成一个\emph{相对密度误差}。
除这一掩码外，网格还携带另外四组掩码（欧姆接触、栅、背面、无栅表面，
以及载流子钉零的缓冲），它们的互斥性与对边界的覆盖由测试断言。

\section{验证与基准}
\label{sec:dd-verification}

验证纪律与 \crefchap{ch:sp1d} 一致：**先对闭式解，再谈不变量，
最后才看 IV 长什么样**。设计稿给出六个级别 T0--T5，
全部在 \fpath{tests/test_transport_dd.py} 中以断言形式固化。
制造解方法（MMS）的使用遵循其标准表述 \cite{roache2002,salari2000}，
验证与确认的总体框架见 \cite{oberkampf2010}。

\subsection{通量的闭式解与偏导}

\begin{veriftable}{Scharfetter--Gummel 通量：闭式解、反射恒等式与偏导}{tab:dd-flux}
\vrow{Bernoulli 函数对各区间}{$10^{-12}$ 相对}{\crefeq{eq:bernoulli} 的闭式解在宽区间内吻合}
\vrow{反射恒等式 $B(-x)=B(x)+x$}{恒等}{通过（含 $x\to0$ 与溢出区）}
\vrow{$B'(x)$ 对中心差分}{相对误差小}{通过（$10^{-12}$ 量级）}
\vrow{四个通量偏导对中心差分}{相对误差小}{通过（\emph{无}场依赖）}
\vrow{四个通量偏导对中心差分（启用 Caughey--Thomas）}{相对误差小}{
  通过（含 \crefeq{eq:dd-mobility-chain} 的 $\dd\mumu/\dd E$ 项）}
\vrow{势偏导反对称}{$\partial J/\partial\psi_L=-\partial J/\partial\psi_R$}{通过}
\vrow{密度偏导与密度无关}{\crefeq{eq:dd-flux-partials} 第一行}{通过（跨 $17$ 个数量级）}
\vrow{异质结台阶上的零通量}{$\abs{J}/\text{scale}<10^{-12}$}{通过}
\vrow{迁移率极限与导数}{低场 $\to\mumu_0$；$\dd\mumu/\dd E$ 对差分}{
  低场偏差 $<10^{-6}$；$E=0$ 处斜率 $<10^{-6}$}
\end{veriftable}

\subsection{T0：均匀电阻闭式解}

把网格退化为直通道、关断极化与肖特基之后，问题有闭式解
$J=qN_D\mumu V/L$。

\begin{veriftable}{T0\ \ 均匀电阻的闭式解（判据：相对误差 $<0.1\%$）}{tab:dd-t0}
\vrow{电流对闭式解 $J=qN_D\mumu V/L$}{$<10^{-3}$}{\textbf{$<10^{-6}$（超出判据三个数量级）}}
\vrow{两端电流一致}{一致}{$<10^{-9}$}
\vrow{电流与偏压严格正比}{比值 $=2$}{$2.0\pm10^{-6}$}
\end{veriftable}

\subsection{T1：平衡态不变量}

\begin{veriftable}{T1\ \ 零偏平衡态（判据：通量 $<10^{-12}$\U{A/m}）}{tab:dd-t1}
\vrow{平衡态端点电流}{$\le10^{-12}$\U{A/m}}{\textbf{$2.7\times10^{-15}$\U{A/m}（机器零）}}
\vrow{横向电流剖面}{$\le10^{-12}$\U{A/m}}{$\sim10^{-14}$\U{A/m}}
\vrow{Poisson 子问题平衡电荷平衡}{残差小}{$<10^{-10}$}
\vrow{带边算子与态密度台阶自洽}{重建 $=$ \akey{nc_eff}}{相对 $10^{-12}$}
\vnote{设计稿 T1 的另一半——与一维自洽 SP 的带边差 $<1$\U{meV}——\textbf{未做}，
      因为层界落在节点上的抹平使该判据在本离散下原理上不可达。
      这本身是一个诚实的结论，见 \S\ref{sec:dd-limits}。}
\end{veriftable}

\subsection{T2：制造解与二阶收敛}

Poisson 子问题用制造解检验：取 $\psi^{\star}=\sin(\pi z/L)$
（与 $x$ 无关，使横向项恒为零，从而\emph{隔离}竖直算子），
把解析源项加进去，加密四档。

\begin{veriftable}{T2\ \ Poisson 子问题的制造解收敛阶}{tab:dd-t2}
\vrow{竖直间距 $25$\U{nm}}{---}{$L_\infty$ 误差 $1.295\times10^{-2}$}
\vrow{$12.5$\U{nm}}{---}{$3.219\times10^{-3}$（观测阶 $2.008$）}
\vrow{$6.25$\U{nm}}{---}{$8.036\times10^{-4}$（观测阶 $2.002$）}
\vrow{$3.125$\U{nm}}{二阶收敛}{$2.008\times10^{-4}$（观测阶 $2.001$）}
\vnote{观测阶由相邻两档误差比的对数给出。三个值都贴在 $2$ 上，
      说明界面以外的格式确实是二阶的 —— 与界面处的一阶抹平并不矛盾：
      制造解在界面两侧都是光滑的。}
\end{veriftable}

\subsection{T3：端点电流连续性门禁}

\begin{veriftable}{T3\ \ 端点电流守恒（判据：相对误差 $<10^{-6}$，全部偏压点）}{tab:dd-t3}
\vrow{平衡态（$0$\U{V}）}{\crefeq{eq:dd-gate}}{低于绝对地板，按设计不设门槛}
\vrow{$V_{ds}\le1$\U{V} 的全部 $10$ 个偏压点}{$<10^{-6}$}{
  \textbf{$1\times10^{-13}$ 至 $1\times10^{-16}$}}
\vrow{耦合 Jacobian 对中心差分（逐行相对）}{相对误差小}{
  $\sim10^{-10}$（沿 Newton 路径）}
\end{veriftable}

\noindent
\textbf{"逐行相对"是判据的关键。} 本模块吃过一次大亏：$( \psi,\varphi_n)$
改 $(\psi,n)$ 时，Jacobian 里残留一行旧代码，把第二个未知量当成准费米势
代进了态密度公式 —— 内部密度与残差相差 \emph{17 个数量级}，
连续性行因此完全失去对 $\psi$ 的耦合，整个耦合矩阵在 Newton 最需要的
方向上近奇异。而当时的差分校验用\emph{列}最大值归一化，
坏掉的项绝对值很小，相对误差看着只有 $10^{-10}$。
改成\emph{逐行相对}判据后立刻暴露。这条教训值得写进任何后续模块：
\textbf{导数校验要在单元级（通量本身）做，且沿 Newton 路径做，
不能只在参考状态做一次。}

\subsection{T4--T5：延拓无关性与趋势}

\begin{veriftable}{T4\ \ 延拓无关性；T5\ \ 趋势与饱和}{tab:dd-t4-t5}
\vgroup{T4\ \ 初值/延拓无关性（判据：相对差 $<5\times10^{-3}$）}
\vrow{同一偏压点经两条不同阶梯到达}{相对差 $<5\times10^{-3}$}{通过（$0.05$ 与 $0.1$\U{V} 阶梯）}
\vgroup{T5\ \ 趋势合理性}
\vrow{栅压变负 $\Rightarrow$ 电流下降}{单调}{通过}
\vrow{迁移率加倍 $\Rightarrow$ 电流加倍}{比值 $=2$}{$2.0\pm0.1$}
\vrow{\textbf{漏电流饱和}}{高场斜率显著小于低场}{斜率比 $1.90\to1.75$}
\vnote{饱和判据必须用\emph{相对}对照。实测 $V_{ds}:0.25\to0.5\to1.0$\U{V} 时
      $I_d:146\to278\to478$\U{A/m}；同器件在\emph{常数}迁移率下斜率比为
      $2.00\to0.92$ —— 栅控器件本身就带约 $8\%$ 的次线性
      （栅-沟道电位改变二维电子气），所以不能指望对照组是欧姆的。
      测试断言的是"场依赖使斜率比明显更小"这一相对判据。}
\end{veriftable}

\subsection{端到端：演示 Deck}

\begin{veriftable}{演示 Deck 的端到端结果（\fpath{decks/demo_hemt_dd.json}）}{tab:dd-demo}
\vrow{运行状态}{success}{通过（$10$ 个偏压点全部收敛）}
\vrow{墙钟时间}{---}{约 $50$\U{s}（\akey{max_step_V} $=0.05$\U{V}）}
\vrow{二维电子气面密度与极化上界}{\crefeq{eq:ns-bound}}{
  比 $\sigeff/q=9.6\times10^{12}$\U{cm^{-2}} 高约 $\mathbf{1.7}$ 倍}
\vrow{网格掩码互斥且覆盖边界}{互斥 + 覆盖}{通过}
\vnote{1.7 倍这一偏差可归因于层界落在节点上的抹平
      （\crefchap{ch:mesh} 与 \S\ref{sec:dd-discretisation}），
      设计稿已预见，此处首次量化。$0.25$\U{V} 阶梯只需约 $9$\U{s}，
      但末端点不收敛 —— 时间是拿收敛性换来的。}
\end{veriftable}

\section{误差来源与适用边界}
\label{sec:dd-limits}
\index{1concept@概念!validity limits@适用边界}

\limitsheading

\begin{bounds}
  \item \textbf{迁移率是占位参数。} \akey{mu0_cm2_Vs}（$1500$\U{cm^2/(V\cdot s)}）
        与 \akey{vsat_cm_s}（$2\times10^{7}$\U{cm/s}）均未经标定。
        L2 基准之前，\emph{绝对数值不可信}；可用的是趋势、
        守恒性与收敛阶。

  \item \textbf{只覆盖到约 $1$\U{V} 的漏压。} 速度饱和已实现且默认启用，
        但 $V_{ds}\ge2$\U{V} 时 Newton 仍不收敛（欠账）。
        这是模型边界而非调参问题。

  \item \textbf{短栅的速度过冲与弹道成分缺席。}
        $L_g=150$\U{nm} 已在漂移扩散假设开始失效的尺度上。
        方案 §15 的对策（流体动力学模型插件 + 经验修正 + 报告标注）
        原样适用。

  \item \textbf{界面被一个单元宽度抹平。} 量化后果即 $1.7$ 倍的面密度偏差。
        设计稿 T1 中"与一维 SP 的带边差 $<1$\U{meV}"这条判据
        在本离散下\emph{原理上不可达} —— 这是诚实结论，不是待办事项。

  \item \textbf{只实现 Boltzmann 统计。} \akey{statistics} 接受其他取值时
        显式报错，因为简并修正必须同时改通量；在 \crefeq{eq:dd-poisson}
        意义上的高密度沟道中，这会使面密度与电流有系统性偏差。

  \item \textbf{缓冲层是抽象而非器件。} "中性介质 + 载流子钉零 + 零通量壁"
        是单极 HEMT 的标准处理，但它不是真实的半绝缘缓冲：
        深能级、俘获、背栅效应一概没有。

  \item \textbf{无栅电流。} $I_g$ 由构造为零，而非由求解得到。
        因此一切与栅泄漏相关的现象（栅漏电、击穿）不在模型内。

  \item \textbf{无陷阱、无自热、无瞬态。} 电流坍塌、gate-lag、
        drain-lag、晶格温升一概没有；RF 大电流点的预测因此偏乐观。

  \item \textbf{无栅表面钉扎是工程近似。} 真实的无栅表面由施主态钉在
        特定能级；本模块把它钉在 Schottky 参考电位上。
        它同时承担"给势垒电位一个锚"的数值职责 ——
        去掉它，势垒电荷近似为零时 Poisson 退化为纯 Laplace，
        势垒电位存在近乎自由的模式。

  \item \textbf{PTC 与阻尼的参数是工程值。} $\tau_0$（$10^{-12}$\U{s}）
        与减半次数不是物理量；换网格尺度可能需要重调。
        测试只锁\emph{行为}，不锁这些参数。
\end{bounds}

\section{接口契约与小结}
\label{sec:dd-interface}
\index{5artifact@数据通道!iv}\index{5artifact@数据通道!current}

{\small
\begin{longtable}{@{}p{0.06\textwidth} p{0.24\textwidth} p{0.60\textwidth}@{}}
\toprule
\textbf{方向} & \textbf{键} & \textbf{本模块使用/提供的字段}\\
\midrule
\endfirsthead
\toprule
\textbf{方向} & \textbf{键} & \textbf{本模块使用/提供的字段}\\
\midrule
\endhead
\bottomrule
\endfoot
入 & \akey{mesh} & \akey{x_nm}, \akey{z_nm}, \akey{nx}, \akey{nz}, \akey{layers[]}（含 \akey{node_slice}）\\
入 & \akey{structure} & \akey{layers[].{role, material, x_al, doping_cm3, doping_kind}},
  \akey{gate.barrier_eV}\\
入 & \akey{materials} & 介电常数、有效质量、带隙（含合金插值）\\
入 & \akey{polarization} & $\sigeff$、势垒厚度、$\Delta E_g$（$\dEc$ 的优先来源）\\
入 & \akey{quantum}（可选） & \akey{potential_eV} 或 \akey{psi_V}，仅作\emph{初值}\\
出 & \akey{electrostatics} & \akey{psi_V}、极值、坐标轴、\akey{gate_reference_V}、偏压\\
出 & \akey{carriers} & \akey{density_m3}、极值、\akey{ns_cm2}、\akey{statistics}、迁移率\\
出 & \akey{current} & 逐点记录（收敛、迭代数、残差、$I_d$/$I_s$/$I_g$、
  连续性误差与门禁结果）、最差连续性误差、失败清单\\
出 & \akey{iv} & \akey{vds_V}/\akey{vgs_V}/\akey{id_A_per_m} 曲线，
  单位为 \U{A/m}（每单位器件宽度）\\
出 & \akey{ns} & 最后偏压点的面密度（\U{cm^{-2}}）\\
\end{longtable}
}

\noindent
契约上有两处刻意的设计：\textbf{场数组只记最后一点}，而\emph{每个}
偏压点的摘要都在 \akey{current} 里 —— 二维场数组随偏压点线性增长，
全存会撑爆上下文，而"哪一点没收敛、为什么"才是诊断真正需要的。
另一处是 \akey{iv} 的逐点完整记录（设计稿只要求曲线，实现多加了
\akey{points}）：门禁失败必须能定位到\emph{具体偏压点}，只给曲线不够用。

\subsection*{小结}

本章的结论可以概括为三点。

其一，\textbf{未知量的选择是一个物理判断，不是数值偏好。}
取 $n$ 而非 $\varphi_n$ 作未知量，使宽禁带势垒与耗尽区的连续性方程
保持良态（$\partial J/\partial n$ 与 $n$ 无关）—— 代价是动态范围极大，
由\emph{固定物理参考量}的行缩放处理，而不是自适应对角均衡。

其二，\textbf{三层保险解耦了三件不同的事}：延拓负责给出好的初值，
线搜索负责保证每一步都下降，PTC 负责在 Newton 暂时失效时改变
\emph{迭代矩阵}而不改变\emph{被解的方程}。稳态残差始终是判据本身，
因此"沿延拓路径的表现"永远不会被误当成"稳态方程的解"。

其三，\textbf{验证的层级必须先于"像不像器件"。}
T0--T5 的顺序不是形式：闭式解先证明\textbf{离散是对的}，
不变量再证明\textbf{守恒是对的}，收敛阶证明\textbf{格式是二阶的}，
只有在这些之后，IV 曲线上的数字才值得被引用 ——
而在迁移率标定之前，那些数字仍然只是\emph{趋势}。

\cleardoublepage
           % 第 11 章  transport.dd                  v0.1.0
% module_id: extract.fom
% module_version: 0.1.0
% plugin_dir: huntian/plugins/extract_fom
% source_files: huntian/plugins/extract_fom/__init__.py
% tests: tests/test_plugins.py

\chapter{电荷控制阈值与优值提取}
\label{ch:fom}

\begin{quote}
\small
\noindent\textbf{模块标识}\quad\pid{extract.fom}\qquad
\textbf{版本}\quad 0.1.0\qquad
\textbf{类别}\quad extract
\end{quote}

\noindent
本章记述流水线的最后一个阶段：把算出的电荷状态打包成一份结果记录，
并补上一个闭式的阈值电压估计。它的主要价值不在于算法难度 ——
本章的公式都是代数的 —— 而在于两件事：一是把\textbf{结果契约}
固定下来，二是把一个容易被误解的\textbf{分工事实}写在明处：
本模块报告的阈值电压与量子模型\emph{无关}。

\section{模块契约}
\label{sec:fom-contract}
\index{4module@模块!extract.fom}

\begin{modulecard}{extract.fom}{0.1.0}
  \cardrow{类别}{extract（结果提取，流水线阶段）}
  \cardrow{能力}{extract, extract.fom}
  \cardrow{依赖 artifact}{ns, quantum, polarization}
  \cardrow{产出 artifact}{results.fom}
  \cardrow{验证状态}{\textbf{unverified}（占位实现，见下）}
  \cardrow{源码}{\fpath{huntian/plugins/extract_fom/}}
  \cardrow{测试}{\fpath{tests/test_plugins.py}（阈值区间、电容区间、压降自洽、趋势）}
  \cardrow{可标定参数}{\akey{vth_correction_V}\calib、\akey{interface_charge_cm2}\calib}
  \cardrow{模型}{全耗尽电荷控制（闭式）、氧化层电容、热电压}
\end{modulecard}

本模块的验证状态是 \texttt{unverified}，且是\textbf{按设计如此}的。
与其余八个模块不同，它没有闭式解可以对、没有收敛阶可以断言 ——
它本身就是一个\emph{模型的闭式解}。因此它的"验证"只能是与实测阈值对比，
而那需要 L2 基准（Id--Vg 实测），目前尚未进行。

模块说明书对此的措辞值得原文保留：\emph{"这不是从真实 Id--Vg 扫描中提取
$V_{\text{th}}$ 的替代品。它是电荷控制极限，作为一个有用的合理性锚点，
以及完整 IV / C--V / S 参数提取链的占位。"}

\section{模型：全耗尽电荷控制}
\label{sec:fom-model}
\index{1concept@概念!threshold voltage@阈值电压}

\subsection{阈值电压}
\label{sec:fom-vth}

在"势垒层完全耗尽、沟道电荷恰好被耗尽掉"这一经典假定下，
栅压需要同时承担三项：肖特基势垒、导带带阶，以及把极化面电荷
$\sigeff$ 在厚度 $d$ 上降掉的电压。于是

\begin{equation}
\boxed{\;
V_{\text{th}} \;=\; \phib \;-\; \dEc \;-\;
\frac{\sigeff\,d}{\epsr\epsO}
\;+\; V_{\text{corr}}
\;}
\label{eq:fom-vth}
\end{equation}

其中各项量纲一致：\crefeq{eq:fom-vth} 中两个能量项以 \U{eV} 计、
其余以 \U{V} 计，二者数值相等，故不出现显式换算因子。
$V_{\text{corr}}$ 是参数 \akey{vth_correction_V}（默认 $0$），
一个吸收"电荷控制估计与实测阈值之差"的经验偏移。

第二项与第三项之间有一个值得点明的分工：$\dEc$ 来自
\emph{材料}（带隙差乘一个比例），而势垒压降来自\emph{静电}
（面电荷乘厚度除以介电常数）。它们都是势垒层的性质。

\subsection{势垒压降与氧化层电容}
\label{sec:fom-cox}
\index{2qty@物理量!oxide capacitance@氧化层电容}

\crefeq{eq:fom-vth} 的第三项单独命名并写进结果记录：

\begin{equation}
V_{\text{drop}} \;=\; \frac{\sigeff\,d}{\epsr\epsO},
\label{eq:fom-drop}
\end{equation}

它是本器件"极化电荷把势垒层压出多少电压"的直接度量 —— 演示器件下
约 $3.76$\U{V}，是三项中最大的一项，也因此决定了阈值是负的。

势垒层的单位面积电容为

\begin{equation}
C_{\text{ox}} \;=\; \frac{\epsr\epsO}{d},
\label{eq:fom-cox}
\end{equation}

它在数值上是 $V_{\text{drop}}$ 与 $\sigeff$ 之间的换算系数
（$V_{\text{drop}}=\sigeff/C_{\text{ox}}$）。本模块把它也写进契约，
因为它是后续跨导与频率特性的基本量：$g_m$ 的第一近似就是
$C_{\text{ox}}W v_{\text{sat}}$ 量级。

热电压同样作为诊断量输出：

\begin{equation}
V_T \;=\; \frac{kT}{q}.
\label{eq:fom-vt}
\end{equation}

\subsection{面密度的界面电荷修正}

真实的界面存在固定的净电荷，它们会抵消一部分二维电子气。
本模块用一个显式参数处理：

\begin{equation}
\ns^{\text{rep}} \;=\; \ns^{\text{quantum}} \;-\; n_{\text{int}},
\label{eq:fom-ns}
\end{equation}

$n_{\text{int}}$ 为参数 \akey{interface_charge_cm2}（默认 $0$）。
同时保留未修正值 \akey{ns_uncorrected_cm2}，使这一处修正始终可被
审计 —— 若只报告修正后的数，"到底改了多少"就无从追查。

\section{一个必须写在明处的分工事实}
\label{sec:fom-division}
\index{1concept@概念!model independence@模型无关性}

本模块报告两个数：$V_{\text{th}}$ 与 $\ns$。它们来自\emph{不同的}模型，
且只有后者经过量子求解：

\begin{center}
\small
\begin{tabular}{@{}p{0.20\textwidth} p{0.34\textwidth} p{0.36\textwidth}@{}}
\toprule
\textbf{量} & \textbf{来自} & \textbf{对量子模型的依赖}\\
\midrule
$V_{\text{th}}$ & \crefeq{eq:fom-vth}，只依赖势垒静电学 &
  \textbf{无}。把量子槽从解析三角阱换成自洽 SP，$V_{\text{th}}$ 一个字都不变\\
$\ns$ & \akey{quantum} artifact 的 \akey{ns_cm2} &
  完全依赖。\crefchap{ch:analytic_tw} 与 \crefchap{ch:sp1d} 给出的值
  相差约 $5\%$\\
\bottomrule
\end{tabular}
\end{center}

这一事实有一个具体的验收后果：在 \crefchap{ch:sp1d} 的即插即用替换中，
下游 \pid{extract.fom} 之所以"一行未改"，正是因为它只认 artifact 键
而不认实现（\crefchap{ch:sp1d}）。而\emph{另一个}后果是使用上的：
读者不应把本模块的 $V_{\text{th}}$ 当作量子模型的输出 ——
它对量子模型的改进\emph{完全不敏感}。把它写上契约，是为了防止
一个看似合理的误解：看到 $V_{\text{th}}$ 与 $\ns$ 在同一个记录里，
就以为它们来自同一次自洽求解。

\section{验证与基准}
\label{sec:fom-verification}

\begin{veriftable}{阈值、电容与压降的自洽性与趋势}{tab:fom-verify}
\vgroup{自洽性}
\vrow{压降与电荷控制关系一致}{\crefeq{eq:fom-drop} 复现报告的压降}{
  精确到 $10^{-6}$（用 artifact 中的 $\sigeff$、$d$、$\epsr$ 复算）}
\vrow{面密度与量子 artifact 一致}{\akey{ns} 等于 \akey{quantum.ns_cm2}}{
  精确到 $10^{-6}$ 相对}
\vgroup{量级合理性}
\vrow{阈值落在物理区间}{$-5<V_{\text{th}}<-1$\U{V}}{演示器件 $-3.43$\U{V}}
\vrow{氧化层电容落在物理区间}{$0.2<C_{\text{ox}}<0.8$\U{\textmu F/cm^2}}{
  演示器件 $0.410$\U{\textmu F/cm^2}}
\vgroup{趋势方向}
\vrow{屏蔽增强 $\Rightarrow$ 阈值更负}{单调减}{
  $s:0.6\to0.9$ 时 $V_{\text{th}}$ 更低}
\vnote{区间判据（而非等值判据）在这里是恰当的：本模块的验证状态是
      \texttt{unverified}，它的常数是占位的，因此只能断言"落在量级
      合理的带内"与"与其他量自洽"。等值断言留给 L2 基准。}
\end{veriftable}

\subsection{演示器件的算例}
\label{sec:fom-worked}

\creftab{tab:fom-worked} 给出演示器件（Al$_{0.25}$Ga$_{0.75}$N，$20$\U{nm} 势垒，
$s=0.8$，$\phib=0.94$\U{eV}）逐项代入的结果。全部中间量都可由
前几章的 artifact 复算。

\begin{veriftable}{演示器件的阈值链逐步计算}{tab:fom-worked}
\vrow{势垒压降 $V_{\text{drop}}$，\crefeq{eq:fom-drop}}{$<5$\U{V}}{$3.762$\U{V}}
\vrow{导带带阶 $\dEc$（比例 $0.7$）}{---}{$0.6055$\U{eV}}
\vrow{肖特基势垒高度 $\phib$（取自 Deck）}{---}{$0.940$\U{eV}}
\vrow{阈值电压 $V_{\text{th}}$，\crefeq{eq:fom-vth}}{$-5<V_{\text{th}}<-1$\U{V}}{
  $-3.428$\U{V}}
\vrow{氧化层电容 $C_{\text{ox}}$，\crefeq{eq:fom-cox}}{$0.2$--$0.8$\U{\textmu F/cm^2}}{
  $0.410$\U{\textmu F/cm^2}}
\vrow{热电压 $V_T$（$300$\U{K}）}{---}{$25.85$\U{mV}}
\vnote{压降是第一项的量级主导者：三项中 $3.762$ 远大于 $0.940$ 与 $0.6055$，
      这正是常关型与常开型 GaN HEMT 之间的分界线所在 ——
      增大 $\sigeff d/(\epsr\epsO)$（更厚的势垒或更强的极化）
      会把阈值推得\emph{更负}，即器件更深地处于常开状态。}
\end{veriftable}

\section{误差来源与适用边界}
\label{sec:fom-limits}
\index{1concept@概念!validity limits@适用边界}

\limitsheading

\begin{bounds}
  \item \textbf{本模块是占位实现，验证状态为 \texttt{unverified}。}
        它的 $V_{\text{th}}$ 是电荷控制极限，不是从 Id--Vg 扫描提取的阈值。
        两者在有短沟效应、界面态或串联电阻的真实器件上会明显不同。

  \item \textbf{全耗尽假定不总成立。} 势垒层未被完全耗尽时
        （例如极厚势垒或栅压过正），\crefeq{eq:fom-vth} 的线性关系失效。

  \item \textbf{$V_{\text{th}}$ 与量子模型无关。} 见
        \S\ref{sec:fom-division}。这不是缺陷，但必须被知道，
        否则会误以为改进量子求解器会改进阈值。

  \item \textbf{无串联电阻、无短沟效应、无陷阱色散。}
        源漏接入区电阻、漏致势垒降低、界面态导致的阈值漂移
        一概不在模型内。

  \item \textbf{$V_{\text{th}}$ 的绝对精度受制于输入常数。}
        $\phib$（源自 Deck 或解析拟合）、$\dEc$（比例 $0.7$ 是模型选择）、
        以及 $\sigeff$（依赖未标定的屏蔽因子）三者共同决定它。
        在 L2 之前，这个数的可信度不高于它的输入。

  \item \textbf{栅电流不参与。} \akey{ig_A_per_m} 由构造为零
        （\crefchap{ch:transport_dd} 无栅泄漏路径），
        因此本模块无从给出与栅相关的优值。

  \item \textbf{没有真正的优值链。} 完整的提取链（IV、C--V、S 参数、
        $g_m$、$f_T$、$f_{\max}$）尚未实现；本模块只固定了
        \akey{results.fom} 这个契约位置。
\end{bounds}

\section{接口契约与小结}
\label{sec:fom-interface}
\index{5artifact@数据通道!results.fom}

{\small
\begin{longtable}{@{}p{0.06\textwidth} p{0.22\textwidth} p{0.62\textwidth}@{}}
\toprule
\textbf{方向} & \textbf{键} & \textbf{本模块使用/提供的字段}\\
\midrule
\endfirsthead
\toprule
\textbf{方向} & \textbf{键} & \textbf{本模块使用/提供的字段}\\
\midrule
\endhead
\bottomrule
\endfoot
入 & \akey{quantum} & \akey{ns_cm2}, \akey{phi_b_eV}, \akey{delta_ec_eV}, \akey{ef_eV}\\
入 & \akey{ns} & 标量面密度（\U{cm^{-2}}），用于与 \akey{quantum} 的一致性核对\\
入 & \akey{polarization} & \akey{barrier_thickness_nm}, \akey{eps_r_barrier},
  \akey{sigma_eff_C_m2}, \akey{sigma_pol_cm2}, \akey{sigma_eff_cm2},
  \akey{screening}, \akey{x_al}\\
出 & \akey{results.fom} &
  \akey{ns_cm2}, \akey{ns_uncorrected_cm2}, \akey{vth_est_V},
  \akey{barrier_drop_V}, \akey{phi_b_eV}, \akey{delta_ec_eV}, \akey{ef_eV},
  \akey{cox_F_m2}, \akey{cox_uF_cm2}, \akey{thermal_voltage_V},
  \akey{temperature_K}, 以及从极化透传的 $\sigpol$、$\sigeff$、屏蔽因子、
  Al 组分与势垒厚度\\
\end{longtable}
}

\noindent
把 $\sigpol$、$\sigeff$、$s$、$x$、$d$ 一并透传进结果记录，是为了让
\emph{最终报告可以自证}：读者拿到一份 \akey{results.fom}，
无须回头翻别的 artifact 就能复算 $V_{\text{th}}$ 与 $C_{\text{ox}}$。
这也正是 \creftab{tab:fom-worked} 的算例得以逐行核对的原因。

\subsection*{小结}

本章的结论可以概括为三点。

其一，\textbf{公式简单不等于结论简单。} \crefeq{eq:fom-vth} 只有三项，
但第一项的绝对精度依赖未标定的屏蔽因子，第二项依赖一个模型比例，
第三项依赖材料常数 —— 因此它的\emph{趋势}可信，\emph{数值}待标定。

其二，\textbf{分工必须在契约里写明。} $V_{\text{th}}$ 与 $\ns$ 出自
不同模型、对量子求解器的敏感度截然不同；把它们放在同一条记录里
并不构成它们同源的证据。本章把这件事写成一张表（\S\ref{sec:fom-division}）
而不是一句脚注。

其三，\textbf{占位实现也要有名字和被断言的性质。}
本模块的 \texttt{unverified} 状态、区间判据、以及与上游 artifact 的
自洽性断言，共同保证了它即便只是占位，也\emph{不会悄悄给出
自相矛盾的结果} —— 这比"假装它已经验证过"要有用得多。

\cleardoublepage
            % 第 12 章  extract.fom                   v0.1.0

% ---- 第五部分 -------------------------------------------------------
%  服务模块用变体模板：§x.1 卡片改记 capabilities 与调用契约
%  （无 requires/produces），§x.2 换为「问题类定义」。
\part{数值服务}
% module_id: nonlinear.newton
% module_version: 1.0.0
% plugin_dir: plugins/nonlinear_newton
% source_files: plugins/nonlinear_newton/{__init__,solver}.py
% tests: tests/test_nonlinear_newton.py
%
% 服务模块变体：§x.1 记 capabilities 与调用契约（无 requires/produces），
% §x.2 换为「问题类定义」。其余槽位与非服务模块一致。

\chapter{通用阻尼 Newton 求解器}
\label{ch:newton}

\begin{quote}
\small
\noindent\textbf{模块标识}\quad\pid{nonlinear.newton}\qquad
\textbf{版本}\quad 1.0.0\qquad
\textbf{类别}\quad solver
\end{quote}

\noindent
本章记述一个\emph{不包含任何物理}的模块：一个通用的稠密
Newton--Raphson 迭代器，连同回溯线搜索与对角正则化。
它的存在本身就是架构主张的一部分 —— \textbf{迭代策略是策略，不是物理}，
因此它不该被写死在任何一个物理模块里，而应当作为一个可替换的协作者
被按需解析出来。

\section{模块契约}
\label{sec:newton-contract}
\index{4module@模块!nonlinear.newton}

\begin{modulecard}{nonlinear.newton}{1.0.0}
  \cardrow{类别}{solver（数值服务，\textbf{不是}流水线阶段）}
  \cardrow{能力}{nonlinear, nonlinear.newton}
  \cardrow{依赖 artifact}{\textbf{无}（服务模块不读上下文数据）}
  \cardrow{产出 artifact}{\textbf{无}（服务模块不写上下文数据）}
  \cardrow{验证状态}{partially-verified}
  \cardrow{源码}{\fpath{plugins/nonlinear_newton/}}
  \cardrow{测试}{\fpath{tests/test_nonlinear_newton.py}}
  \cardrow{可标定参数}{\textbf{无}（求解器参数是数值选择，不是器件物理量）}
  \cardsep
  \cardrow{调用契约}{\texttt{solve(residual, jacobian, x0, options=) -> NewtonResult}；
    \texttt{verify\_jacobian(residual, jacobian, x) -> JacobianCheck}}
  \cardrow{解析方式}{经 Deck 的 \akey{nonlinear} 槽配置，由物理模块通过
    \texttt{ctx.services} 按需取得}
  \cardrow{已知消费者}{\pid{physics.quantum.sp1d}（\crefchap{ch:sp1d}）}
\end{modulecard}

\creftab{tab:newton-vs-service} 对比了服务模块与普通数据模块在契约上的差别。
其中\textbf{产物为空}这一点最容易引起误解，值得明说：它不产生 artifact，
但它\emph{确实}被 Deck 配置、确实参与一次运行、也确实影响结果 ——
只不过它的影响是通过\emph{调用}而非通过\emph{数据通道}传递的。
正因如此，把 \texttt{execute()} 当作流水线阶段调用会被显式拒绝，
错误信息里写明"这是一个服务"。

\begin{veriftable}{服务模块与数据模块的契约差别}{tab:newton-vs-service}
\vrow{流水线阶段}{是}{否（\texttt{execute()} 抛错，信息含 "service"）}
\vrow{读/写 artifact}{是}{否}
\vrow{Deck 中的位置}{\akey{physics} / \akey{modules} 槽}{\akey{modules} 下的服务槽
  （\akey{nonlinear} / \akey{linear} / \akey{backend}）}
\vrow{取得方式}{由 Pipeline 按依赖顺序调度}{
  由 \texttt{ServiceRegistry} 惰性实例化并按配置缓存}
\vrow{参数校验}{同左}{同左（未知参数同样被拒绝）}
\end{veriftable}

\subsection{服务槽机制}

\akey{nonlinear} 槽里的模块\textbf{不是}流水线阶段：它们不读 artifact、
不写 artifact，而是在另一个模块\emph{内部}被调用。它们通过上下文上的
\texttt{ServiceRegistry} 按需解析，参数由 Deck 配置：

\begin{center}
\small
\begin{tabular}{@{}p{0.44\textwidth} p{0.46\textwidth}@{}}
\toprule
\textbf{机制} & \textbf{行为}\\
\midrule
配置 & \texttt{configure(module\_id, version, params)}，注册中心里不存在即报错\\
实例化 & 惰性：首次 \texttt{get()} 时才建实例\\
缓存 & 按 \texttt{(模块标识 $|$ 版本 $|$ 规范化参数)} 作键缓存\\
覆盖 & 调用方传入的 \texttt{params} 覆盖 Deck 槽里的同名参数\\
释放 & \texttt{release\_all()} 统一 \texttt{terminate()}，单个失败不阻断其余\\
\bottomrule
\end{tabular}
\end{center}

缓存键里带\emph{规范化参数}是必要的：同一模块以两组不同参数被请求时，
必须得到两个独立实例，否则后一次调用会静默继承前一次的配置。
而"调用方覆盖"这一条，正是物理模块注入\emph{自己量纲}的容差的方式 ——
例如 \pid{physics.quantum.sp1d} 传入的是以 \U{eV} 计的容差。

\section{问题类定义}
\label{sec:newton-problem}
\index{1concept@概念!service slot@服务槽}

本模块求解的是如下问题类：给定连续的 $\mat{R}:\mathbb{R}^{n}\to\mathbb{R}^{n}$
及其导数 $\mat{J}=\partial\mat{R}/\partial\bm{x}$（稠密可算），求

\begin{equation}
\mat{R}(\bm{x}^{\star}) \;=\; \bm{0}.
\label{eq:newton-problem}
\end{equation}

三个限定词划定了它的能力边界：\textbf{连续}（线搜索依赖残差可比大小）、
\textbf{稠密}（线性步用稠密 LU，见 \S\ref{sec:newton-cost}）、
\textbf{导数可得}（不做割线或拟 Newton）。

\subsection{调用接口}

残差与 Jacobian 以\emph{可调用对象}传入，模块对问题本身一无所知：

\begin{equation}
\mat{R},\ \text{aux} \;=\; \texttt{residual}(\bm{x}),
\qquad
\mat{J} \;=\; \texttt{jacobian}(\bm{x},\ \text{aux}).
\label{eq:newton-callbacks}
\end{equation}

\crefeq{eq:newton-callbacks} 中的 \texttt{aux} 通道是接口设计上的一个刻意安排：
残差求值通常包含一件昂贵的事（对 \crefchap{ch:sp1d} 的自洽 SP 系统而言，
是薛定谔方程的本征分解），而 Jacobian 恰好需要\emph{同一个}中间状态。
把 \texttt{aux} 从残差传给 Jacobian，使这次昂贵计算\emph{每步只做一次}。
若没有这条通道，Jacobian 要么重算一遍，要么把状态藏在模块全局里 ——
前者使代价翻倍，后者使模块不可重入。

\section{算法}
\label{sec:newton-algorithm}

\subsection{Newton 步与阻尼}
\label{sec:newton-step}

核心迭代是标准的：

\begin{equation}
\mat{J}(\bm{x}_k)\,\bm{\delta}_k \;=\; -\mat{R}(\bm{x}_k),
\qquad
\bm{x}_{k+1} \;=\; \bm{x}_k + \lambda_k\,\bm{\delta}_k .
\label{eq:newton-iteration}
\end{equation}

\subsection{回溯线搜索}
\label{sec:newton-line-search}
\index{3method@方法!line search@回溯线搜索}

阻尼因子 $\lambda_k$ 不是固定值，而是由回溯搜索确定的：从 $\lambda=1$ 起，
每次减半，直到满足 Armijo 型充分下降条件

\begin{equation}
\bigl\lVert\mat{R}(\bm{x}_k+\lambda\bm{\delta}_k)\bigr\rVert_\infty
\;\le\;
\bigl(1 - c\,\lambda\bigr)\,
\bigl\lVert\mat{R}(\bm{x}_k)\bigr\rVert_\infty ,
\qquad c = 10^{-4}\ (\akey{armijo}),
\label{eq:newton-armijo}
\end{equation}

线搜索因此是\emph{单调性保证}：接受的每一步都不增大残差范数。
阻尼序列被完整记录在 \texttt{damping\_history} 中，可供诊断 ——
"每一步都是全额步"与"一直在折半"是两种完全不同的收敛故事。

搜索有三个终止口：接受、$\lambda < \akey{min_damping}$（默认 $10^{-4}$）、
或折半次数达到 \akey{max_backtracks}（默认 $12$）。
连续 \akey{stall_limit}（默认 $3$）次搜索失败即中止，
并在消息中明确指出\textbf{"Jacobian 可能与残差不一致"} ——
这是线性搜索失败最常见的根因，把猜测写进错误信息比让人去猜要好。

\subsection{收敛判据：更新范数，不是残差范数}
\label{sec:newton-convergence}

收敛由\emph{更新范数}判定：

\begin{equation}
\bigl\lVert\lambda_k\bm{\delta}_k\bigr\rVert_\infty \;\le\; \texttt{tolerance}
\quad(\text{默认 } 10^{-8}).
\label{eq:newton-convergence}
\end{equation}

为什么不用残差范数？因为对于\emph{行均衡}过的系统
（\crefchap{ch:transport_dd} 与 \crefchap{ch:sp1d} 都是），
残差里最弱的几行会因相消而停在一个地板上，
此时\emph{任何}步都无法把残差范数再降低。若以残差范数判定，
迭代会在地板上空转到上限；以更新范数判定则立刻认定"迭代已停止移动"，
这是正确的语义。

另有一个可选的 \akey{residual_tolerance}：设定后它成为\emph{额外的}
接受条件（允许在残差已足够小时提前接受一个不满足 Armijo 的步），
默认不设。

\subsection{对角正则化（PTC）}
\label{sec:newton-ptc}
\index{3method@方法!pseudo-transient continuation@伪瞬态延拓}

\akey{ptc_shift}（默认 $0$）给出一个常数对角位移：

\begin{equation}
\bigl(\mat{J}(\bm{x}_k) + \sigma\mat{I}\bigr)\bm{\delta}_k \;=\; -\mat{R}(\bm{x}_k).
\label{eq:newton-ptc}
\end{equation}

它与 \crefchap{ch:transport_dd} 的伪瞬态延拓是同一个数学操作
（那里的 $\sigma = V/\tau$），但这里的形式更弱：$\sigma$ 是
\emph{固定的}，不做自适应调度。因此本模块提供的是"可用的兜底"，
而不是 \crefchap{ch:transport_dd} 那套随成败增减 $\tau$ 的状态机。

\subsection{Jacobian 校验器}
\label{sec:newton-checker}

本模块另提供一个诊断钩子：把解析 Jacobian 与中心差分比较。

\begin{equation}
\mat{J}^{\text{FD}}_{:,j} \;=\;
\frac{\mat{R}(\bm{x}+h\bm{e}_j)-\mat{R}(\bm{x}-h\bm{e}_j)}{2h},
\qquad
h = 10^{-6}\max\!\bigl(1,\norm{\bm{x}}_\infty\bigr),
\label{eq:newton-fd}
\end{equation}

\begin{equation}
e_j \;=\;
\frac{\bigl\lVert \mat{J}^{\text{FD}}_{:,j} - \mat{J}_{:,j}\bigr\rVert_\infty}
{\max\bigl(\bigl\lVert \mat{J}^{\text{FD}}_{:,j}\bigr\rVert_\infty,\
\bigl\lVert \mat{J}_{:,j}\bigr\rVert_\infty,\ 10^{-300}\bigr)},
\qquad
\text{判定}\ \max_j e_j < 10^{-4}.
\label{eq:newton-checker}
\end{equation}

阈值 $10^{-4}$ 的选取有明确的理由：正确的 Jacobian 在步长适当时
通常吻合到 $10^{-6}$ 至 $10^{-10}$，因此这个阈值\emph{能抓住符号错误
与缺项，同时容忍差分公式本身的截断误差}。返回结构里带
\akey{worst_column}，使"错在哪一列"可以直接读出。

\crefeq{eq:newton-checker} 的分母取的是\textbf{逐列}的范数。
这一点在 \crefchap{ch:transport_dd} 的语境下\emph{不够} —— 那里记录过
一条教训：耦合 $(\psi,n)$ 系统里一个绝对值很小但完全错误的项，
在列归一化下看起来只有 $10^{-10}$ 的相对误差，改成逐行相对判据后
才立刻暴露。本模块的校验器对它自己的问题类（良缩放的稠密系统）
是够用的，而\crefchap{ch:transport_dd} 因此另写了自己的判据。
这条差异不是缺陷，但它是一个\textbf{必须被知道}的适用边界，
见 \S\ref{sec:newton-limits}。

\section{代价与复杂度}
\label{sec:newton-cost}
\index{1concept@概念!dense factorisation@稠密分解}

每一步的主要代价是稠密 LU 分解，为 $O(n^{3})$ 时间与 $O(n^{2})$ 内存。
此外每步至少一次残差求值（线搜索每次回溯再加一次）。

这个代价决定了本模块的\emph{适用范围上界}。\crefchap{ch:transport_dd}
的演示网格有 $2n_xn_z = 132\,922$ 个未知量 —— 稠密分解需要约
$1.4\times10^{10}$ 个浮点字的存储，完全不可行。因此
\crefchap{ch:transport_dd} \textbf{不走}这个服务，而是经 Deck 的
\akey{linear} 槽解析 \pid{linear.blocktridiag}（\crefchap{ch:blocktridiag}），
用自己的 Newton 循环加块三对角线性步。

把 \pid{nonlinear.newton} 也迁到 \akey{linear} 槽上（即让它
\emph{不}自己稠密分解，而是把线性步也交给服务）仍是欠账。
它不影响任何已实现的功能，但它意味着：换一个新的线性代数后端，
需要改两处而不是一处。

\section{验证与基准}
\label{sec:newton-verification}

验证策略是"用有闭式解或已知收敛行为的问题把算法逼到各个分支上"。
所有基准在 \fpath{tests/test_nonlinear_newton.py} 中固化。

\begin{veriftable}{Newton 迭代的收敛性与分支覆盖}{tab:newton-verify}
\vgroup{收敛性}
\vrow{线性系统}{$\le2$ 步收敛}{通过，解精确到 $10^{-10}$}
\vrow{标量根 $x^{2}-2$}{局部阶 $>1.7$}{通过（收敛到 $\sqrt2$，精确到 $10^{-12}$）}
\vrow{耦合自洽 SP 系统（\crefchap{ch:sp1d} 实测）}{二阶}{局部阶 $1.99$，
  末位残差 $1.3\times10^{-16}$}
\vgroup{线搜索}
\vrow{远起点恢复（$x^{3}-1$，自 $x=0.1$）}{收敛}{通过（精确到 $10^{-9}$，
  且阻尼历史中确实出现 $<1$ 的因子）}
\vrow{残差单调性}{逐点不增}{通过（含 $10^{-12}$ 的数值余量）}
\vrow{良好起点无需回溯}{全部阻尼因子 $=1$}{通过}
\vrow{\textbf{关闭线搜索后远起点失败}}{不收敛}{
  通过（$x=40$、上限 $4$ 步 —— 用反例证明线搜索的必要性）}
\vgroup{边界与诊断}
\vrow{迭代上限被报告}{消息含 "iteration limit"}{通过}
\vrow{奇异 Jacobian 被报告而非抛出}{不收敛 + 消息含 "linear solve failed"}{通过}
\vrow{非方阵 Jacobian 被拒绝}{抛 \texttt{ValueError}}{通过}
\vrow{PTC 位移改变步长}{步长不同}{通过}
\vrow{结果可序列化}{含四个关键字段}{通过}
\vrow{未知选项被拒绝}{抛 \texttt{TypeError}}{通过}
\vgroup{Jacobian 校验器}
\vrow{正确 Jacobian}{$<10^{-8}$}{通过}
\vrow{错一个元素（$4.0$ 应为 $3.0$）}{被检出}{
  $>10^{-2}$、\akey{ok=False}、\akey{worst_column} $=1$}
\vrow{列子集被尊重}{\akey{per_column} 长度等于选择数}{通过}
\vgroup{模块契约}
\vrow{声明 \akey{nonlinear} 能力；\akey{produces} 为空}{是}{通过}
\vrow{未知参数被拒绝}{抛 \texttt{ModuleConfigError}}{通过}
\vrow{调用方选项覆盖模块配置}{覆盖生效}{通过}
\vrow{作为流水线阶段运行被拒绝}{抛 \texttt{RuntimeError} 且含 "service"}{通过}
\end{veriftable}

\noindent
\textbf{"关闭线搜索后远起点失败"这一条是刻意的反例基准。}
它记录的不是实现的能力，而是\emph{另一个配置下的失败}，
从而证明线搜索不是装饰。这类"负向基准"在本文档中不多见，
但只要某个机制的价值只能通过对比才能说明，就应当有它。

\section{误差来源与适用边界}
\label{sec:newton-limits}
\index{1concept@概念!validity limits@适用边界}

\limitsheading

\begin{bounds}
  \item \textbf{线性步是稠密的。} $O(n^{3})$ 时间与 $O(n^{2})$ 内存
        把适用范围限制在数千个未知量以内。\crefchap{ch:transport_dd}
        的 $1.3\times10^{5}$ 个未知量它装不下，因此 DD 走
        \crefchap{ch:blocktridiag} 的线性服务而不是这里。

  \item \textbf{Jacobian 校验器用逐列归一化。} 见 \S\ref{sec:newton-checker}。
        对耦合系统（\crefchap{ch:transport_dd}）这个判据已被证明不够，
        那里另写了逐行相对判据。使用本校验器时应当知道这一区别。

  \item \textbf{不做拟 Newton 或割线法。} Jacobian 必须由调用方提供，
        且必须是精确的（或至少足够接近）。若调用方只有残差，
        本模块无能为力。

  \item \textbf{PTC 是静态的。} \akey{ptc_shift} 是常数，
        不随迭代成败自适应。真正需要延拓策略的问题应当像
        \crefchap{ch:transport_dd} 那样自建状态机。

  \item \textbf{没有信赖域。} 线搜索保证单调下降，但不保证全局收敛：
        残差曲面若有多个盆地，从坏初值出发可能收敛到另一个根，
        或在线搜索中停滞。此时唯一的补救是更好的初值。

  \item \textbf{收敛判据依赖问题自身的缩放。} 更新范数
        $\norm{\lambda\bm{\delta}}_\infty$ 的量纲是\emph{未知量}的量纲。
        因此容差必须由调用方按自己的量纲给出
        （\crefchap{ch:sp1d} 传的是以 \U{eV} 计的容差）。
        默认值 $10^{-8}$ 只是一个中性起点，不是通用推荐值。

  \item \textbf{无并行、无 GPU。} 全部为 NumPy 稠密运算，
        且每步的分解是串行的。
\end{bounds}

\section{接口契约与小结}
\label{sec:newton-interface}

\subsection{调用契约}

{\small
\begin{longtable}{@{}p{0.20\textwidth} p{0.24\textwidth} p{0.48\textwidth}@{}}
\toprule
\textbf{方向} & \textbf{名字} & \textbf{形状与语义}\\
\midrule
\endfirsthead
\toprule
\textbf{方向} & \textbf{名字} & \textbf{形状与语义}\\
\midrule
\endhead
\bottomrule
\endfoot
入 & \texttt{residual} & \texttt{x -> (R, aux)}；\texttt{R} 与 \texttt{x} 同形（一维），
  \texttt{aux} 为不透明对象\\
入 & \texttt{jacobian} & \texttt{(x, aux) -> (n, n)} 稠密数组；形状错误即抛 \texttt{ValueError}\\
入 & \texttt{x0} & 一维初值\\
入 & \texttt{options} & \texttt{NewtonOptions} 或映射；未知键即抛 \texttt{TypeError}\\
出 & \texttt{NewtonResult} & \texttt{x}, \texttt{converged}, \texttt{iterations},
  \texttt{residual\_norm}, \texttt{residual\_history}, \texttt{damping\_history},
  \texttt{linear\_solves}, \texttt{message}\\
出 & \texttt{JacobianCheck} & \texttt{max\_relative\_error}, \texttt{worst\_column},
  \texttt{per\_column}, \texttt{step}, \texttt{ok}\\
\end{longtable}
}

\noindent
\texttt{NewtonResult} 里的两个历史列表（残差与阻尼）不是调试残留，
而是\textbf{可验证性的前提}：\crefchap{ch:sp1d} 的收敛阶断言
（局部阶 $1.99$）正是从 \texttt{residual\_history} 上算出来的。
一个只返回"收敛了/没收敛"的求解器，会让调用方\emph{无法}
在测试中证明它收敛得对。

\subsection*{小结}

本章的结论可以概括为三点。

其一，\textbf{把迭代策略做成服务，换策略与换物理是同一种操作。}
一个未来的 Krylov 或信赖域求解器只要声明 \akey{nonlinear} 能力、
暴露同样的 \texttt{solve()}，就能顶上，物理模块一行不用改。

其二，\textbf{收敛判据的语义比它的数值更重要。}
用\emph{更新范数}而不是残差范数判定收敛，是对"行均衡系统的残差
存在相消地板"这一事实的正确回应；用错了不会报错，
只会让迭代在地板上空转到上限。

其三，\textbf{诊断信息是可验证性的一部分。}
残差历史、阻尼历史、\akey{worst_column} 三者合起来，
使"这个 Jacobian 对不对""这次收敛是二次的吗"成为\emph{可断言}的问题，
而不是靠输出一堆数字让人目测。

\cleardoublepage
       % 第 13 章  nonlinear.newton              v1.0.0
% module_id: linear.blocktridiag
% module_version: 0.1.0
% plugin_dir: plugins/linear_blocktridiag
% source_files: plugins/linear_blocktridiag/{__init__,solver}.py
% tests: tests/test_linear_blocktridiag.py
%
% 服务模块变体：§x.1 记 capabilities 与调用契约（无 requires/produces），
% §x.2 换为「问题类定义」。

\chapter{块三对角直接求解器}
\label{ch:blocktridiag}

\begin{quote}
\small
\noindent\textbf{模块标识}\quad\pid{linear.blocktridiag}\qquad
\textbf{版本}\quad 0.1.0\qquad
\textbf{类别}\quad solver
\end{quote}

\noindent
本章记述第二个数值服务：一个针对\emph{块三对角}结构的直接求解器。
它是 \crefchap{ch:transport_dd} 的线性步提供者，也是本项目"零第三方依赖"
这一约束下最有信息量的一个工程决策 —— 设计稿原本指定用
\texttt{scipy.sparse}，而部署环境只有 NumPy。本章说明为什么不退回稠密求解，
以及一个纯 NumPy 的块消元为什么恰好够用。

\section{模块契约}
\label{sec:block-contract}
\index{4module@模块!linear.blocktridiag}

\begin{modulecard}{linear.blocktridiag}{0.1.0}
  \cardrow{类别}{solver（数值服务，\textbf{不是}流水线阶段）}
  \cardrow{能力}{linear, linear.blocktridiag}
  \cardrow{依赖 artifact}{\textbf{无}}
  \cardrow{产出 artifact}{\textbf{无}}
  \cardrow{验证状态}{partially-verified}
  \cardrow{源码}{\fpath{plugins/linear_blocktridiag/}}
  \cardrow{测试}{\fpath{tests/test_linear_blocktridiag.py}}
  \cardrow{可标定参数}{\textbf{无}（\akey{regularization} 是数值选项）}
  \cardsep
  \cardrow{调用契约}{\texttt{solve(lower, diagonal, upper, rhs, regularization=) -> (n,m)}
    ；\texttt{dense(lower, diagonal, upper) -> (nm,nm)}（诊断用）}
  \cardrow{解析方式}{经 Deck 的 \akey{linear} 槽配置；
    也可由 \texttt{get\_by\_capability("linear")} 解析}
  \cardrow{已知消费者}{\pid{transport.dd}（\crefchap{ch:transport_dd}）}
  \cardrow{第三方依赖}{\textbf{无}（纯 NumPy；替代部署环境中缺失的 SciPy）}
\end{modulecard}

\section{问题类定义}
\label{sec:block-problem}

本模块求解的是如下形式的线性系统：

\begin{equation}
\mat{A}\,\bm{x} \;=\; \bm{b},
\qquad
\mat{A} \;=\; \operatorname{tridiag}\bigl(\mat{L}_j,\ \mat{D}_j,\ \mat{U}_j\bigr),
\label{eq:block-system}
\end{equation}

其中 $\mat{D}_j\in\mathbb{R}^{m\times m}$ 是 $n$ 个对角块，
$\mat{L}_j,\mat{U}_j\in\mathbb{R}^{m\times m}$（$j=1,\dots,n-1$）
分别是下、上次对角块；右端按块行组织为 $\bm{b}\in\mathbb{R}^{n\times m}$。

这一结构不是抽象的：\crefchap{ch:transport_dd} 的未知量按
$(\psi,n)$ 交错、横向索引排在前，竖直方向只与相邻行耦合，
于是\emph{块大小 $m=2n_x$、块行数 $n=n_z$}，Jacobian 天然是
\crefeq{eq:block-system} 的形式。演示网格下 $n=1621$、$m=82$，
总未知量 $132\,922$。

\creftab{tab:block-vs-dense} 对比两条路线。稠密分解的代价随
$(nm)^{3}$ 增长，对 $1.3\times10^{5}$ 个未知量完全不可行；
块三对角结构把它降到 $n\,m^{3}$ —— 对同一个问题约六个数量级的差距。

\begin{veriftable}{块三对角直接解与稠密解的代价对比（演示网格 $n=1621$, $m=82$）}{tab:block-vs-dense}
\vrow{分解代价}{$O\bigl((nm)^{3}\bigr)$}{$O\bigl(n\,m^{3}\bigr)$}
\vrow{按上式的量级比}{$1$}{$\approx 1/(nm)^{2}\approx 5.6\times10^{-11}$}
\vrow{存储}{$O\bigl((nm)^{2}\bigr)\approx1.2\times10^{11}$ 字}{
  $O\bigl(n\,m^{2}\bigr)\approx1.1\times10^{7}$ 字（三条带）}
\vnote{表中数字只用于说明量级，不是实测计时。实际内存还取决于是否
      保留中间消元结果；实现保留（\S\ref{sec:block-algorithm}），
      因此实际占用是三条带的常数倍。}
\end{veriftable}

\section{算法}
\label{sec:block-algorithm}
\index{3method@方法!block Thomas elimination@块 Thomas 消元}

\subsection{块消元}

算法是标量 Thomas 消元（三对角消元，经典表述见 \cite{golub2013}）
在块层面的推广。前消：

\begin{equation}
\begin{aligned}
\mat{D}'_1 &= \mat{D}_1, &\qquad \bm{b}'_1 &= \bm{b}_1,\\
\mat{D}'_j &= \mat{D}_j - \mat{L}_{j-1}\,\bigl(\mat{D}'_{j-1}\bigr)^{-1}\mat{U}_{j-1},
&\qquad
\bm{b}'_j &= \bm{b}_j - \mat{L}_{j-1}\,\bigl(\mat{D}'_{j-1}\bigr)^{-1}\bm{b}_{j-1},
\end{aligned}
\qquad j=2,\dots,n,
\label{eq:block-forward}
\end{equation}

回代：

\begin{equation}
\bm{x}_n = \bigl(\mat{D}'_n\bigr)^{-1}\bm{b}'_n,
\qquad
\bm{x}_j = \bigl(\mat{D}'_j\bigr)^{-1}\bigl(\bm{b}'_j - \mat{U}_j\,\bm{x}_{j+1}\bigr),
\qquad j=n-1,\dots,1.
\label{eq:block-backward}
\end{equation}

\paragraph{一次分解，$m+1$ 个右端。} \crefeq{eq:block-forward} 的两行共享
同一个矩阵 $\bigl(\mat{D}'_{j-1}\bigr)^{-1}$，实现上把
$\bigl[\mat{U}_{j-1}\ \big|\ \bm{b}'_{j-1}\bigr]$ 拼成一个
$m\times(m+1)$ 的右端块一次解出。这不是微优化：它把每步的
\emph{分解次数}减半，而分解正是块消元的主导代价。同一个技巧在
\crefchap{ch:sp1d} 的秩-1 亏格消元里也出现过（那里是 $N$ 列右端）。

\paragraph{块主元由 LAPACK 选取。} 每个块主元的求逆走
\texttt{numpy.linalg.solve}，即带\emph{部分主元}的 LU。
这一点值得强调：它意味着本求解器\textbf{不依赖 Jacobian 是 M-矩阵}，
也不依赖任何对角占优假设。对耦合 Newton 系统这是必要的 ——
一旦通量依赖电位，M-矩阵性质很容易丢失，而"假设它成立"的求解器
会在没有任何征兆的情况下给出错误的解。

\subsection{对角正则化}

可选的正则化把每个对角块平移：

\begin{equation}
\bigl(\mat{A} + \sigma\mat{I}\bigr)\bm{x} \;=\; \bm{b}
\qquad\Longleftrightarrow\qquad
\mat{D}_j \leftarrow \mat{D}_j + \sigma\mat{I}.
\label{eq:block-regularization}
\end{equation}

\crefeq{eq:block-regularization} 与 \crefchap{ch:transport_dd} 的伪瞬态项
$\sigma = V/\tau$ 是同一个操作。把它做成求解器的\emph{参数}而不是让
调用方重建矩阵，使调用方可以在迭代过程中调整 $\tau$ 而不必重新装配
整个 Jacobian —— \crefchap{ch:transport_dd} 的 PTC 状态机正是这样用的。

\section{代价与复杂度}
\label{sec:block-cost}
\index{1concept@概念!M-matrix}

\begin{equation}
\text{时间}\ O\bigl(n\,m^{3}\bigr),
\qquad
\text{内存}\ O\bigl(n\,m^{2}\bigr)\ \text{（三条带）}.
\label{eq:block-complexity}
\end{equation}

\subsection{两个诚实的限制}

\textbf{（i）块是稠密存储的，因此横向加密是最贵的方向。}
一条横向有 $200$ 个未知量的行，其每个带每 $1000$ 个竖直行约占
$1.6$\U{MB}（$200^2\times8$ 字节 $\times1000$）。这意味着
"加密 $n_x$"的代价是 $m^{3}$ 而"加密 $n_z$"只是线性 ——
若将来需要横向细化，带状（banded）或嵌套剖分（nested dissection）
分解是显然的下一步。

\textbf{（ii）刻意\emph{不}提供迭代回退。}
实现里没有 BiCGSTAB + ILU(0) 之类的迭代解法，理由是：
它会引入\emph{第二套收敛性故事}（迭代解的收敛与外层 Newton 的收敛
纠缠在一起），而两套收敛性同时需要调试时，定位问题会变得困难得多。
若直接分解真的成为瓶颈，那才是伸手去取迭代解法的地方，
而 Deck 的 \akey{linear} 槽正是替换件的绑定位置
（迭代解法的教材参考见 \cite{saad2003}）。

\section{验证与基准}
\label{sec:block-verification}

验证策略非常直接：\textbf{直接解可以精确对照}。每个算例都
物化出稠密矩阵、用 \texttt{numpy.linalg.solve} 求解，再逐元素比较。
构造用的是经典 M-矩阵方式（行和占优，位移 $4.0$），
因此被比较的双方都是良态的。

\begin{veriftable}{块三对角解对稠密解的逐例一致}{tab:block-verify}
\vgroup{数值正确性（判据：与稠密解之差 $<10^{-10}$）}
\vrow{$(n,m)=(2,2)$}{$<10^{-10}$}{通过}
\vrow{$(n,m)=(5,3)$}{$<10^{-10}$}{通过}
\vrow{$(n,m)=(9,4)$}{$<10^{-10}$}{通过}
\vrow{$(n,m)=(17,6)$}{$<10^{-10}$}{通过}
\vrow{单块（$n=1$）退化为稠密解}{$<10^{-12}$}{通过}
\vrow{正则化等价于平移对角块}{$<10^{-10}$}{
  通过（与显式构造的 $\mat{D}+\sigma\mat{I}$ 逐元素一致）}
\vgroup{错误处理}
\vrow{奇异块主元报出\textbf{所在行}}{异常信息含行号}{
  通过（构造 $D_2=0$ 且无耦合，信息含 "row 1"）}
\vrow{形状不合法被拒绝（\texttt{BlockTridiagonalError}）}{抛出异常}{
  通过（三条带与右端的形状组合各测一例）}
\vgroup{布局正确性}
\vrow{\texttt{to\_dense} 的块布局}{上/下/对角块落在正确位置}{
  通过（并断言远角为零）}
\vgroup{服务契约}
\vrow{声明 \akey{linear} 能力}{是}{通过（版本 $0.1.0$）}
\vrow{按能力解析得到可用求解器}{可解且 \akey{last_solve} 记录未知量数}{
  通过（$5$ 块 $\times3=15$ 个未知量）}
\vrow{作为流水线阶段运行被拒绝}{抛 \texttt{RuntimeError}}{通过}
\vrow{空注册中心报告能力缺失}{抛异常}{通过}
\vnote{"奇异块主元报出所在行"这一条的价值不在报错本身，而在\emph{行号}：
      它把"矩阵奇异"这一无从下手的诊断，变成"第 1 个块行有问题"
      这一可操作的信息。}
\end{veriftable}

\subsection{在 \texorpdfstring{\crefchap{ch:transport_dd}}{第 11 章} 中的间接证据}

本模块还有一层不在自身测试里的验证：它是 \crefchap{ch:transport_dd}
的线性步提供者，因此后者的每一项端到端基准都同时是对本模块的检验。
其中与线性代数直接相关的是耦合 Jacobian 的差分校验
（$\sim10^{-10}$，逐行相对）—— 若块消元有符号或布局错误，
Newton 根本不会收敛，更不会以二阶收敛。

\section{误差来源与适用边界}
\label{sec:block-limits}
\index{1concept@概念!validity limits@适用边界}

\limitsheading

\begin{bounds}
  \item \textbf{只处理块三对角，不处理一般稀疏。} 结构由调用方保证；
        若把非三对角的耦合塞进三条带里，会得到\emph{无提示的错误解}
        （矩阵装配正确但截断了耦合）。这是直接解固有的风险：
        结构假设不被检查。

  \item \textbf{块按稠密存储。} 横向加密的代价是 $m^{3}$，
        见 \S\ref{sec:block-cost}。$n_x$ 很大的问题需要带状或
        嵌套剖分分解。

  \item \textbf{没有迭代回退，也没有条件数估计。} 病态系统会表现为
        主元奇异（被捕获并报出行号），或者\emph{不报错但精度损失}。
        后者无法从本模块内部察觉，只能由外层 Newton 的收敛行为暴露。

  \item \textbf{正则化是标量且均匀的。} \crefeq{eq:block-regularization}
        只支持一个全局 $\sigma$；逐块或逐未知量的正则化
        （例如仅对连续性行加时间项）无法通过本接口表达。
        \crefchap{ch:transport_dd} 因此把 PTC 项加在\emph{残差}层面
        而不是线性步层面。

  \item \textbf{无部分主元的跨块选主元。} 主元只在\emph{块内}由 LAPACK
        选取；块之间不做选主元。因此它不像完全主元的高斯消元那样
        对任意非奇异矩阵稳定，而是建立在"块对角占优或结构良好"
        的常规前提上（这与标量 Thomas 消元的前提相同 \cite{golub2013}）。

  \item \textbf{单线程、无 GPU。} 前消与回代都是串行循环，
        每个块主元是一次独立的 LAPACK 调用。
\end{bounds}

\section{接口契约与小结}
\label{sec:block-interface}

{\small
\begin{longtable}{@{}p{0.18\textwidth} p{0.24\textwidth} p{0.50\textwidth}@{}}
\toprule
\textbf{方向} & \textbf{名字} & \textbf{形状与语义}\\
\midrule
\endfirsthead
\toprule
\textbf{方向} & \textbf{名字} & \textbf{形状与语义}\\
\midrule
\endhead
\bottomrule
\endfoot
入 & \texttt{lower} & $(n-1,m,m)$；\texttt{lower[k]} 把块行 $k+1$ 耦合到块行 $k$\\
入 & \texttt{diagonal} & $(n,m,m)$；$n$ 个对角块\\
入 & \texttt{upper} & $(n-1,m,m)$；\texttt{upper[k]} 把块行 $k$ 耦合到块行 $k+1$\\
入 & \texttt{rhs} & $(n,m)$；每个块行一个右端向量\\
入 & \texttt{regularization} & 标量，默认 $0$；等价于解 $(\mat{A}+\sigma\mat{I})\bm{x}=\bm{b}$\\
出 & \texttt{solution} & $(n,m)$；形状与 \texttt{rhs} 相同\\
出 & \texttt{last\_solve} & 诊断：\akey{blocks}, \akey{block_size},
  \akey{unknowns}, \akey{regularization}\\
出 & \texttt{BlockTridiagonalError} & 形状不合法或块主元奇异时抛出，含行号\\
\end{longtable}
}

\noindent
契约上有两处设计决定值得点明。\textbf{（i）三条带分开传，不传一个稀疏矩阵对象。}
这使接口不依赖任何稀疏矩阵库，调用方也不必为了用本模块而构造
SciPy 对象 —— 正是这一条让"零第三方依赖"在部署环境中成立。
\textbf{（ii）\texttt{rhs} 是矩阵而不是向量。} 它对应的是
"每个块行一个右端"的自然组织方式，同时也让
\crefeq{eq:block-forward} 的一次分解多右端技巧可以透到接口上。

\subsection*{小结}

本章的结论可以概括为三点。

其一，\textbf{结构就是性能，而结构由上游决定。}
块三对角之所以可用，是因为 \crefchap{ch:mesh} 给出了张量积网格、
\crefchap{ch:transport_dd} 按横向优先排序未知量。求解器本身没有
创造这个结构，它只是\emph{没有浪费}它。

其二，\textbf{用 LAPACK 的块主元换来一个更弱的假设。}
不假设 M-矩阵、不假设对角占优，代价是每步若干次
$m\times m$ 的带主元分解 —— 对耦合 Newton 系统，
这个交换是值得的，因为那些性质在通量依赖电位后很容易丢失。

其三，\textbf{"不做什么"也是设计的一部分。}
不加迭代回退、不加跨块选主元、不支持逐块正则化 ——
每一条都写进了 \S\ref{sec:block-limits}，连同它的后果。
一份只说能力的接口说明，会让使用者在第一次遇到
它的边界时把它当成 bug。

\cleardoublepage
    % 第 14 章  linear.blocktridiag           v0.1.0

% ---------------------------------------------------------------------
%  附录 A--D
%  \appendix 之后 \chapter 的编号变为 A/B/C/D，
%  因此 \crefapp{app:...} 印作"附录 A"。
% ---------------------------------------------------------------------
\appendix
% 附录 A —— 闭式解基准汇总

\chapter{闭式解基准汇总}
\label{app:benchmarks}

本附录把散落在各模块章的闭式解与已知解基准集中在一处，
供两类读者使用：需要快速判断"这套代码验证到什么程度"的人，
以及需要为新增模块挑选基准的人。

\section{按模块汇总}

{\small
\begin{longtable}{@{}p{0.075\textwidth} p{0.36\textwidth} p{0.17\textwidth} p{0.30\textwidth}@{}}
\toprule
\textbf{章} & \textbf{基准} & \textbf{判据} & \textbf{实测}\\
\midrule
\endfirsthead
\toprule
\textbf{章} & \textbf{基准} & \textbf{判据} & \textbf{实测}\\
\midrule
\endhead
\bottomrule
\endfoot
6 & 网格结构不变量（规模、覆盖、单调性、界面坐标、节点切片、接触次序） &
  逐条相等 & 全部通过\\
7 & D0 晶格常数对实验 & $\le1\%$ & GaN $0.345\%$/$0.328\%$，AlN $0.032\%$/$0.060\%$\\
7 & D1 弹性常数对实验 & $\le10\%$ & GaN $0.000\%$/$1.728\%$，AlN $2.778\%$/$1.877\%$\\
7 & D2 自发极化对实验 & $\le15\%$ & GaN $0.000\%$，AlN $0.000\%$\\
7 & D3 带隙（HSE 采纳值）对实验 & $\le5\%$ & GaN $0.877\%$，AlN $3.065\%$\\
7 & PBE 原始带隙不得被采纳 & 合并值 $\neq$ PBE 值 & 通过（原始值仍可查）\\
8 & Ambacher 极化面电荷（$x=0.25$） & $\in(1.10,1.30)\times10^{13}$\U{cm^{-2}} &
  $1.2020\times10^{13}$\U{cm^{-2}}\\
8 & 屏蔽严格按比例 & $10^{-6}$ 位 & 通过\\
8 & 压电极化为负；势垒侧 $p_{\text{sp}}$ 低于沟道侧 & 符号 & 通过\\
9 & 解析不动点的不变量（上界、$\Ef\ge E_0$、收敛） & 逐条 & 通过\\
9 & 面密度落在物理区间 & $4.0$--$12.0\times10^{12}$\U{cm^{-2}} & 通过\\
10 & 双介质平行板（无面电荷） & $<10^{-10}$\U{V} & $<10^{-10}$\U{V}（精确）\\
10 & 双介质平行板 $+$ 界面面电荷 & $<10^{-10}$\U{V} & $<10^{-10}$\U{V}（精确）\\
10 & $\Ec$ 算子 vs.\ $\varphi$ 形式独立离散化 & $<10^{-10}$\U{eV} &
  $2.9\times10^{-13}$\U{eV}\\
10 & 无限深方阱 $E_n=n^2\pi^2\hb^2/(2\meff L^2)$ & 二阶收敛 & 二阶收敛\\
10 & 谐振子 $E_n=\hb\omega(n+\tfrac12)$ & $<10^{-3}$\U{eV} & 通过\\
10 & 三角阱 Airy 零点 $E_n=\abs{\zeta_n}(\hb^2F^2/2\meff)^{1/3}$ & 二阶收敛 &
  $5.5\times10^{-5}\to3.4\times10^{-6}$\\
11 & 均匀电阻 $J=qN_D\mumu V/L$ & $<0.1\%$ & $\mathbf{<10^{-6}}$\\
11 & 均匀电阻两端电流一致 & 一致 & $<10^{-9}$\\
11 & 电流与偏压严格正比 & 比值 $=2$ & $2.0\pm10^{-6}$\\
11 & Bernoulli 函数 $B(x)=x/(\ee^x-1)$ & $10^{-12}$ 相对 & 通过（宽区间）\\
11 & 反射恒等式 $B(-x)=B(x)+x$ & 恒等 & 通过\\
11 & 通量四个偏导对中心差分 & 相对误差小 & 通过（含场依赖迁移率）\\
11 & 异质结台阶上的零通量 & $<10^{-12}$ & 通过\\
11 & 平衡态端点电流 & $<10^{-12}$\U{A/m} & $2.7\times10^{-15}$\U{A/m}\\
13 & 线性系统 & $\le2$ 步 & 通过（精确到 $10^{-10}$）\\
13 & 标量根 $x^2-2$ & 局部阶 $>1.7$ & 通过（收敛到 $\sqrt2$）\\
13 & 正确 Jacobian 被校验器确认 & $<10^{-8}$ & 通过\\
13 & 错一个元素被检出 & $>10^{-2}$ 且 \akey{ok=False} & 通过（\akey{worst_column} 正确）\\
14 & 块三对角解对稠密解（四组尺寸） & $<10^{-10}$ & 通过\\
14 & 单块退化为稠密解 & $<10^{-12}$ & 通过\\
14 & 正则化等价于平移对角块 & $<10^{-10}$ & 通过\\
\end{longtable}
}

\section{按"能证明什么"重排}

同一批基准若按 §\ref{sec:verify-criteria} 的五类判据重排，
可以看出验证覆盖的分布：

\begin{veriftable}{基准的类型分布}{tab:appA-kinds}
\vrow{闭式解（格式精度）}{检验离散格式}{
  第 10 章（4 项）、第 11 章（4 项）、第 14 章（3 项）}
\vrow{收敛阶（阶数正确）}{检验加密行为}{
  第 10 章（3 项）、第 11 章（1 项）}
\vrow{制造解（任意解上的精度）}{检验算子}{
  第 11 章（1 项，观测阶 $2.008/2.002/2.001$）}
\vrow{不变量与趋势}{检验必须成立的关系}{
  第 6、8、9、10、11 章}
\vrow{对实验（材料层）}{检验常数而非格式}{
  第 7 章（D0--D3 四级）}
\vnote{两类判据明显地不均衡：\textbf{格式侧的证据远多于器件侧}。
      这正是 \crefchap{ch:verification-methodology} 中"基准阶梯只走到 L1"
      的具体表现 —— 本表下半部分（对实测）几乎是空的。}
\end{veriftable}

\section{尚未使用的基准类型}

以下基准类型在方案中列出但在本文档当前版本中尚无用例，
记录在此以说明验证体系的\emph{空白}而非仅仅它的成果：

\begin{itemize}[leftmargin=1.7em,itemsep=0.28em,topsep=0.35em]
  \item \textbf{网格收敛指数（GCI）}：本项目的收敛阶断言用的是
        Richardson 三点公式，尚未按 GCI 规程给出离散误差的置信区间；
  \item \textbf{与独立实现的交叉对比}：除 \crefchap{ch:transport_dd}
        的 MMS 与 \crefchap{ch:sp1d} 的双路径 Jacobian 之外，
        尚无与第三方 TCAD 的对比（L3）；
  \item \textbf{不确定度量化}：无敏感性分析、无蒙特卡洛或多项式混沌。
\end{itemize}

\cleardoublepage
            % 附录 A  闭式解基准汇总
% 附录 B —— 可标定参数总表

\chapter{可标定参数总表}
\label{app:parameters}

本附录列出全部标记为\emph{可标定}的参数。它们的共同性质是：
默认值未经 L2 基准（霍尔 / XRD / C--V 实测）验证，使用前必须自行标定。
凡未列入本表的参数都是\emph{数值或几何选择}，不应被自动拟合循环挪动。

\section{总表}

{\small
\setlength{\tabcolsep}{4pt}
\begin{longtable}{@{}p{0.045\textwidth} p{0.175\textwidth} p{0.28\textwidth}
                    p{0.075\textwidth} p{0.135\textwidth} p{0.18\textwidth}@{}}
\toprule
\textbf{章} & \textbf{模块} & \textbf{参数} & \textbf{默认} & \textbf{单位} & \textbf{取值依据（录自 MANIFEST）}\\
\midrule
\endfirsthead
\toprule
\textbf{章} & \textbf{模块} & \textbf{参数} & \textbf{默认} & \textbf{单位} & \textbf{取值依据（录自 MANIFEST）}\\
\midrule
\endhead
\bottomrule
\endfoot
8 & \pid{physics.polarization.ambacher} & \akey{screening} & $0.8$ & --- &
  \textbf{无公认取值}；未被屏蔽的理论值在实践中高估面密度，须用霍尔数据标定\\
8 & 同上 & \akey{relaxed_fraction} & $1.0$ & --- &
  默认假定赝晶生长；除非 TEM/XRD 另有结论\\
\midrule
9 & \pid{physics.quantum.analytic_tw} & \akey{phi_b_0_eV} & $0.84$ & \U{eV} &
  Ni/AlGaN 势垒拟合 $\phib\approx0.84+0.39x$\\
9 & 同上 & \akey{phi_b_x_eV} & $0.39$ & \U{eV} & 同上\\
9 & 同上 & \akey{delta_ec_ratio} & $0.7$ & --- &
  常用值 $0.7$；实验报道约跨 $0.7$--$1.0$\\
9 & 同上 & \akey{gamma0_eV_m43} & $2.123\times10^{-12}$ & \U{eV\cdot m^{4/3}} &
  GaN 的三角阱预因子，取自文献 \cite{ambacher2000}\\
\midrule
10 & \pid{physics.quantum.sp1d} & \akey{phi_b_0_eV} & $0.84$ & \U{eV} & 同第 9 章\\
10 & 同上 & \akey{phi_b_x_eV} & $0.39$ & \U{eV} & 同第 9 章\\
10 & 同上 & \akey{delta_ec_ratio} & $0.7$ & --- & 同第 9 章\\
10 & 同上 & \akey{gate_bias_V} & $0.0$ & \U{V} & 平衡态求解所处的栅偏压\\
\midrule
11 & \pid{transport.dd} & \akey{temperature_K} & $300.0$ & \U{K} &
  等温晶格温度；热耦合属未来工作\\
11 & 同上 & \akey{mu0_cm2_Vs} & $1500.0$ & \U{cm^2/(V\cdot s)} &
  \textbf{占位值}，待迁移率插件与霍尔标定\\
11 & 同上 & \akey{vsat_cm_s} & $2.0\times10^{7}$ & \U{cm/s} &
  \textbf{占位值}，同上；非正值退回常数迁移率\\
11 & 同上 & \akey{delta_ec_ratio} & $0.7$ & --- &
  仅在无量子模块提供 $\dEc$ 时使用\\
\midrule
12 & \pid{extract.fom} & \akey{vth_correction_V} & $0.0$ & \U{V} &
  吸收电荷控制估计与实测阈值之差的经验偏移\\
12 & 同上 & \akey{interface_charge_cm2} & $0.0$ & \U{cm^{-2}} &
  从二维电子气面密度中扣除的净固定界面电荷\\
\end{longtable}
}

\section{两个零与一个"无"}

表中有三处值得强调的取值，它们都是方法论上的选择而非疏漏。

\subsection*{\pid{materials.dft}：一个可标定参数都没有}

第 7 章的模块\textbf{没有任何可标定参数}，且这一条被测试固化为断言
（\akey{calibratable_params()} 返回空元组）。理由是：
\akey{ecut_wfc_ry}、\akey{kspacing_1_A}、\akey{pseudopotential_family}
这些量是\emph{数值方法的选择}，不是器件的物理参数。若允许自动拟合循环
去调整它们，就等于\emph{为了对上数据而调网格} —— 那会把标定与
收敛性混为一谈，并且系统性地破坏两者的可信度。

\subsection*{\pid{extract.fom}：两个零}

\akey{vth_correction_V} 与 \akey{interface_charge_cm2} 默认为 $0$，
含义是"暂时不修正"。这两个参数的存在本身就是一种诚实：
它们承认电荷控制估计与实测阈值之间\emph{必然}存在差异，
并把差异\emph{显式地}放在一个可标定位置上，
而不是把它藏进某个物理常数里。

\subsection*{$\akey{vth}$ 与 $\ns$ 的可标定性不对称}

第 12 章记录了一个容易被误解的事实：本模块报告的阈值电压
\emph{只依赖势垒静电学}，与量子模型无关；只有面密度来自量子求解。
因此若用霍尔数据标定 $\ns$，受影响的参数在
第 8 章（\akey{screening}）；而标定 $V_{\text{th}}$ 影响的是
第 12 章的两个参数。\textbf{两类标定针对不同的参数，不应混做一次。}

\section{标定的建议顺序}

基于依赖关系与不确定度大小，建议的标定顺序是：

\begin{enumerate}[leftmargin=1.9em,itemsep=0.4em,topsep=0.4em]
  \item \textbf{\akey{screening}（第 8 章）} —— 它是极化链上
        最大的不确定度来源，且直接影响 $\ns$ 与 $V_{\text{th}}$ 两侧；
  \item \textbf{\akey{mu0_cm2_Vs} 与 \akey{vsat_cm_s}（第 11 章）} ——
        它们不影响平衡态电荷，只影响电流量级与饱和行为，
        因此可以在电荷标定之后独立进行；
  \item \textbf{\akey{vth_correction_V} 与 \akey{interface_charge_cm2}
        （第 12 章）} —— 在前两步之后，用 Id--Vg 实测收尾；
  \item \textbf{\akey{relaxed_fraction}、\akey{delta_ec_ratio}、
        \akey{gamma0_eV_m43}} —— 这三者应当优先用\emph{结构}表征
        （XRD、TEM）与文献共识确定，而不是用器件电学数据反推。
\end{enumerate}

最后一条尤其重要：$f_{\text{relax}}$ 与 $r$ 都是可以用非电学手段
独立确定的量。用电学数据去反推它们，会让多个参数吸收同一个误差，
从而失去可解释性。

\section{未列入本表的参数}

以下参数\emph{不}应被标定，各章已给出理由：

\begin{itemize}[leftmargin=1.7em,itemsep=0.28em,topsep=0.35em]
  \item \textbf{网格参数}（\akey{nx}、\akey{target_spacing_nm} 等）：
        属于离散分辨率，应通过收敛性研究确定（\crefchap{ch:mesh}）；
  \item \textbf{迭代控制参数}（\akey{max_iterations}、\akey{tolerance}、
        \akey{max_step_V}、\akey{ptc_tau0_s} 等）：
        属于数值选择，只影响收敛路径与代价（\crefchap{ch:newton}、
        \crefchap{ch:transport_dd}）；
  \item \textbf{DFT 数值设置}（\akey{ecut_wfc_ry}、\akey{kspacing_1_A}、
        \akey{pseudopotential_family}）：属于方法选择，
        且被陈旧判据保护（\crefchap{ch:dft}）；
  \item \textbf{\akey{density_clamp}}：默认关闭的数值守卫，
        其理由是数值的而非物理的（\crefchap{ch:sp1d}）。
\end{itemize}

\cleardoublepage
            % 附录 B  可标定参数总表
% 附录 C —— 版本、指纹与 artifact 键矩阵

\chapter{版本、指纹与 artifact 键矩阵}
\label{app:contracts}

\section{模块版本记录}

{\small
\begin{longtable}{@{}p{0.055\textwidth} p{0.26\textwidth} p{0.085\textwidth} p{0.50\textwidth}@{}}
\toprule
\textbf{章} & \textbf{模块标识} & \textbf{版本} & \textbf{版本语义说明}\\
\midrule
\endfirsthead
\toprule
\textbf{章} & \textbf{模块标识} & \textbf{版本} & \textbf{版本语义说明}\\
\midrule
\endhead
\bottomrule
\endfoot
6 & \pid{mesh.structured2d} & 0.1.0 & 首个版本；横向栅边缘加密列为后续计划\\
7 & \pid{materials.dft} & 0.1.0 & 首个版本；\akey{table} 后端已落地，\akey{qe_*} 未实现\\
8 & \pid{physics.polarization.ambacher} & 0.1.0 & 首个版本；仅 Ga 面极性\\
9 & \pid{physics.quantum.analytic_tw} & 0.1.0 & 首个版本；作为基线保留\\
10 & \pid{physics.quantum.sp1d} & \textbf{1.1.1} &
  1.1.0 $\to$ 1.1.1：Jacobian 做到严格精确（全谱波函数响应 + 质量响应），
  数值行为改变，\emph{因此必须升版本}\\
11 & \pid{transport.dd} & 0.1.0 & 首个版本；收尾时补齐场依赖迁移率\\
12 & \pid{extract.fom} & 0.1.0 & 首个版本；占位实现（\texttt{unverified}）\\
13 & \pid{nonlinear.newton} & 1.0.0 & 首个稳定版本\\
14 & \pid{linear.blocktridiag} & 0.1.0 & 首个版本；替代设计稿中的 SciPy 稀疏步\\
\end{longtable}
}

\noindent
表中第 10 章一行是版本语义的一个实例：\textbf{数值行为改变必须升版本}，
因为指纹不捕获数值（见 \crefchap{ch:architecture}）。

\section{指纹记录}

\begin{veriftable}{已知的指纹值与语义}{tab:appC-fingerprints}
\vrow{解析量子 Deck（\fpath{demo_hemt_sp1d.json} 之前的身形）}{
  两次运行一致}{\texttt{bcfbb07632737b44}}
\vrow{自洽量子 Deck（\fpath{demo_hemt_sp1d.json}）}{
  两次运行一致}{\texttt{9d716641ad967853}}
\vrow{同一 Deck 两次运行}{相等}{通过}
\vrow{两个不同实现的 Deck}{不等}{通过}
\vrow{材料槽替换后的 Deck}{不等}{通过（\crefchap{ch:dft}）}
\vnote{指纹的组成是"参数 + 模块版本 + 材料库指纹 + 标记"。
      \textbf{artifact 数值不进指纹}，因此上表第二行与第一行的差异
      并不是"数值变了"这一事实本身，而是模块版本从 1.1.0 升到 1.1.1
      的结果。这条语义在第 \ref{ch:architecture} 章有完整说明。}
\end{veriftable}

\section{artifact 键与单位矩阵}

本附录与 \creftab{tab:data-matrix}（契约矩阵：谁产谁用）互补 ——
这里给出\emph{每个键的单位}，因为跨模块最常见的误解就是单位。

{\small
\setlength{\tabcolsep}{4pt}
\begin{longtable}{@{}p{0.23\textwidth} p{0.275\textwidth} p{0.12\textwidth} p{0.275\textwidth}@{}}
\toprule
\textbf{artifact} & \textbf{键} & \textbf{单位} & \textbf{说明}\\
\midrule
\endfirsthead
\toprule
\textbf{artifact} & \textbf{键} & \textbf{单位} & \textbf{说明}\\
\midrule
\endhead
\bottomrule
\endfoot
\akey{structure} & \akey{layers[].thickness_nm} & \U{nm} & 层厚\\
 & \akey{layers[].x_al} & --- & AlN 摩尔分数，$\in[0,1]$\\
 & \akey{layers[].doping_cm3} & \U{cm^{-3}} & 净掺杂浓度\\
 & \akey{gate.length_um} & \U{\mu m} & 栅长\\
\midrule
\akey{mesh} & \akey{x_nm[]}, \akey{z_nm[]} & \U{nm} & 两个一维轴\\
 & \akey{interfaces_nm[]} & \U{nm} & 层界坐标（落在节点上）\\
 & \akey{vertical_spacing_nm} & \U{nm} & 竖直间距极值（\akey{{min,max}} 两个子键）\\
 & \akey{contacts_nm.*} & \U{nm} & 源/栅/漏接触位置\\
 & \akey{layers[].node_slice} & --- & 层 $\to$ 节点区间（可反查）\\
\midrule
\akey{materials} & \akey{a_lattice_A}, \akey{c_lattice_A} & \U{\AA} & 晶格常数\\
 & \akey{p_sp_C_m2} & \U{C/m^2} & 自发极化\\
 & \akey{e31_C_m2}, \akey{e33_C_m2} & \U{C/m^2} & 压电常数\\
 & \akey{C13_GPa}, \akey{C33_GPa} & \U{GPa} & 弹性常数\\
 & \akey{eps_r}, \akey{m_eff} & --- & 相对介电常数、有效质量比\\
 & \akey{Eg_eV} & \U{eV} & 带隙\\
\midrule
\akey{polarization} & \akey{barrier_thickness_nm} & \U{nm} & 势垒厚度\\
 & \akey{sigma_pol_C_m2}, \akey{sigma_eff_C_m2} & \U{C/m^2} & 极化 / 有效面电荷\\
 & \akey{sigma_pol_cm2}, \akey{sigma_eff_cm2} & \U{cm^{-2}} & 同一量的电荷密度表示（除以 $q$）\\
 & \akey{strain}, \akey{relaxed_fraction}, \akey{screening} & --- & 无量纲\\
 & \akey{band_offset_eV} & \U{eV} & 带隙差 $\Delta E_g$（不是 $\dEc$）\\
\midrule
\akey{quantum} & \akey{ns_m2} & \U{m^{-2}} & 面密度\\
 & \akey{ns_cm2} & \U{cm^{-2}} & 同一量的工程单位\\
 & \akey{ns_upper_bound_cm2} & \U{cm^{-2}} & $\sigeff/q$\\
 & \akey{ef_eV}, \akey{e0_eV} & \U{eV} & 费米能级、基态\\
 & \akey{delta_ec_eV} & \U{eV} & $\dEc=r\,\Delta E_g$\\
\akey{ns} & （标量） & \U{cm^{-2}} & 面密度的便捷通道\\
\midrule
\akey{electrostatics} & \akey{psi_V} & \U{V} & 二维电势场（最后偏压点）\\
 & \akey{gate_reference_V} & \U{V} & 栅参考电位\\
\akey{carriers} & \akey{density_m3} & \U{m^{-3}} & 二维电子密度场\\
 & \akey{ns_cm2} & \U{cm^{-2}} & 横向平均后的面密度\\
 & \akey{mobility_m2_Vs} & \U{m^2/(V\cdot s)} & 低场迁移率\\
\akey{current} & \akey{id_A_per_m}, \akey{is_A_per_m} & \U{A/m} &
  \textbf{每单位器件宽度}\\
 & \akey{continuity_error} & --- & 相对电流不平衡\\
\akey{iv} & \akey{vds_V[]}, \akey{vgs_V[]} & \U{V} & 偏压轴\\
 & \akey{id_A_per_m[]} & \U{A/m} & 漏电流（每单位宽度）\\
\midrule
\akey{results.fom} & \akey{vth_est_V} & \U{V} & 电荷控制阈值估计\\
 & \akey{barrier_drop_V} & \U{V} & $\sigeff d/(\epsr\epsO)$\\
 & \akey{cox_F_m2} & \U{F/m^2} & 势垒层单位面积电容\\
 & \akey{cox_uF_cm2} & \U{\mu F/cm^2} & 同一量的工程单位\\
 & \akey{thermal_voltage_V} & \U{V} & $kT/q$\\
\midrule
\akey{dft_provenance} & \akey{numerics.*} & 混合 & 截断能（\U{Ry}）、$k$ 间距（\U{1/\AA}）等\\
 & \akey{entries.*.uncertainty} & 混合 & 逐值区间 $[v_{\text{lo}},v_{\text{hi}}]$\\
\end{longtable}
}

\noindent
表中三处\textbf{同一量的双单位}（\akey{sigma_*_C_m2} 与
\akey{sigma_*_cm2}、\akey{ns_m2} 与 \akey{ns_cm2}、
\akey{cox_F_m2} 与 \akey{cox_uF_cm2}）是有意为之：
国际单位制用于计算，工程单位用于报告，两者同时提供以避免
下游各自换算。这也意味着\textbf{改动其中之一必须同时改另一个}，
一致性由各章的验证表锁定。

\section{契约的机械可校验性}

本文档的契约矩阵与单位的最终真相在代码里（\texttt{MANIFEST} 与各模块的
\texttt{ctx.put} 调用）。当前两者的一致性靠人工维护；
把"文档里的键" 与"代码里的键"做机械比对是一件
\emph{明确可行但尚未实现}的工程化工作，
它与 \crefchap{ch:intro} 提到的"章首模块标识与 MANIFEST 逐字段比对"
是同一件事。

\cleardoublepage
             % 附录 C  版本、指纹与 artifact 键矩阵
% 附录 D —— 已知边界与欠账清单

\chapter{已知边界与欠账清单}
\label{app:debts}

本附录汇总全文的"不知道"与"还没做"。它与附录 \ref{app:benchmarks}
（已证明的）成对：一份只说成果的文档，会让使用者在第一次撞到边界时
把它当成缺陷。

\section{原理上不可达的判据}

有一类"欠账"不是待办事项，而是\emph{在本离散下不可能达成}的判据。
把它们与普通欠账分开，是因为继续投入不会使它们消失。

\begin{veriftable}{原理上不可达的判据}{tab:appD-unreachable}
\vrow{与一维自洽 SP 的带边差 $<1$\U{meV}}{设计稿 T1 的一半}{
  \textbf{不可达}：网格把层界放在节点上，界面处被一个单元宽度抹平
  （\crefchap{ch:mesh} 与 \crefchap{ch:transport_dd}）}
\vrow{单个波函数是物理的束缚态}{物理性}{
  \textbf{不成立}：势垒以上的态是有限域盒本征态。它们是三维连续态的
  合法求积（域深度依赖 $<10^{-5}$），但单个波函数不是束缚态
  （\crefchap{ch:sp1d}）}
\vnote{这两条都写在对应章的 §x.6 里，此处汇总以便一次性看到全貌。
      第一条若要达成，需要让输运网格也采用面基界面 ——
      那会改变第 6 章的输出契约，属于架构级改动，不在当前路线图中。}
\end{veriftable}

\section{全局欠账（按优先级）}

\subsection*{第一优先：使绝对数值可信}

\begin{enumerate}[leftmargin=1.9em,itemsep=0.4em,topsep=0.4em]
  \item \textbf{L2 基准 + 回归门禁。} 用霍尔测量对标二维电子气面密度，
        把测试套件接入 CI 门禁。这是\emph{一切绝对值}的前提：
        在此之前，本文档中所有"实测"列的数字都只能作为
        \emph{格式与守恒性}的证据，不能作为器件数值的证据。
  \item \textbf{迁移率独立插件。} Deck 的 \akey{mobility} 槽当前为空，
        \crefchap{ch:transport_dd} 把 $\mumu_0$ 与 $v_{\text{sat}}$
        作为自身参数（占位值）。一个独立的迁移率插件（如场依赖模型、
        合金散射、界面粗糙度散射）是让电流数值可信的必要条件。
  \item \textbf{材料常数的独立验证。} \crefchap{ch:dft} 的
        \akey{table} 后端是\emph{文献转录}，其查找表的
        \texttt{honesty\_note} 明确声明了这一点，且全部条目
        \akey{convergence: null}。真正的硬度提升要等
        \akey{qe_local} 重算。
\end{enumerate}

\subsection*{第二优先：扩展物理范围}

\begin{enumerate}[leftmargin=1.9em,itemsep=0.4em,topsep=0.4em,start=4]
  \item \textbf{陷阱与自热。} 电流坍塌、gate-lag、drain-lag、
        晶格温升一概没有；因此射频大电流点的预测偏乐观。
        对应 \akey{traps} 与 \akey{thermal} 两个空槽。
  \item \textbf{短栅的流体动力学修正。} $L_g=150$\U{nm} 已在漂移扩散
        假设开始失效的尺度上；速度过冲与弹道成分缺席。
        方案给出的对策是 \pid{transport.hydro} 插件加经验修正，
        并\textbf{在报告中标注}。
  \item \textbf{高偏压收敛。} \crefchap{ch:transport_dd} 只覆盖到约
        $1$\U{V} 的漏压；$V_{ds}\ge2$\U{V} 时 Newton 仍不收敛。
        这是模型边界而非调参问题。
  \item \textbf{\pid{materials.dft} 的相 2 与相 3。}
        \akey{qe_local} 驱动（结构弛豫 → SCF → 应力/Born → 后处理）
        与 \akey{qe_slurm} + SQS 合金扫描均未实现；
        当前只完成可用性探测与清晰的失败信息。
\end{enumerate}

\subsection*{第三优先：提取链与性能}

\begin{enumerate}[leftmargin=1.9em,itemsep=0.4em,topsep=0.4em,start=8]
  \item \textbf{提取链。} 完整的 IV / C--V / S 参数提取
        （$g_m$、$f_T$、$f_{\max}$）尚未实现；
        \crefchap{ch:fom} 只固定了 \akey{results.fom} 这个契约位置，
        并明确标记为 \texttt{unverified}。
  \item \textbf{精确 Jacobian 的装配代价。} \crefchap{ch:sp1d} 的
        精确 Jacobian 使装配贵约一个数量级、整解贵 $1.5$--$1.7$ 倍。
        两条提速路线已明确：(a) 三对角逆的\emph{半可分}（semiseparable）
        Green 函数两步扫描降到 $O(N)$ 次向量化运算；
        (b) 线性步改为\emph{矩阵无关}的 Krylov 方法，
        响应只需对向量作用。
  \item \textbf{把 \pid{nonlinear.newton} 迁到 \akey{linear} 槽。}
        它当前的线性步是自己做的稠密 LU，因此它装不下
        \crefchap{ch:transport_dd} 的 $1.3\times10^{5}$ 个未知量。
        迁移之后，"换线性代数后端"才只需改一处
        （\crefchap{ch:newton}）。
  \item \textbf{横向细化能力。} \crefchap{ch:blocktridiag} 的块是稠密
        存储的，横向加密代价为 $O(m^3)$；带状或嵌套剖分分解是
        显然的下一步。
\end{enumerate}

\section{逐模块边界索引}

每章的 §x.6 列出了该模块自己的适用边界。此处给出索引，
使读者可以按模块定位：

{\small
\begin{longtable}{@{}p{0.075\textwidth} p{0.20\textwidth} p{0.60\textwidth}@{}}
\toprule
\textbf{章} & \textbf{模块} & \textbf{边界的要点（详见该章 §x.6）}\\
\midrule
\endfirsthead
\toprule
\textbf{章} & \textbf{模块} & \textbf{边界的要点（详见该章 §x.6）}\\
\midrule
\endhead
\bottomrule
\endfoot
6 & \pid{mesh.structured2d} &
  界面落节点（$1.7$ 倍面密度偏差）；横向均匀无栅边缘加密；
  张量积几何表达力有限；$N_{\min}$ 可能造成局部过密；只覆盖外延栈\\
7 & \pid{materials.dft} &
  文献转录非自算；全部条目缺两点收敛记录；0\,K 理想晶体；
  无温度修正；合金为端点插值；\akey{qe_*} 未实现；
  不确定度区间不是置信区间\\
8 & \pid{physics.polarization.ambacher} &
  屏蔽因子无公认取值；材料常数未独立验证；只支持 Ga 面；
  无表面态模型；无应变弛豫动力学；合金为端点插值\\
9 & \pid{physics.quantum.analytic_tw} &
  不含任何微分方程；忽略波函数穿透与多亚带；
  态密度质量取势垒侧参数（已知不一致）；$\gamma_0$ 是拟合常数；
  不产生势剖面与波函数\\
10 & \pid{physics.quantum.sp1d} &
  材料常数未验证；屏蔽因子无公认取值；有效质量近似忽略非抛物性；
  只有 $\Gamma$ 谷；盒本征态；一维近似；Anderson 回退较脆弱；
  装配较贵；无陷阱/自热\\
11 & \pid{transport.dd} &
  迁移率是占位参数；只覆盖到约 $1$\U{V}；短栅速度过冲缺席；
  界面抹平；只实现 Boltzmann 统计；缓冲层是抽象；
  无栅电流；无陷阱/自热/瞬态；表面钉扎是工程近似；
  PTC 参数是工程值\\
12 & \pid{extract.fom} &
  占位实现（\texttt{unverified}）；全耗尽假定；
  $V_{\text{th}}$ 与量子模型无关；无串联电阻与短沟效应；
  无栅电流；没有真正的优值链\\
13 & \pid{nonlinear.newton} &
  稠密线性步限制规模；校验器用逐列归一化（对耦合系统不够）；
  不做拟 Newton；PTC 是静态的；无信赖域；
  收敛判据依赖问题缩放；无并行\\
14 & \pid{linear.blocktridiag} &
  只处理块三对角（结构假设不被检查）；块稠密存储；
  无迭代回退与条件数估计；正则化是标量且均匀；
  不做跨块选主元；单线程\\
\end{longtable}
}

\section{工程化欠账（文档侧）}

本文档自身也有尚未兑现的工程化工作，如实列出：

\begin{enumerate}[leftmargin=1.9em,itemsep=0.4em,topsep=0.4em]
  \item \textbf{章首模块标识与 \texttt{MANIFEST} 的机械比对。}
        本文档在每个模块章开头以注释给出 \akey{module_id}、
        \akey{module_version}、\akey{plugin_dir} 等锚点，
        但比对尚未自动进行。
  \item \textbf{契约键的机械比对。} 附录 \ref{app:contracts} 的单位矩阵
        与代码中的 \texttt{ctx.put} 调用之间靠人工维护。
  \item \textbf{数字的溯源表。} 正文中每个"实测"值目前靠
        "可回查到测试文件"这一人工约定；一个把数字与测试断言
        直接关联的溯源表尚未建立。
  \item \textbf{验证表的列宽等版式约束}已记录在
        \fpath{docs/tex/README.md}，属工程细节而非方法学问题。
\end{enumerate}

\section{一句话总结}

\textbf{骨架可信，量子求解器与 Newton 求解器都验证到闭式解，
网格收敛与 Newton 收敛阶都修到了二阶（Jacobian 严格精确）；
漂移扩散已能跑出守恒到 $10^{-13}$ 的 IV 曲线，通过 T0 闭式解
（$<10^{-6}$）、T2 二阶收敛与电流饱和判据；
但绝对数值未标定，$V_{ds}\ge2$\U{V} 仍不收敛，
基准阶梯只走到 L1。}

\cleardoublepage
                 % 附录 D  已知边界与欠账清单

% =====================================================================
\backmatter
% =====================================================================
%  只印被正文引用过的文献。refs.bib 中为后续章节预置的条目
%  在对应章节成文之前不会出现在此列表中 —— 参考文献列表反映的是
%  "本文档引用了什么"，不是"库房里有什么"。
\printbibliography[heading=bibintoc,title={参考文献}]
\printindex

\end{document}
